\documentclass[11pt,a4paper]{article}

% ============================================
% PACKAGES
% ============================================
\usepackage[utf8]{inputenc}
\usepackage[T1]{fontenc}
\usepackage{amsmath,amsthm,amssymb}
\usepackage{mathtools}
\usepackage{array}
\usepackage{enumitem}
\usepackage{booktabs}
\usepackage{geometry}
\usepackage{hyperref}
\usepackage{cleveref}
\usepackage{float}
\usepackage{microtype}
\usepackage{xcolor}
\usepackage{stmaryrd}

% MiKTeX + stmaryrd can trigger non-scalable font expansion failures.
% Keep microtype enabled but disable expansion globally for portability.
\microtypesetup{expansion=false}

\geometry{margin=1in}

\hypersetup{
    colorlinks=true,
    linkcolor=blue!70!black,
    citecolor=green!50!black,
    urlcolor=blue!70!black,
    hypertexnames=false
}
\emergencystretch=2em
\pdfstringdefDisableCommands{%
  \def\leq{<=}%
  \def\geq{>=}%
  \def\to{->}%
  \def\otimes{otimes}%
  \def\nabla{nabla}%
  \def\wedge{wedge}%
  \def\flat{flat}%
  \def\sharp{sharp}%
  \def\bigcirc{next}%
  \def\Diamond{eventually}%
}

% ============================================
% THEOREM ENVIRONMENTS
% ============================================
\theoremstyle{plain}
\newtheorem{theorem}{Theorem}[section]
\newtheorem{lemma}[theorem]{Lemma}
\newtheorem{proposition}[theorem]{Proposition}
\newtheorem{corollary}[theorem]{Corollary}

\theoremstyle{definition}
\newtheorem{definition}[theorem]{Definition}
\newtheorem{example}[theorem]{Example}
\newtheorem{axiom}[theorem]{Axiom}

\theoremstyle{remark}
\newtheorem{remark}[theorem]{Remark}
\newtheorem{observation}[theorem]{Empirical Observation}

% ============================================
% CUSTOM COMMANDS
% ============================================
\newcommand{\N}{\mathbb{N}}
\newcommand{\Z}{\mathbb{Z}}
\newcommand{\R}{\mathbb{R}}
\newcommand{\D}{\mathbb{D}}
\newcommand{\U}{\mathcal{U}}
\newcommand{\Disc}{\mathrm{Disc}}

% ============================================
% TITLE
% ============================================
\title{\textbf{Geometric Type-Theoretic Library Growth\\
via the Principle of Efficient Novelty}}

\author{Halvor Lande\\
\texttt{hsl@awc.no}}

\date{March 2026}

% ============================================
% DOCUMENT
% ============================================
\begin{document}

\maketitle

% ============================================
% ABSTRACT
% ============================================
\begin{abstract}
Automated discovery in formal mathematics traditionally focuses on searching for proofs to human-specified conjectures. This paper introduces an alternative paradigm: treating the foundational growth of the mathematical library itself as the optimization target. We present the Principle of Efficient Novelty (PEN), an executable metatheoretic model that iteratively extends an initially empty accepted library over a fixed intensional type-theoretic signature. At each step, PEN selects the admissible candidate that maximizes structural efficiency ($\rho = \nu/\kappa$, the ratio of marginal derivational capacity to irreducible specification clauses) while clearing a history-dependent integration bar.

The manuscript separates three claim layers. At the \emph{PEN-theorem layer}, we prove that in fully coupled intensional Homotopy Type Theory (specifically CCHM/HoTT), maintaining global canonicity imposes a Coherence Window of depth $d=2$. This two-step memory forces the interface debt of new structures to scale according to the Fibonacci sequence, raising the survival threshold at the rate of the Golden Ratio. Consequently, purely discrete structures (e.g., natural numbers, classical logic) yield only linear novelty and are permanently rejected, while geometric and topological structures naturally dominate because their higher path constructors generate the superlinear combinatorial capacity required to outpace the exponentially rising bar.

At the \emph{executable layer}, the current Rust \texttt{pen-atomic} workspace performs bounded live atomic strict search over anonymous Minimal Binary Type Theory (MBTT) telescopes and recovers a deterministic 15-step trajectory through step~15. The hot path is structural and name-free: admissibility, atomic candidate generation, exact evaluation, semantic minimality, and minimal-positive-overshoot selection are all computed without target-conditioned semantic rewards. The engine also writes deterministic run artifacts, frontier-backed resume evidence, and observational Agda exports. At the \emph{semantic-interpretation layer}, labels such as cohesion, metric, Hilbert, or DCT are post-hoc readings of selected AST skeletons rather than targets encoded in the search. Under that interpretation, the sequence bootstraps minimal type theory, ascends through core homotopy theory, reaches modal and differential API shells, and halts at a temporal-cohesive shell whose semantic completion can be read as the Dynamical Cohesive Topos (DCT). The machine-verified strict canon reported here is therefore the output of the current bounded live atomic artifact under a fixed shared vocabulary and fixed admissibility policy: the accepted library starts empty, but the signature does not. The claim is conditional structural selection from that fixed expressive basis, not invention of modal or temporal primitives from scratch.
\end{abstract}

\tableofcontents
\clearpage


% ============================================
% INTRODUCTION
% ============================================
\section{Introduction}
\label{sec:introduction}

Automated discovery in formal mathematics has two coupled problems: how to \emph{generate} plausible typed candidates, and how to \emph{select} among them without target-conditioned scoring, theorem injection, or retrospective evaluator tuning.
This paper presents an executable answer via the \emph{Principle of Efficient Novelty} (PEN): from an initially empty accepted library over a fixed intensional type-theoretic signature, extend by the admissible candidate that clears a history-dependent efficiency bar with minimal positive overshoot, where
\[
  \rho = \nu / \kappa
\]
and $\nu$ counts marginal derivational expansion while $\kappa$ counts irreducible specification clauses.

The key point is algorithmic: candidate search, scoring, and selection are implemented as one reproducible loop.
The current executable path is MBTT-first, atomic, and type-directed over well-typed anonymous ASTs, with live exact-band search through step~15 in the currently supported bounded lane; semantic names are assigned only post hoc for interpretability.
The resulting trajectory is deterministic under fixed settings and yields a specific sequence of geometric type-theoretic structures.

\paragraph{Three claim layers.}
The paper distinguishes three layers that should not be conflated.
\emph{Executable finding:} the current Rust artifact recovers a specific 15-step trajectory in a bounded live atomic strict lane under bounded search.
\emph{PEN-level structural theorem:} under PEN's fixed admissibility policy, cost model, and fully coupled CCHM/HoTT lane, the coherence dynamics force Fibonacci integration debt and privilege high-interaction geometric structure over discrete alternatives.
\emph{Post-hoc semantic interpretation:} the engine searches over MBTT ASTs, while names such as ``metric,'' ``Hilbert,'' or ``DCT'' are human semantic readings of the selected structural skeletons rather than semantic targets supplied to the evaluator.
Late-stage discovery in the current executable artifact therefore means objective filtering over a bounded live atomic MBTT lane, not theorem-targeted search or open-ended invention of a new mathematical vocabulary.

This differs from the dominant theorem-prover use pattern.
Interactive provers and ATP systems typically consume a human-specified conjecture and produce or check a proof object for that target.
PEN instead treats \emph{library growth} itself as the optimization target: the engine generates admissible structural extensions, scores them with a fixed objective, and selects without goal-specific theorem injection in the evaluator.
The mathematically central claim is therefore metatheoretic rather than anthropomorphic: under maximal interface density, a fixed MBTT vocabulary, and the fully coupled CCHM/HoTT lane studied here, the proof-theoretic growth dynamics force a narrow efficiency ordering in which cohesive and differential structure dominate discrete alternatives.
The current Rust atomic engine is the computational verification of this claim, not an argument that the system ``invented'' its own vocabulary.

\paragraph{Fixed shared vocabulary.}
The search does \emph{not} begin from a semantically empty token inventory.
The shared MBTT grammar already fixes foundational, modal, and temporal primitives, including $\flat$, $\sharp$, $\Disc$, $\Pi_{\!\infty}$, $\bigcirc$, and $\Diamond$.
The current live lane is nevertheless bounded: it searches only the currently admissible region of the anonymous MBTT space under fixed exact-band controls, windowed connectivity requirements, and deterministic retention policies.
The defensible claim is narrower: PEN gives these symbols no target-conditioned reward, theorem hint, or hand-written selection bonus.
They survive only if typed combinations built from that fixed vocabulary clear the same structural objective as every other admissible candidate.

\paragraph{Vocabulary Dependence.}
For reviewer-facing clarity, the manuscript's ``empty library'' language means that the \emph{accepted library state} starts empty; it does \emph{not} mean that the underlying signature is devoid of advanced primitives.
The current MBTT grammar already contains cohesive and temporal operators, so the executable claim is not ex nihilo invention of time or cohesion from the void.
It is a narrower and still substantive claim: under a fixed expressive basis, fixed admissibility policy, and fixed structural objective, the engine recovers a specific trajectory without target-conditioned scoring, theorem injection, or hand-written rewards for the late semantic readings.
Different signatures could in principle yield different trajectories; the contribution here is a conditional structural-selection result over the disclosed basis, not a claim that the system independently originated the modal vocabulary itself.

\paragraph{Conjecture versus finding.}
The paper makes one conjectural claim and one implemented claim, and they should not be conflated.
\emph{Conjecture:} with sufficiently large compute and materially weaker pruning, the same PEN objective over the same admissible MBTT space would recover the same 15-step structural trajectory, including the same temporal-cohesive endpoint later semantically read as DCT.
\emph{Finding:} the current implementation recovers that trajectory in a bounded live atomic strict lane suitable for standard desktop hardware by combining exact-band admissibility, atomic MBTT enumeration, hybrid novelty evaluation, semantic minimality filtering, and deterministic minimal-overshoot selection.
These tractability devices are part of the implementation status report, not part of the mathematical theorem.
Legacy structural audits, donor Haskell artifacts, and reference telescopes remain diagnostics or provenance only and do not override the strict canon reported in this manuscript.

The key structural input is the \emph{Coherence Window}---the depth of historical context required to stabilize irreducible multi-layer coherence obligations when a new structure is sealed against the library.
We prove that intensional type theory, in the fully coupled CCHM/HoTT lane studied here, has Coherence Window $d = 2$: sealing a new type requires interaction with the two most recent integration layers, but not deeper.
This two-step memory forces the integration cost to obey the Fibonacci recurrence $\Delta_{n+1} = \Delta_n + \Delta_{n-1}$, so the selection threshold rises at the rate of the Golden Ratio $\varphi \approx 1.618$.

Against this exponentially rising bar, only structures with \emph{superlinear} novelty growth can survive.
Discrete structures---natural numbers, ordinals, classical logic---have at most linear novelty ($\nu_H = 0$, no path constructors), so they are permanently rejected once the bar exceeds their efficiency; \S\ref{sec:exponentiality} promotes this to a general asymptotic proposition for $0$-groupoid-style candidates.
Geometric structures, by contrast, carry nontrivial homotopy content: path constructors generate $\mathcal{O}(d^2)$ independent computation rules, producing the superlinear novelty needed to outpace the Fibonacci threshold.
The algorithmic objective $\rho$ thus inherently biases the search away from discrete structures, mathematically demonstrating why continuous, topological frameworks naturally dominate complexity-bounded synthesis.

Starting from an initially empty accepted library under this fixed signature, the selection process produces a 15-step \emph{PEN Synthesis Trajectory}:
\begin{enumerate}[nosep]
    \item \textbf{Bootstrap} (steps 1--4): A universe, unit type, witness, and dependent function/sum types---the minimal infrastructure of type theory.
    \item \textbf{Geometric Ascent} (steps 5--9): The circle $S^1$, propositional truncation, the sphere $S^2$, the H-space $S^3$, and the Hopf fibration---the core objects of homotopy theory.
    \item \textbf{Framework Abstraction} (steps 10--14): Cohesion, connections, curvature, an endomorphic operator bundle (metric-equivalent semantic reading), and Hilbert-functional structure.
    \item \textbf{Synthesis} (step 15): A temporal-cohesive shell, later semantically completed as the Dynamical Cohesive Topos (DCT), fuses spatial logic with Nakano temporal modalities and explicit exchange axioms, clears the selection bar with the largest absolute overshoot in the sequence ($\rho = 12.88$, bar $= 7.40$), and triggers PEN's local halting condition.
\end{enumerate}
Termination is not a failure of the search: the Step-15 temporal-cohesive shell internalizes continuous evolution on the connected component of the Univalent Universe, so every subsequent continuous proposal has zero marginal novelty up to equivalence, while disconnected axiomatic jumps cannot amortize the required Step-16 interface debt $\Delta_{16}=F_{16}=987$ under PEN's bar.
We call this boundary the \emph{Univalent Horizon}---the connected-component boundary of the type-theoretic moduli space beyond which only external axiomatic jumps can reach.

The model produces several structural results:
\begin{itemize}[nosep]
    \item The \textbf{Coherence Window Theorems} establish $d = 1$ for extensional type theory (constant costs, stagnation) and $d = 2$ for the fully coupled CCHM/HoTT intensional lane analyzed here (Fibonacci costs, sustained evolution), via three independent arguments from categorical, homological, and topological perspectives.
    \item The \textbf{Complexity Scaling Theorem} shows that this $d = 2$ coherence regime forces $\Delta_n = F_n$ (the Fibonacci sequence), making the Golden Ratio $\varphi$ the dominant eigenvalue of mathematical evolution.
    \item The \textbf{Combinatorial Novelty Theorem} proves that novelty grows superlinearly, ensuring that efficiency permanently outpaces the rising selection bar.
    \item The \textbf{Terminal Fixed-Point Analysis} shows that the temporal-cohesive endpoint, whose semantic completion is read as DCT, is selected with a large deterministic overshoot in the current bounded live atomic lane and realizes the unique high-interaction structural fixed point available inside the connected component explored by the engine.
\end{itemize}
Computational verification is provided by the current Rust atomic engine, which recovers the 15-step sequence in its current bounded live strict lane.
The manuscript's canonical values come from that strict run; auxiliary audits are used only as consistency checks and never override the strict canon.
Formal proofs in Cubical Agda verify the Fibonacci recurrence and the abstraction barrier at steps 8--9.

\paragraph{Scope and limitations.}
A common objection to automated mathematical discovery is teleology---that target structures were smuggled into candidate generation.
Our current implementation addresses this partly by making MBTT-first typed AST synthesis the default runtime path, but the search basis is \emph{not} semantically empty: the shared grammar already contains foundational, modal, and temporal primitives, and the live lane remains bounded by explicit admissibility and runtime controls.
The strict implementation still uses bounded search and explicit admissibility controls, and its main discovery path is the bounded live atomic lane described in \S\ref{sec:verification}.
The PEN evaluator remains structural: it computes efficiency mechanically from the candidate and library state, and the names in Table~\ref{tab:genesis} are post-hoc semantic identifications of selected syntactic winners.
The search therefore uses a fixed shared vocabulary rather than an unbiased empty grammar; what is absent is target-conditioned scoring, theorem injection, and semantic rewards for any particular modal or geometric label.
PEN therefore derives a concrete trajectory of geometric type-theoretic library extensions, not a claim of global uniqueness across all logics or all objective functions.
The model still assumes \emph{maximal interface density} (each candidate seals against the exported interface of the past $d$ layers).
For foundational constructors and modalities, that burden is tied to canonicity preservation: without explicit interaction clauses against prior eliminators, normalization can get stuck and the $d=2$ Fibonacci latency theorem loses its logical basis.
For late API specifications, the same burden is enforced by an economic-topological squeeze: an uncoupled API cannot clear PEN's bar, while any profitable interaction with the univalent library must extend functorially across the full $2$-dimensional boundary of the active interface.

\paragraph{Main contributions.}
\begin{enumerate}[nosep]
    \item A formal selection model (PEN) with explicit axioms, costs, and admissibility.
    \item An executable strict MBTT-first synthesis algorithm implementing the model end-to-end.
    \item Structural theorems linking coherence depth to Fibonacci integration dynamics and bar growth.
    \item A deterministic 15-step PEN trajectory plus auditable omitted alternatives under the same scoring rule, together with a bounded live atomic recovery of that trajectory in the current implementation.
    \item Open, reproducible artifact support: source, scripts, and instructions for independent reruns.
\end{enumerate}

\paragraph{Outline.}
\Cref{sec:related} situates PEN against AI theorem proving, automated discovery systems, and HoTT/univalent foundations.
\Cref{sec:genesis} presents the complete PEN Synthesis Trajectory and the patterns that govern it.
\Cref{sec:framework} defines the model: state, candidates, the dual-cost framework, Generative Capacity, and the five axioms of selection dynamics.
\Cref{sec:coherence} proves the Coherence Window Theorems ($d = 1$ for extensional, $d = 2$ for the fully coupled CCHM/HoTT intensional lane studied here).
\Cref{sec:scaling} derives the Complexity Scaling Theorem ($\Delta_n = F_n$).
\Cref{sec:exponentiality} establishes superlinear novelty growth, the DCT synthesis mechanism, and the Univalent Horizon.
\Cref{sec:decomposition} consolidates the DCT fixed-point evidence and the mathematical properties that make DCT the terminal endpoint under PEN.
\Cref{sec:verification} describes the computational verification: the Rust atomic engine, the bounded live evaluation loop, auxiliary diagnostics, rejected candidates, and Cubical Agda mechanization.
\Cref{sec:discussion} assesses what is proved, what is assumed, and what remains open.
\Cref{sec:boundaries} collects three boundary clarifications---algorithmic criticality, the Univalent Horizon, and the status of the DCT novelty accounting---plus a short supplementary note linking back to the main-text discrete rejection argument.

% ============================================
% SECTION 2: RELATED WORK
% ============================================
\section{Related Work}
\label{sec:related}

\paragraph{AI theorem proving.}
Recent systems such as AlphaGeometry~\cite{alphageometry} and Lean-centered neural proving pipelines~\cite{leandojo,openai-lean} show strong performance on goal-conditioned proving tasks.
Their core objective is to prove or refute a \emph{given} statement, typically via search policies guided by learned value/policy estimates over proof states.
PEN differs in objective: it has no externally supplied target theorem and instead optimizes library extension under a fixed structural utility $\rho=\nu/\kappa$.

\paragraph{Theory exploration and lemma discovery.}
QuickSpec~\cite{quickspec}, IsaCoSy~\cite{isacosy}, and Hipster~\cite{hipster} search for equations or lemmas inside a pre-existing theory.
Their success criterion is explanatory power over a fixed signature: infer useful equalities, conjectures, or proof obligations already latent in the background logic.
PEN instead scores candidate \emph{extensions of the background logic/library itself}; its output is a typed API specification together with a proof-theoretic cost/benefit analysis, not a conjecture set.

\paragraph{Program synthesis and type-directed search.}
Synquid~\cite{synquid} and Agda's auto/Agsy-style proof search inhabit user-specified types or refinement signatures.
PEN borrows typed search discipline from that tradition, but changes the optimization target: the system is not asked to inhabit a given type, but to choose which new type former, modality package, or API should enter the library.

\paragraph{Automated mathematical discovery.}
Classical theory-formation systems (e.g., AM/HR lines of work~\cite{lenat-am,colton-atf}) generate concepts and conjectures from heuristic interestingness criteria.
PEN is closer to constructive library synthesis: candidates are typed MBTT specifications, admissibility is type-theoretic, and acceptance is determined by explicit dual-cost dynamics with a history-dependent bar.
The resulting object is an auditable trajectory of realized structures, not just a conjecture list.

\paragraph{Type-theoretic foundations.}
Our setting is intensional HoTT with univalence~\cite{hott}, implemented in a cubical ecosystem~\cite{cubical,cubical-agda}.
This foundation is essential for PEN's novelty calculus because higher-path structure contributes directly to computation-rule growth.
On the metamathematical side, our use of intensionality is aligned with work showing that types support weak $\omega$-groupoid structure~\cite{lumsdaine,vdberg-garner}; PEN differs by turning that higher-dimensional coherence into an explicit cost model for foundational library growth.

% ============================================
% SECTION 3: THE PEN SYNTHESIS TRAJECTORY
% ============================================
\section{The PEN Synthesis Trajectory}
\label{sec:genesis}

\Cref{tab:genesis} displays the complete output of the Principle of Efficient Novelty: a deterministic emergence path of geometric type-theoretic structure from an initially empty accepted library under a fixed intensional signature.
The model operates on abstract Obligation Graphs rather than syntactic artifacts; it generates candidate structures, computes their integration costs, and selects the least supercritical admissible candidate.
The reader should treat the table in two ways: as the mathematical prediction produced by the PEN model, and as the trajectory that the current bounded live atomic implementation recovers under bounded desktop-oriented search policies.

\subsection{Syntactic Skeletons vs.\ Semantic Interpretation}
\label{sec:semantic-firewall}

The manuscript uses a strict firewall between selected syntax, PEN-level theorems, and later semantic interpretation.
At the \emph{syntax layer}, the engine searches over anonymous MBTT ASTs and ranks them by the structural objective $\rho = \nu/\kappa$.
At the \emph{PEN layer}, the theorems concern admissibility, coherence windows, Fibonacci integration debt, and bar-clearing order under a fixed evaluator.
At the \emph{semantic layer}, names in \cref{tab:genesis} are post-hoc human interpretations of the selected AST skeletons.

This separation matters most for the late steps.
PEN does \emph{not} claim that the raw AST in \cref{tab:step13} already contains the full semantic machinery of Riemannian geometry or that the temporal shell in \cref{tab:step15} directly presents a fully developed topos.
The executable finding is narrower: the engine selects compact, highly coupled typed API skeletons that are natural hosts for those interpretations once one passes to an appropriate semantic completion.

Accordingly, \cref{tab:genesis} should be read with three constraints in mind:
\begin{enumerate}[nosep]
    \item \textbf{What PEN derives:} A deterministic sequence of realizable library extensions under fixed axioms, admissibility, vocabulary, and ranking.
    \item \textbf{What PEN does not derive:} A semantic theorem that every selected AST already contains the full body of the human mathematical concept named in the late-row glosses.
    \item \textbf{What the implementation establishes:} Recovery of this trajectory in the current bounded live atomic lane, not exhaustive search over the full admissible MBTT space.
\end{enumerate}
Late rows should therefore be read shell-first: the quantitative canon attaches to the selected AST skeleton, while labels such as ``metric'' and ``DCT'' name semantic completions discussed later in the paper.

\begin{definition}[Three structural components of novelty]
\label{def:nu-components-preview}
For any admissible candidate $X$, we write
\[
  \nu(X) = \nu_G(X) + \nu_C(X) + \nu_H(X),
\]
where $\nu_G$ counts formation/introduction schemas that create new inhabitants, $\nu_C$ counts elimination and cross-interface schemas induced by the candidate's API, and $\nu_H$ counts computation, Kan, and path-reduction schemas.
This notation is used throughout \cref{tab:genesis} and receives its operational interpretation in \S\ref{sec:framework}.
\end{definition}

\begin{table}[tbp]
\centering
\caption{The PEN Synthesis Trajectory.  $\nu$ is the Generative Capacity (count of atomic inference rules added to the derivation logic).  The quantitative canon in this table is the machine-verified output of the current bounded live atomic strict engine (\S\ref{sec:verification}).  Legacy structural and uniform diagnostics are discussed later only as secondary cross-checks and do not override these values.  Rows 10--15 should be read as selected AST skeletons with post-hoc semantic readings, not as literal theorem-level presentations of the named human theories; in particular, Step~13 receives the metric-equivalent reading and Step~15 receives the semantic DCT completion.}
\label{tab:genesis}
\small
\setlength{\tabcolsep}{3pt}
\begin{tabular}{@{}cr l rrrr rrr@{}}
\toprule
$n$ & $\tau$ & Selected AST skeleton & $\Delta_n$ & $\nu$ & $\kappa$ & $\rho$ & $\Phi_n$ & $\Omega_{n-1}$ & Bar \\
\midrule
1  & 1    & Universe $\U_0$              & 1   & 1   & 2 & 0.50  & ---  & ---  & ---  \\
2  & 2    & Unit type $\mathbf{1}$       & 1   & 1   & 1 & 1.00  & 1.00 & 0.50 & 0.50 \\
3  & 4    & Witness $\star : \mathbf{1}$ & 2   & 2   & 1 & 2.00  & 2.00 & 0.67 & 1.33 \\
4  & 7    & $\Pi$/$\Sigma$ types         & 3   & 5   & 3 & 1.67  & 1.50 & 1.00 & 1.50 \\
\midrule
5  & 12   & Circle $S^1$                 & 5   & 7   & 3 & 2.33  & 1.67 & 1.29 & 2.14 \\
6  & 20   & Propositional truncation     & 8   & 8   & 3 & 2.67  & 1.60 & 1.60 & 2.56 \\
7  & 33   & Sphere $S^2$                 & 13  & 10  & 3 & 3.33  & 1.62 & 1.85 & 3.00 \\
8  & 54   & H-space $S^3$                & 21  & 18  & 5 & 3.60  & 1.62 & 2.12 & 3.43 \\
9  & 88   & Hopf fibration               & 34  & 17  & 4 & 4.25  & 1.62 & 2.48 & 4.01 \\
\midrule
10 & 143  & Modal shell (cohesion reading)              & 55  & 19  & 4 & 4.75  & 1.62 & 2.76 & 4.46 \\
11 & 232  & Connection shell                            & 89  & 26  & 5 & 5.20  & 1.62 & 3.03 & 4.91 \\
12 & 376  & Curvature shell                             & 144 & 34  & 6 & 5.67  & 1.62 & 3.35 & 5.42 \\
13 & 609  & Endomorphic operator bundle (metric reading) & 233 & 46  & 7 & 6.57  & 1.62 & 3.70 & 5.99 \\
14 & 986  & Hilbert-functional shell                    & 377 & 62  & 9 & 6.89  & 1.62 & 4.13 & 6.68 \\
\midrule
15 & 1596 & Temporal-cohesive shell (semantic DCT completion) & 610 & 103 & 8 & 12.88 & 1.62 & 4.57 & 7.40 \\
\bottomrule
\end{tabular}
\end{table}

\subsection{Reading the Table}

Each row records a realized extension together with its integration debt $\Delta_n$, novelty $\nu$, clause cost $\kappa$, efficiency $\rho=\nu/\kappa$, and active bar $\Phi_n\Omega_{n-1}$.
Three features matter for the rest of the paper:
\begin{enumerate}[nosep]
    \item \textbf{Fibonacci timing.} $\Delta_n$ follows $F_n$ and $\tau_n=F_{n+2}-1$.
    \item \textbf{Selective survival.} Every realized structure clears the bar, with the tightest margin at Step~6.
    \item \textbf{Four phases.} The trajectory separates into bootstrap, geometric ascent, framework abstraction, and terminal synthesis.
    \item \textbf{Shell-first late rows.} Steps~10--15 are named by selected AST skeletons first; their human mathematical names are semantic glosses supplied only after selection.
\end{enumerate}
A worked arithmetic expansion of Steps~5--6 and an expanded reading guide are deferred to \cref{app:worked-arithmetic}.

\subsection{The Efficiency Peak}

The fifteenth selected shell is temporal-cohesive; DCT is its semantic completion.
Syntactically, the selected AST combines spatial logic (cohesion) with temporal logic (LTL~\cite{nakano}) and explicit exchange laws only; the infinitesimal reading is derived later as a semantic consequence rather than postulated as a primitive clause.
Its executable efficiency is $\rho = 12.88$, clearing the Step-15 bar of $7.40$ by an absolute overshoot of $5.48$, the largest in the sequence.
The current evaluator reaches this value by scoring the temporal shell at conceptual cost $\kappa = 8$ and measuring its mechanically available temporal-cohesive interaction payload against the existing library.
The manuscript canon is the executable value reported in \cref{tab:genesis}; explanatory rule-group reconstructions later in the paper are subordinate to that total.
We define the Step-15 AST and its semantic DCT reading in \S\ref{sec:exponentiality} and detail the computational verification in \S\ref{sec:verification}.

After this temporal-cohesive step, no candidate in PEN's strict fully coupled lane can clear the positive Step-16 bar: internal continuations contribute zero marginal novelty up to equivalence, while disconnected jumps face the full interface debt $\Delta_{16}=987$.
The sequence halts algorithmically.

% ============================================
% SECTION 4: THE MODEL
% ============================================
\section{The Model}
\label{sec:framework}

We model automated emergence of geometric type-theoretic structure as a discrete-time optimization process operating on a state $\mathcal{B}$ (the ``Library'').
At each step, the system generates candidate extensions, calculates their \emph{Efficiency} $\rho = \nu / \kappa$, and selects the optimal candidate.

\subsection{State and Candidates}

\begin{definition}[State]
\label{def:state}
A \emph{State} $\mathcal{B}$ is a monotone context closed under derivability.
An evolution step $\mathcal{B}_n \leadsto \mathcal{B}_{n+1}$ is an extension by a single sealed structure.
\end{definition}

\begin{definition}[Candidate]
\label{def:candidate}
A \emph{Candidate} $X$ is a pair $(X_{\mathrm{core}}, \mathcal{G}_{\mathrm{obl}})$, where $X_{\mathrm{core}}$ is the definitive data (type formers, constructors) and $\mathcal{G}_{\mathrm{obl}}$ is the \emph{Obligation Graph}: the set of atomic coherence obligations required to seal $X$ against the history.
\end{definition}

\subsection{The Dual-Cost Model}

\begin{definition}[Integration Latency]
\label{def:latency}
The \emph{Integration Latency} $\Delta(X \mid \mathcal{B}) := |\mathcal{G}_{\mathrm{obl}}|$ counts the coherence witnesses required to seal $X$ against the library.
\end{definition}

\begin{definition}[Construction Effort]
\label{def:effort}
A \emph{specification} $\mathcal{S}$ of candidate~$X$ over library~$\mathcal{B}$ is a set of atomic statements (type formations, term introductions, equations) that determines the inference rules of~$X$ relative to~$\mathcal{B}$.
The \emph{Construction Effort} of a specification is its clause count:
\begin{equation}
    \kappa(\mathcal{S}) := |\mathcal{S}|
\end{equation}
Different specifications of the same mathematical type may include different amounts of structure (e.g., a bare higher inductive type vs.\ one equipped with algebraic operations), yielding different Generative Capacities $\nu(\mathcal{S})$ and hence different efficiencies $\rho(\mathcal{S}) = \nu(\mathcal{S}) / \kappa(\mathcal{S})$.
The selection mechanism (Axiom~\ref{ax:selection}) jointly evaluates all admissible specifications.

Each clause is represented as a term in a prefix-free binary encoding (\emph{Minimal Binary Type Theory}, MBTT; see \cref{tab:mbtt-grammar-main,sec:mbtt-audit}), so the grammar used for clause accounting is part of the model rather than appendix-only notation.
Referencing an existing structure $L_i \in \mathcal{B}$ incurs an information-theoretic pointer cost of $O(\log_2 i)$ bits via Elias gamma coding.
The clause count $\kappa$ is primary; the MBTT bit-length is an audit/tiebreak layer used only when two specifications yield the same efficiency~$\rho$.
\end{definition}

\begin{definition}[Definitional derivability vs.\ constructive irreducibility]
\label{def:derivability}
Fix a prior library~$\mathcal{B}$ with operational semantics $\rightsquigarrow_{\mathcal{B}}$ generated by its existing eliminators, $\beta$-rules, and judgmental equalities.
For an operation symbol $o$ proposed over~$\mathcal{B}$:
\begin{enumerate}[nosep]
    \item $o$ is \emph{definitionally derivable} if there exists a closed term $t_o$ in the signature of~$\mathcal{B}$ such that $t_o \rightsquigarrow_{\mathcal{B}}^{*} v_o$ to a canonical value and all required computation clauses of $o$ are inherited from the existing eliminator fragment.  In that case $o$ contributes $\kappa_{\mathcal{B}}(o)=0$.
    \item $o$ is \emph{constructively irreducible} if every term intended to realize its specification contains a new neutral head or a stuck eliminator relative to $\rightsquigarrow_{\mathcal{B}}$; equivalently, orienting the intended behavior requires adding at least one new clause to the operational semantics.  In that case $o$ contributes $\kappa_{\mathcal{B}}(o)>0$.
\end{enumerate}
The intended contrast is standard: addition on $\mathbb{N}$ is derivable from $\mathsf{rec}_{\mathbb{N}}$, whereas a Levi-Civita connection or Hodge star requires new computation/API clauses not reducible in the prior library.
\end{definition}

The MBTT node inventory used for clause accounting is fixed in \cref{tab:mbtt-grammar-main}.
Primitive nodes for $\flat$, $\sharp$, $\bigcirc$, and $\Diamond$ are ambient vocabulary in that shared grammar, not target-conditioned rewards: they still incur explicit local cost and are selected only if they clear the same structural objective as every other admissible combination.

\begin{table}[tbp]
\centering
\caption{MBTT grammar used in the model for clause accounting and bit-cost audit.}
\label{tab:mbtt-grammar-main}
\scriptsize
\begin{tabular}{@{}llrl@{}}
\toprule
Node & Prefix & Bits & Description \\
\midrule
\multicolumn{4}{@{}l}{\textbf{Core constructors}} \\
\textsc{App}$(f, x)$   & \texttt{00}    & 2 & Application \\
\textsc{Lam}$(b)$      & \texttt{01}    & 2 & Abstraction \\
\textsc{Pi}$(A, B)$    & \texttt{100}   & 3 & Dependent function type \\
\textsc{Sigma}$(A, B)$ & \texttt{1010}  & 4 & Dependent sum type \\
\textsc{Univ}          & \texttt{1011}  & 4 & Universe $\U_0$ \\
\midrule
\multicolumn{4}{@{}l}{\textbf{Variables and references}} \\
\textsc{Var}$(i)$      & \texttt{110} + $\gamma(i)$ & $3 + |{\gamma(i)}|$ & Bound variable \\
\textsc{Lib}$(i)$      & \texttt{111} + $\gamma(i)$ & $3 + |{\gamma(i)}|$ & Library pointer \\
\midrule
\multicolumn{4}{@{}l}{\textbf{HoTT extensions}} \\
\textsc{Id}$(A,x,y)$  & \texttt{11100}   & 5 & Identity type \\
\textsc{Refl}$(a)$    & \texttt{11101}   & 5 & Reflexivity \\
\textsc{Susp}$(A)$    & \texttt{11110}   & 5 & Suspension \\
\textsc{Trunc}$(A)$   & \texttt{111110}  & 6 & Propositional truncation \\
\textsc{PathCon}$(d)$ & \texttt{111111} + $\gamma(d)$ & $6 + |{\gamma(d)}|$ & Path constructor (dim $d$) \\
\midrule
\multicolumn{4}{@{}l}{\textbf{Modal and temporal operators}} \\
$\flat(A)$             & \texttt{1111100}   & 7 & Flat modality \\
$\sharp(A)$            & \texttt{1111101}   & 7 & Sharp modality \\
\textsc{Disc}$(A)$     & \texttt{1111110}   & 7 & Discrete reflector \\
\textsc{Shape}$(A)$    & \texttt{11111110}  & 8 & Cohesive shape \\
$\bigcirc(A)$          & \texttt{111111110} & 9 & Next modality \\
$\Diamond(A)$          & \texttt{111111111} & 9 & Eventually modality \\
\bottomrule
\end{tabular}
\end{table}

\begin{remark}[Kolmogorov invariance of the encoding bias]
\label{rem:kappa-invariance}
The exact binary description length of a specification depends on the chosen prefix-free universal grammar, here MBTT.
By the Invariance Theorem of Kolmogorov Complexity~\cite{kolmogorov}, replacing MBTT by any other universal grammar changes that description length by at most an additive constant $O(1)$.
PEN uses clause count as the coarse primary denominator and MBTT bit-length as an audit/tiebreak layer, so the encoding choice contributes only bounded bias at the granularity used here.
Because the $d=2$ survival bar is later shown to grow with the Fibonacci debt sequence (\cref{cor:fibonacci}), this $O(1)$ grammar dependence can affect only finitely many low-complexity edge cases and is asymptotically negligible.
\end{remark}

\begin{definition}[API dependency graph and structural bundle]
\label{def:api-graph}
Let $\mathcal{S}=\{c_1,\dots,c_r\}$ be a candidate specification over~$\mathcal{B}$.
Its \emph{API dependency graph} $G_{\mathcal{S}}$ has one vertex per clause and a directed edge $c_i \to c_j$ when the type, defining equation, or computation rule of $c_i$ mentions a symbol introduced by~$c_j$.
A \emph{structural bundle} is an induced strongly connected component $C \subseteq G_{\mathcal{S}}$ that is closed under well-typedness: every symbol referenced by a clause of $C$ is either inherited from~$\mathcal{B}$ or introduced by another clause of~$C$.
\end{definition}

\begin{definition}[Minimal Complete API]
\label{def:minimal-api}
For a library specification at steps~10--14, a \emph{Minimal Complete API} is a specification $\mathcal{S}$ satisfying:
\begin{enumerate}[nosep]
    \item \emph{Irreducibility:} every clause of $\mathcal{S}$ is constructively irreducible in the sense of \cref{def:derivability}.
    \item \emph{Completeness:} every irreducible operation required for the package's intended computation rules is included.
    \item \emph{Structural unity:} $\mathcal{S}$ forms a single structural bundle in the sense of \cref{def:api-graph}.
    \item \emph{Minimality:} no proper structural sub-bundle of $\mathcal{S}$ is simultaneously well-typed, complete, and able to clear the active selection bar.
\end{enumerate}
This criterion immunizes the clause count~$\kappa$ against accusations of cherry-picking: the boundary between ``derivable definition'' ($\nu = 0$) and ``irreducible API clause'' ($\nu > 0$) is the type-theoretic constructibility of the operation from the prior library (see \cref{app:metric-spec} for the detailed application to the Metric at Step~13).
\end{definition}

\begin{lemma}[SCC minimality of the Step-13 metric package]
\label{lem:metric-scc}
The seven Step-13 metric clauses form a structural bundle whose proper sub-bundles either leave referenced operators undefined or satisfy $\rho \le 3.0 < 5.99$.
In particular, the bare tensor clause $g : TX \otimes_s TX \to \mathbb{R}$ has $\kappa=1$ and fails the Step-13 bar, while operators such as $\Delta_g$ and $\mathrm{Ric}_g$ depend on the Levi-Civita and curvature clauses and therefore cannot be split off as independent admissible candidates.
\end{lemma}

\begin{proof}[Proof sketch]
Edges in the dependency graph record the chain $g \to \nabla_g \to R \to \mathrm{Ric}_g$ together with the shared dependence of $\star_g$ and $\Delta_g$ on the metric-induced volume and connection data.
Removing downstream clauses breaks well-typedness or leaves required symbols uninterpreted; removing all but the tensor clause yields the stripped specification analyzed in \cref{app:metric-spec}, whose efficiency is at most~$3.0$.
Hence survival at Step~13 requires the super-additive cross-interactions generated by the full bundle rather than any proper fragment.
\end{proof}

\begin{axiom}[Constructive Irreducibility Boundary]
\label{ax:constructive-irreducibility}
Let $\mathcal{S}$ be a candidate specification over a prior library~$\mathcal{B}$.
\begin{enumerate}[nosep]
    \item If an operation is definitionally derivable in the sense of \cref{def:derivability}, i.e.\ represented by a closed term reducing under $\rightsquigarrow_{\mathcal{B}}$ via existing eliminators, it must be excluded from~$\mathcal{S}$ and contributes $\nu=0$ as derived structure.
    \item If realizing the operation would otherwise leave a new neutral term or stuck eliminator relative to $\rightsquigarrow_{\mathcal{B}}$, then the operation is constructively irreducible and must appear explicitly as an API clause in~$\mathcal{S}$.
\end{enumerate}
This boundary is mandatory for all steps and removes curator discretion in clause inclusion: $\N$-addition is excluded, whereas Levi-Civita/Hodge/Laplacian clauses are required.
\end{axiom}

\begin{remark}[Specification selection and the $S^3$ example]
\label{rem:kolmogorov-ambiguity}
Different specifications of the same type compete within the selection dynamics.
Consider $S^3$ at step~8, where the bar is $3.43$:
\begin{center}
\small
\begin{tabular}{@{}lrrrl@{}}
\toprule
Specification & $\kappa$ & $\nu$ & $\rho$ & Overshoot \\
\midrule
Bare 3-sphere ($\Sigma S^2$) & 1 & 15 & 15.00 & 11.57 \\
Native HIT (3-cell)          & 3 & 15 &  5.00 &  1.57 \\
H-space $S^3$ (with multiplication)    & 5 & 18 &  3.60 &  0.17 \\
\bottomrule
\end{tabular}
\end{center}
The bare 3-sphere (via suspension or native HIT) contributes 15~inference rules ($\nu_G = 5$, $\nu_H = 10$, $\nu_C = 0$).
The H-space specification additionally equips $S^3$ with multiplication and unit laws, yielding 3~extra Elimination rules ($\nu_C = 3$, total $\nu = 18$) at the cost of 2~additional specification clauses ($\kappa = 5$).
All three clear the bar, but minimal overshoot (\cref{ax:selection}) selects the H-space $S^3$ specification.
This is intentionally weaker than full Lie-group ($A_\infty$-level) axiomatization, and therefore more economical: the rising bar at step~8 selects the minimal multiplicative extension that still clears.
Read through \cref{def:primality}, the selected H-space package is the constructively prime multiplicative factor, whereas richer algebraic packings would introduce surplus structure not forced by the current debt.
\end{remark}

\begin{definition}[Interface Basis]
\label{def:interface}
For a foundation with Coherence Window $d$, the interface available for sealing $X_{n+1}$ is:
\begin{equation}
    I^{(d)}_n := \biguplus_{j=0}^{d-1} S(L_{n-j})
\end{equation}
where $S(L_k)$ denotes the schemas exported by integration layer $L_k$.
\end{definition}

\begin{theorem}[Integration Trace Principle]
\label{lem:trace}
Each integration layer $L_k$ exports exactly $|S(L_k)| = \Delta_k$ schemas.
These schemas are the \emph{integration trace}: the set of resolved obligations from sealing $X_k$.
Each resolved obligation becomes one opaque export (\emph{elimination duality}), and each export generates one obligation for the subsequent candidate (\emph{one-per-face correspondence}).
\end{theorem}

\begin{proof}
\emph{Linearity of Elimination.}
The recursor $\mathrm{rec}_X$ distributes over type formers: it acts independently on each constructor and each path constructor of~$X_k$.
Consequently, the behavior of $\mathrm{rec}_X$ is determined by an \emph{atomic basis}---the set of clauses, one per cell of~$X_k$---and the obligation graph $\mathcal{G}_{\mathrm{obl}}$ decomposes into independent atoms indexed by this basis.
For maps (as at step~9), the two functorial operations (pullback and postcomposition) play the role of $\mathrm{rec}_X$, distributing over domain and codomain cells respectively.

\emph{Context Extension Principle.}
The active interface $I^{(d)}_n$ forms a \emph{context telescope}: an ordered sequence of typed entries, each well-formed in the context of its predecessors.
Sealing $X_{n+1}$ against this telescope requires exactly one clause per entry:
\emph{at least one} by constructive completeness (the eliminator would be stuck on an unhandled entry),
\emph{at most one} by confluence (the clause is uniquely determined by the entry's type).
Hence $|\mathcal{G}_{\mathrm{obl}}| = |I^{(d)}_n|$.

\emph{Sealing Encapsulation} (\cref{rem:encapsulation}): resolved obligations become opaque exports of the sealed layer; future types interact with $L_k$'s interface, not with underlying layers.

\emph{Verification scope:} steps~3--9, uniform across HITs (steps~3--8) and maps (step~9, the Hopf fibration);
machine-checked abstraction barrier at steps~8--9 in Cubical Agda (\texttt{Saturation/\allowbreak AbstractionBarrier.agda}, \texttt{AbstractionBarrier9.agda}); see \S\ref{sec:verification}.
\end{proof}

\begin{theorem}[Canonicity preservation under maximal interface density]
\label{thm:canonicity-preservation}
Let $\mathcal{B}$ be a Fully Coupled Type Theory (FCTT) library and let $X$ be a new constructor, HIT-style former, or operational modality sealed against the active interface $I_n^{(d)}$.
Then the extension preserves global computational closure, strong normalization, and canonicity only if, for every eliminator exported by $I_n^{(d)}$, the specification of $X$ provides a defining interaction clause.
Equivalently, a depth-$d$ extension requires exactly $|I_n^{(d)}|$ sealing clauses.
\end{theorem}

\begin{proof}
Normalization reduces eliminators against constructor or computation clauses.
Suppose some eliminator $e$ exported by $I_n^{(d)}$ has no interaction clause with~$X$.
Then the extended theory contains a well-typed term $e_X(\vec{u})$ whose head symbol is neutral relative to the old reduction system and whose reduction cannot proceed, because no $\beta$- or computation rule orients the interaction.
This produces an irreducible neutral term that is not a canonical form of the target type, so canonicity fails; the same stuck term witnesses loss of global computational closure.

Conversely, if each eliminator in $I_n^{(d)}$ receives exactly one interaction clause, reduction on the active telescope is total: every eliminator-constructor encounter is oriented by a unique equation, so normalization and canonicity are preserved on the extension.
The ``at least one'' direction follows from the stuck-neutral argument, and the ``at most one'' direction follows from confluence, since multiple distinct interaction clauses for the same eliminator head would create critical pairs.
Hence the number of obligatory sealing clauses is exactly $|I_n^{(d)}|$.
\end{proof}

\begin{theorem}[Univalent functoriality and API integration debt]
\label{thm:univalent-functoriality}
For axiomatic and API library extensions in the 2D/univalent lane (steps~10--14 of the PEN trajectory), the Fibonacci integration debt
\[
  \Delta_{n+1} = \Delta_n + \Delta_{n-1}
\]
is not forced by foundational canonicity alone.
Rather, once an API candidate must couple to the prior library in order to clear the PEN bar, univalent functoriality and higher parametricity force its interaction data to extend over the full 2-dimensional boundary of the active interface.
In the $d=2$ regime, that boundary has size $\Delta_n + \Delta_{n-1}$, so the required sealing burden is again Fibonacci.
\end{theorem}

\begin{proof}
In an extensional 1-dimensional type theory, an API candidate $X$ can remain largely isolated: it may specify a local record or operation and ignore most of the prior library while paying only $\mathcal{O}(1)$ interface debt.
Within PEN's selection dynamics, however, such an isolated API generates negligible cross-interaction novelty.
The concrete Step-13 stress test makes this explicit: a bare metric tensor has $\nu_G = 1$ and $\nu_C \le 2$, hence $\rho \le 3.0$, far below the Step-13 bar $5.99$.
Thus any late API that survives the exponentially rising bar must couple actively to the existing geometric library in order to harvest the combinatorial amplification of $\nu_C$.

Once $X$ couples to a univalent higher library, the cost of that coupling is no longer discretionary.
The universe $\U$ and its geometric inhabitants are nontrivial higher groupoids, so any cross-structural operation must be functorial: it cannot act only on 0-cells (points), but must also extend coherently over 1-cells (paths/equivalences) and 2-cells (homotopies/squares).
Whether the ambient presentation is cubical, simplicial, HoTT-style, or higher observational, the checker enforces the same abstract requirement: transport witnesses, naturality squares, or extension fillers must be supplied for the exposed 2-dimensional boundary data of the active interface.
Without those witnesses, the proposed interaction is ill-typed, non-fibrant, or observationally incomplete and therefore cannot be accepted as a compiled library extension.

By \cref{thm:spectral-degeneracy}, the active independent coherence data truncate at dimension~2 in the lane studied here.
The exposed boundary is therefore parameterized exactly by the integration traces of the previous two layers, of sizes $\Delta_n$ and $\Delta_{n-1}$.
Hence the full cost of a profitable API interaction is
\[
  |I_n^{(2)}| = \Delta_n + \Delta_{n-1},
\]
and the same Fibonacci recurrence governing the foundational steps persists across the phase transition from constructors to API specifications.
\end{proof}

\begin{remark}[Resolution of the API coupling gap]
\label{rem:maximal-coupling}
Within the fully coupled PEN lane, maximal interface density is not a free heuristic knob.
For the foundational and HIT steps, full coupling is forced proof-theoretically by \cref{thm:canonicity-preservation}: omitting a prior eliminator interaction produces a stuck neutral term and breaks global canonicity.
For the API steps (10--14), full coupling is forced by a different but equally rigid squeeze captured by \cref{thm:univalent-functoriality}: an uncoupled API dies economically because its $\nu_C$ is too small, while a coupled API must pay the exact cost of covering the full 2-dimensional boundary of the active interface.

Human mathematics often avoids this squeeze by modularity and selective decoupling.
A mathematician can define a metric tensor without immediately specifying its interaction with every prior geometric package.
PEN deliberately models a monolithic, maximally dense foundational core optimized for structural efficiency.
Inside that regime, the objective function $\rho$ economically forces an API to interact, and the 2-dimensional logic topologically fixes the exact cost of that interaction.
This is why extensional MLTT stagnates under the same objective while the 2D/univalent lane exhibits the Fibonacci recurrence all the way from foundational constructors to API shells.
\end{remark}

\begin{lemma}[Latency Recurrence]
\label{lem:recurrence}
By the Integration Trace Principle (\cref{lem:trace}), each layer exports $|S(L_k)| = \Delta_k$ schemas, so:
\begin{equation}
    \Delta_{n+1} = \sum_{j=0}^{d-1} \Delta_{n-j}
\end{equation}
For $d = 2$: $\Delta_{n+1} = \Delta_n + \Delta_{n-1}$, i.e., $\Delta_n = F_n$.
\end{lemma}

\subsection{Novelty: Generative Capacity}

We define novelty as the expansion of the logical capacity of the library.

\begin{definition}[Generative Capacity]
\label{def:novelty}
Let $\mathcal{L}(\mathcal{B})$ be the set of derivation schemas---type-inhabitation patterns classifiable as Introduction, Elimination, or Computation rules---available in library $\mathcal{B}$.
The \emph{Generative Capacity} of a candidate $X$ is the marginal expansion of the library's operational repertoire:
\begin{equation}
    \nu(X \mid \mathcal{B}) \;:=\; |\mathcal{L}(\mathcal{B} \cup \{X\})| - |\mathcal{L}(\mathcal{B})|
\end{equation}
measured at coherence depth $d$.
In the current executable implementation, $\nu$ is operationalized by the strict evaluator described in \S\ref{sec:verification}: native structural novelty plus explicit topological witness payload where present.
\end{definition}

\begin{definition}[Formal components of novelty]
\label{def:nu-components}
For every admissible candidate $X$, the novelty measure decomposes as
\[
  \nu(X) = \nu_G(X) + \nu_C(X) + \nu_H(X),
\]
where:
\begin{enumerate}[nosep]
    \item $\nu_G$ counts new formation/introduction schemas.
    \item $\nu_C$ counts elimination, projection, and cross-interface schemas induced by the specification.
    \item $\nu_H$ counts computation schemas arising from path constructors, Kan operations, and their independent reductions.
\end{enumerate}
The current executable audit extracts these three components directly from the telescope AST; \cref{def:nu-components-preview} introduced the notation earlier because it is needed to read \cref{tab:genesis}.
\end{definition}

\begin{remark}[Library extensions vs.\ foundational primitives]
\label{rem:library-api}
PEN evaluates the optimal \emph{mathematical library} to build within a type theory, not merely extensions of the core judgment forms.
For steps~1--9, the library grows by adding foundational primitives (new type formers, inductive types, fibrations); each primitive directly enlarges the set of inference rules.
For steps~10--14, the library grows by adding \emph{API specifications}: packages of types, operations, and axioms (e.g., a metric $g : TX \otimes_s TX \to \R$ with its associated Levi-Civita connection, geodesics, and curvature operators) that are not derivable from prior structures but are stated within the existing logic.
The candidate's $\kappa$ is the Kolmogorov complexity (clause count) of the API specification; $\nu$ is the combinatorial explosion of derivation schemas the specification natively unlocks.
A uniform before/after inhabitation diagnostic is also available (\S\ref{sec:uniform-nu}) and is informative on the foundational steps, but the manuscript's canonical totals are set by the strict hybrid evaluator described in \S\ref{sec:verification}.
\paragraph{Scalar provenance policy (strict first-use).}
Any API that mentions $\R$ must include an explicit construction route or explicit scalar axiom in the candidate specification; no ambient scalar base type is assumed.
Under the strict V1 policy used in this manuscript, scalar prerequisites are charged at first use: if a local API has cost $\kappa_{\mathrm{local}}$ and introduces scalar burden $\kappa_{\mathrm{scalar}}$, the scored denominator is $\kappa' = \kappa_{\mathrm{local}} + \kappa_{\mathrm{scalar}}$.
This enforces the empty-library premise at the clause-accounting level and removes any off-ledger scalar subsidy.
\end{remark}

\begin{observation}[DCT fixed-point preview]
\label{thm:spectral-preview}
The terminal candidate selected at step~15 is the temporal-cohesive shell whose semantic completion is read as DCT, with executable values $\nu=103$, $\kappa=8$, and $\rho=12.88$, clearing the bar by the largest absolute overshoot in the sequence.
Section~\ref{sec:decomposition} consolidates the evidence that this endpoint is the terminal structural fixed point selected by the objective rather than by target-conditioned theorem search.
\end{observation}

\begin{example}[Foundation steps]
\label{ex:foundation-steps}
The Generative Capacity definition captures novelty that is invisible to type-inhabitation methods:
\begin{itemize}[nosep]
    \item \emph{Witness} ($\nu = 2$): $\star : \mathbf{1}$ is an Introduction rule ($\nu_G = 1$); pattern matching on $\mathbf{1}$ is an Elimination rule ($\nu_C = 1$).
    \item \emph{$\Pi/\Sigma$ types} ($\nu = 5$): $\lambda$-abstraction and pair formation are Introduction rules ($\nu_G = 2$); application, fst, snd are Elimination rules ($\nu_C = 3$).
    \item \emph{Circle $S^1$} ($\nu = 7$): Five newly inhabited type schemas are Introduction rules ($\nu_G = 5$); \texttt{loop} adds two Computation rules ($\nu_H = 2$).
\end{itemize}
\end{example}

\begin{lemma}[Adjoint Completion Principle]
\label{lem:adjoint-completion}
In any hyperdoctrine/comprehension-category semantics of the library growth process, elimination structure is forced by the adjunctions interpreting substitution and quantification.
Granting elimination credit for introductions is therefore a completeness requirement of the semantics, not a heuristic bonus.
\newline
Specifically:
\begin{enumerate}[nosep]
    \item Reindexing/substitution sits in the standard adjoint triple $\Sigma_f \dashv f^* \dashv \Pi_f$, so introduction and elimination rules are paired by the left/right adjoints of the underlying comprehension structure.
    \item A type former with Introduction rule $\Gamma \vdash \mathrm{intro} : F(\Gamma)$ therefore determines an Elimination rule $\Gamma, x : F(\Gamma) \vdash \mathrm{elim}(x) : C$ satisfying the $\beta$- and $\eta$-laws required by the adjunction.
    \item No independent heuristic enumeration of these eliminators is needed: once the adjoint data is fixed, the corresponding elimination capability is semantically mandatory.
\end{enumerate}
This is the categorical content of Lawvere's adjointness thesis~\cite{lawvere-adjoint}: the unit/counit of the adjunction determine the $\beta$- and $\eta$-laws via the triangle identities.
Because elimination rules are term-level operations rather than new inhabited types, they are invisible to pure inhabitation counts.
The Adjoint Completion Principle closes that observational gap categorically, so counting the induced eliminators is required by the semantics of the logical framework itself.
The uniform diagnostic (\S\ref{sec:uniform-nu}) matches these two bootstrap steps exactly.
\end{lemma}

\begin{theorem}[Telescopic elimination capacity for API bundles]
\label{thm:telescopic-elimination}
Let $\mathcal{S}$ be a Minimal Complete API in the sense of \cref{def:minimal-api}, and let $\mathrm{Refs}(\mathcal{S})=\{L_1,\dots,L_k\}$ be the set of distinct prior library interfaces that occur as direct imports in its explicit reference DAG.
Assume that these imports telescope over a dependency DAG with a unique maximal imported interface under reachability, denoted $L_{\max}$.
Then the elimination and cross-interaction capacity natively unlocked by $\mathcal{S}$ satisfies
\[
  \nu_C(\mathcal{S})
  =
  \nu(L_{\max}) + \kappa(\mathcal{S}) + (k-1).
\]
This is an exact proof-theoretic decomposition of the API bundle's native eliminative surface area, not an evaluator-side heuristic.
\end{theorem}

\begin{proof}
The three terms come from three distinct structural sources.

\emph{Functorial subsumption:} the imported interfaces form a dependency DAG rather than a disjoint sum.
Summing the $\nu$-values of all imports would therefore double-count shared ancestry, violating weakening and substitution along the import telescope.
Because $\mathcal{S}$ is evaluated over a dependent context, it acts as a fibrant extension of the deepest imported derivation tree.
Every derivation schema available in the unique maximal import $L_{\max}$ canonically lifts along substitution to the parameterized bundle defined by $\mathcal{S}$.
Hence the inherited base contribution is exactly $\nu(L_{\max})$, not $\sum_i \nu(L_i)$.

\emph{Internal adjoint duality:} by \cref{def:minimal-api}, the $\kappa(\mathcal{S})$ clauses form a single structurally unified record-like telescope, not a loose collection of axioms.
Each constructively irreducible field introduces one new internal surface that must remain computationally accessible inside the bundle.
By the Adjoint Completion Principle (\cref{lem:adjoint-completion}), each such field forces exactly one canonical internal projection, destructor, or computation-facing accessor.
Therefore the bundle contributes exactly $\kappa(\mathcal{S})$ native internal elimination schemas.

\emph{Context amalgamation:} the direct imports $\mathrm{Refs}(\mathcal{S})$ are distinct namespaces before assembly.
To fuse $k$ such branches into one structurally unified telescope, the dependency graph must acquire the minimal number of bridge edges needed to make the imported components interact as a single strongly connected bundle.
Graph-theoretically, this is a spanning tree on $k$ vertices and therefore contributes exactly $k-1$ amalgamation edges.
Proof-theoretically, each such edge manifests as one required compatibility schema, exchange law, bridge equation, or naturality square connecting two previously separate imports inside the new bundle.

Adding the inherited base, the $\kappa(\mathcal{S})$ internal eliminators, and the $k-1$ bridge schemas yields
\[
  \nu_C(\mathcal{S})
  =
  \nu(L_{\max}) + \kappa(\mathcal{S}) + (k-1).
\qedhere
\]
\end{proof}

\begin{example}[Late-stage API bundles]
\label{ex:telescopic-elimination}
The theorem is the exact accounting rule used by the manuscript's explicit-reference late API bundles.
\begin{itemize}[nosep]
    \item \emph{Metric (Step~13):} the bundle has $\kappa=7$ and directly imports Curvature and Connections, so $k=2$ and $L_{\max}=\mathrm{Curvature}$ with $\nu(L_{\max})=34$.  Hence
    \[
      \nu_C = 34 + 7 + (2-1) = 42.
    \]
    The extra bridge term is the unique cross-interface synchronization linking the metric shell to both imported layers inside one structurally unified telescope.
    \item \emph{Hilbert (Step~14):} the bundle has $\kappa=9$ and directly imports Metric, Curvature, and Connections, so $k=3$ and $L_{\max}=\mathrm{Metric}$ with $\nu(L_{\max})=46$.  Hence
    \[
      \nu_C = 46 + 9 + (3-1) = 57.
    \]
    Here the two bridge terms account for the minimal amalgamation needed to bind the analytic shell to the geometric stack as one package.
\end{itemize}
\end{example}

\begin{definition}[Efficiency]
\label{def:efficiency}
$\rho(X) := \nu(X) / \kappa(X)$.
\end{definition}

\subsection{Selection Dynamics}

\begin{definition}[Structural Inflation]
\label{def:inflation}
$\Phi_n := \Delta_n / \Delta_{n-1} = F_n/F_{n-1} \to \varphi$.
\end{definition}

\begin{remark}
\label{rem:phi-not-tau}
Inflation is defined from marginal integration cost, not cumulative elapsed time: if one used $\tau_n / \tau_{n-1}$, then at $n = 4$ the value $7/4 = 1.75$ would raise the bar above $\rho_4 = 1.67$ and block the infrastructure phase. The marginal ratio $\Phi_4 = 3/2 = 1.50$ preserves admissibility at Step~4. Inflation therefore measures the \emph{marginal} growth of interface debt, not the cumulative burden.
\end{remark}

\begin{definition}[Cumulative Baseline]
\label{def:omega}
$\Omega_{n-1} := \sum_{i=1}^{n-1} \nu_i / \sum_{i=1}^{n-1} \kappa_i$.
\end{definition}

\subsection{The Five Axioms}

\begin{axiom}[Cumulative Growth]
\label{ax:cumulative}
$R(\tau - 1) \subseteq R(\tau)$.
The library only grows; realized structures are never removed.
\end{axiom}

\begin{axiom}[Horizon Policy]
\label{ax:horizon}
After each realization: $H \leftarrow 2$.
After each idle tick: $H \leftarrow H + 1$.
\end{axiom}

\begin{axiom}[Admissibility]
\label{ax:admissibility}
Candidate $X$ is admissible iff derivable from $\mathcal{B}$ and $\kappa(X) \leq H$.
\end{axiom}

\begin{remark}[Admissibility constraints for axiomatic extensions]
\label{rem:admissibility-constraints}
For steps in the axiomatic regime (Steps~10--14), where candidates are type-theoretic axioms rather than explicit constructions, admissibility imposes four structural constraints:
\begin{enumerate}[nosep]
    \item \emph{Structural unity:} Each candidate introduces a single structural bundle in the sense of \cref{def:api-graph,def:minimal-api}.  (For example, Cohesion's four modalities $\flat \dashv \Disc \dashv \sharp$ form an indivisible adjoint string; proposing $\flat$ alone would violate the adjunction.)
    \item \emph{Well-formedness:} All type signatures in the specification must be well-typed in the current library~$\mathcal{B}$.
    \item \emph{Non-redundancy:} The extension must not be derivable from~$\mathcal{B}$.
    \item \emph{Minimality:} $\kappa$ counts the clauses of the selected specification (\cref{def:effort}); the selection mechanism jointly evaluates $\nu$ and $\kappa$.
\end{enumerate}
Constraints~(1)--(3) prevent ``axiom packing''---the artificial inflation of $\nu$ by bundling unrelated axioms or declaring vacuous content.  All 15 steps of the PEN Synthesis Trajectory satisfy these constraints: each is a structurally unified extension whose components are interdependent (see \cref{rem:axiom-packing}).
\end{remark}

\begin{remark}[Defense against axiom packing]
\label{rem:axiom-packing}
A natural objection: what prevents an opaque axiom $\mathtt{StandardModel} : \mathcal{U}$ with $\kappa = 1$ from claiming unbounded~$\nu$?  The answer is that Generative Capacity measures \emph{atomic inference rules added to the derivation logic}, which requires structured interaction with the library:

\begin{itemize}[nosep]
    \item \emph{Opaque constants generate negligible~$\nu$.}  An axiom $X : \mathcal{U}$ with no structured type signature has $\nu_G \leq 1$ (one formation rule) and $\nu_C = 0$ (no elimination principle), giving total $\nu \leq 1$ and $\rho \leq 1$.  This fails the selection bar after Step~2.
    \item \emph{Structured axioms generate moderate~$\nu$.}  Cohesion's four modalities have signatures referencing~$\mathcal{U}$, yielding $\nu_G = 2$ (unit maps) and $\nu_C = 17$ (counits, structural rules, cross-interactions with existing library types).  The amplification requires the axiom to \emph{reference existing library structure} in its type signature.
    \item \emph{Bounded amplification.}  By Combinatorial Schema Synthesis (\cref{thm:tensor}), an axiom whose specification references~$k$ library types generates at most $O(k^d)$ new schemas at depth~$d$.  This bounds the $\nu/\kappa$ ratio by the library's combinatorial capacity.
\end{itemize}

\noindent The defense is structural: Generative Capacity rewards \emph{inferential interaction with the library}, not mere existence in~$\mathcal{U}$.  Axiom packing increases $\kappa$ without proportionally increasing $\nu$, because opaque content provides no elimination principles and composes with nothing.  Minimal overshoot then acts on the surviving candidates as a constructive-primality filter rather than as a reward bonus.
\end{remark}

\begin{definition}[Constructive primality]
\label{def:primality}
Let $\mathcal{S}$ be a bar-clearing candidate specification with API dependency graph $G_{\mathcal{S}}$.
We call $\mathcal{S}$ \emph{constructively composite} at step~$n$ if there exists a proper induced subgraph $C \subsetneq G_{\mathcal{S}}$ that is a union of terminal strongly connected components such that the induced sub-specification $\mathcal{S}|_C$ is well-typed over~$\mathcal{B}_n$, constructively irreducible, and already satisfies $\rho(\mathcal{S}|_C) \ge \mathrm{Bar}_n$.
Otherwise $\mathcal{S}$ is \emph{constructively prime}.
\end{definition}

\begin{axiom}[Selection via Variational Optimization]
\label{ax:selection}
The Selection Bar is $\mathrm{Bar}_n := \Phi_n \cdot \Omega_{n-1}$. Rather than a greedy, unconstrained maximization of novelty, PEN operates via the minimization of a structural objective functional over the moduli space of candidate MBTT telescopes. For any candidate $X$, we define the cost functional:
\begin{equation}
    \mathcal{J}_n(X) = \mathcal{P}_{\mathrm{adm}}(X) + \mathcal{P}_{\mathrm{bar}}(X) + \gamma \big( \rho(X) - \mathrm{Bar}_n \big) + \epsilon \kappa(X)
\end{equation}
Where:
\begin{itemize}[nosep]
    \item $\mathcal{P}_{\mathrm{adm}}(X)$ applies an infinite penalty to ill-typed or non-canonical telescopes (enforcing absolute constraint satisfaction).
    \item $\mathcal{P}_{\mathrm{bar}}(X)$ applies an infinite penalty if $\rho(X) < \mathrm{Bar}_n$ (the sub-criticality barrier).
    \item $\gamma > 0$ weights the positive efficiency surplus, and $0 < \epsilon \ll \gamma$ acts as a Kolmogorov tie-breaker favoring minimal syntax.
\end{itemize}
The selected candidate $X_{n+1}$ is the global minimum $\arg\min_X \mathcal{J}_n(X)$. This uniquely and analytically enforces the \emph{minimal positive overshoot} criterion.
\end{axiom}

\begin{remark}[Minimal overshoot as least structural action]
Formulating selection as the minimization of $\mathcal{J}_n(X)$ replaces heuristic sorting with formal constrained optimization. If the engine greedily maximized $\rho$, it would be vulnerable to ``axiom packing''—absorbing multiple independent components into a single composite update. By selecting the global minimum of the objective functional, the system dynamically tracks the frontier of the optimization landscape. It settles into the nearest admissible structural ground state, absorbing exactly the required capacity to clear the historical debt while leaving no unmotivated structural excitations.
\end{remark}

\begin{remark}[Algorithmic API assembly and compositional synergy]
\label{rem:api-assembly}
The structural unity of steps~10--14 is not treated as a semantic preference in scoring; it is an algorithmic consequence of the bar dynamics.
Isolated low-clause additions typically have weak cross-interaction ($\nu_C$) and fail the rising threshold, whereas interdependent typed clauses can increase novelty superlinearly relative to additive effort.
In practice, the selected API-like bundles are the candidates whose internal cross-references and library interactions produce enough compositional amplification in $\nu$ to remain above the bar.
A concrete Step~13 stress test makes this explicit: a bare metric tensor alone ($\kappa=1$) yields only weak novelty ($\nu_G=1$, $\nu_C\leq 2$), so $\rho\leq 3.0$, far below the Step~13 bar $5.99$.
By \cref{lem:metric-scc}, only the coupled Levi-Civita/Hodge/Laplacian/Ricci package forms a surviving structural bundle, so bundling is selected by proof-theoretic survival pressure rather than curator discretion.
\end{remark}

\begin{axiom}[Coherent Integration]
\label{ax:integration}
When $X_{n+1}$ is realized, it produces an integration layer $L_{n+1}$ with gap $\Delta_{n+1} := \kappa(L_{n+1})$.
\end{axiom}

\subsection{Realization Time}

\begin{definition}
\label{def:tau}
$\tau_n := \sum_{i=1}^{n} \Delta_i = F_{n+2} - 1$ for $d = 2$.
\end{definition}

\subsection{Executable PEN Loop}
\label{sec:pen-loop}

The implementation follows the same abstract loop at every step:
\begin{enumerate}[nosep]
    \item Compute strict admissibility from the active two-layer interface debt.
    \item Enumerate exact-band anonymous MBTT telescopes in the currently admissible structural lane.
    \item Type-check candidates, deduplicate by canonical key, and evaluate each candidate at exact structural cost.
    \item Apply semantic minimality by testing proper terminal SCC sub-bundles against the same active bar.
    \item Select by minimal positive overshoot (tie-break by lower $\kappa$, then deterministic key).
    \item Persist the winner together with bounded frontier evidence, then advance the discovered library.
\end{enumerate}

In pseudocode:
\begin{verbatim}
for step n:
  A <- StrictAdmissibility(B_n)
  for exact band k in prioritized bands:
    C <- EnumerateAtomicTelescopes(B_n, k)
    C <- TypeCheckAndDedup(C, B_n)
    E <- StrictEvaluate(C, B_n)  # exact kappa, hybrid nu
    V <- SemanticMinimalityFilter({x in E | rho(x) >= Bar_n})
    RecordRetainedEvidence(E, V)
    if V != empty: break
  if no viable V: halt
  x* <- argmin_{x in V} (rho(x)-Bar_n, kappa(x), stable_key(x))
  B_{n+1} <- InsertWinner(B_n, x*)
\end{verbatim}

This is the canonical path used in the current bounded live atomic lane: no paper-value substitution, no step-indexed reference injection in scoring, and no target-conditioned semantic reward.
The key tractability moves are structural rather than semantic: exact admissibility, canonical dedupe, SCC minimality, bounded frontier retention, and deterministic artifact-backed resume.

\subsection{Exact-Band Atomic Search}
\label{sec:mdp-mcts}

The executable search is best understood as bounded exact-band atomic synthesis over anonymous MBTT telescopes rather than as blind AST fuzzing.
Each atomic candidate now has three relevant representations:
\begin{itemize}[nosep]
    \item \textbf{Charged telescope.} The anonymous MBTT object charged by the selector's exact $\kappa$.
    \item \textbf{Canonical structural key.} The deterministic quotient used for deduplication and tie-breaking.
    \item \textbf{Persisted evidence surface.} Candidate-level traces, retained-valid summaries, frontier retention, and resume artifacts written downstream of selection.
\end{itemize}

The current live lane is bounded by exact admissibility bands, active-window connectivity constraints, deterministic retention policy, and the presently supported live range through step~15.
This does \emph{not} mean the selector is choosing from a hand-authored late-stage menu.
It means the current executable claim is about objective filtering over a bounded atomic region of the anonymous MBTT space under one fixed admissibility and selection rule.
Older molecular and Haskell artifacts remain donor provenance only.

The manuscript canon is fixed by bounded live atomic search together with hybrid strict evaluation:
\[
  \nu_{\mathrm{strict}} = (\nu_G+\nu_C)_{\mathrm{native}} + \nu_{H,\mathrm{explicit}},
\]
where the first term is computed on the evaluated telescope and the second counts only explicit topological witness payload whose framework-derived computations are present in the surviving witness basis.

% ============================================
% SECTION 5: COHERENCE WINDOWS
% ============================================
\section{Coherence Windows}
\label{sec:coherence}

The magnitude of the Integration Latency is determined by the foundation's \emph{Coherence Window}---the depth of historical context required to stabilize structural obligations.

\subsection{Induced Obligations}

\begin{definition}[Induced Obligations]
\label{def:obligations}
Let $\mathcal{O}^{(k)}(X)$ denote the set of normalized atomic obligations induced when candidate $X$ is sealed against a history of depth $k$.
Because the interface is cumulative:
$\mathcal{O}^{(1)}(X) \subseteq \mathcal{O}^{(2)}(X) \subseteq \mathcal{O}^{(3)}(X) \subseteq \cdots$
\end{definition}

\begin{definition}[Coherence Window]
\label{def:window}
A foundation has Coherence Window $d$ if for all candidates $X$ and all $k \geq d$:
$\mathcal{O}^{(k)}(X) \cong \mathcal{O}^{(d)}(X)$.
\end{definition}

\begin{definition}[Obligation Reduction]
\label{def:reduction}
An obligation $o \in \mathcal{O}^{(k)}(X)$ referencing layers $L_n, \ldots, L_{n-j}$ \emph{reduces at depth~$j$} if it decomposes into obligations referencing only $L_n, \ldots, L_{n-j+1}$.
An obligation is \emph{irreducible at depth~$j$} if it references layer $L_{n-j+1}$ and does not reduce.
\end{definition}

\begin{definition}[Historical support and arity]
\label{def:arity}
For a normalized obligation $o$, let $\mathrm{Supp}(o)$ be the set of historical layers whose exported interface clauses occur irreducibly in the typing derivation of~$o$ after each previously sealed layer is collapsed to its opaque interface.
The \emph{historical arity} of $o$ is $a(o) := |\mathrm{Supp}(o)|$.
A direct lookup to a deep library entry counts as arity~1: it is a substitution morphism into one exported interface, not an irreducible cross-layer coherence.
\end{definition}

\begin{lemma}[Arity-to-dimension correspondence]
\label{lem:arity-dimension}
An irreducible obligation of historical arity~$k$ requires a $k$-dimensional coherence cell to state.
Unary interactions are 1-cells (substitution/transport), binary interactions are 2-cells (naturality squares/homotopies), and ternary or higher interactions require 3-cells or higher.
\end{lemma}

\begin{proof}
The normalized typing derivation of an obligation records how the new candidate must cohere with each element of its historical support.
With one historical input, the required witness is a morphism or transport map, hence a 1-cell.
With two inputs, the witness compares two unary composites and therefore takes the form of a square/homotopy, hence a 2-cell.
More generally, relating $k$ independently supplied historical interfaces requires a $k$-dimensional filler; the Mac~Lane pentagon is the prototype at $k = 3$.
\end{proof}

\begin{definition}[Coherence Obligation by Dimension]
\label{def:dim-obligation}
When candidate $X$ with cell presentation $P = (C_0, C_1, C_2, \ldots)$ is sealed against $\mathcal{B}_n$, the elimination data for a type family $Y : X \to \U$ decomposes:
\begin{itemize}[nosep]
    \item \emph{Dimension~0:} For each $c \in C_0$, a point $d_c : Y(c)$.
    \item \emph{Dimension~1:} For each $p \in C_1$ with $p : a = b$, a path $d_p : \mathrm{transport}^Y(p, d_a) = d_b$.
    \item \emph{Dimension~$k$:} For each $k$-cell $s \in C_k$, a $k$-path witnessing coherence of $(k{-}1)$-dimensional data.
\end{itemize}
\end{definition}

\subsection{\texorpdfstring{Theorem A: Extensional Systems ($d = 1$)}{Theorem A: Extensional Systems (d = 1)}}

\begin{theorem}[Extensional Coherence Window]
\label{thm:ext-window}
In MLTT + UIP (or any type theory where identity types are h-propositions), the Coherence Window is $d = 1$.
\end{theorem}

\begin{proof}
\textbf{Upper bound ($d \leq 1$).}
In MLTT+UIP, all types are h-sets: for any $a, b : A$, the identity type $a =_A b$ is either empty or contractible.

Dimension-0 obligations (point data $d_c : Y(c)$) reference only $L_n$.
Dimension-1 obligations require paths $d_p : \mathrm{transport}^Y(p, d_a) = d_b$; since $Y(b)$ is an h-set, this type is a proposition, so $d_p$ carries no independent data.
Dimension-$k$ obligations for $k \geq 2$ involve coherence between paths in an h-set, which is trivially contractible.
Therefore $\mathcal{O}^{(k)}(X) \cong \mathcal{O}^{(1)}(X)$ for all $k \geq 1$.

\medskip\noindent\textbf{Lower bound ($d \geq 1$).}
$d = 0$ would mean no obligations at all, but sealing any structure requires checking well-typedness against $L_n$.
\end{proof}

\subsection{\texorpdfstring{Theorem B: Intensional Systems ($d = 2$)}{Theorem B: Intensional Systems (d = 2)}}

\begin{theorem}[Intensional Coherence Window]
\label{thm:int-window}
In HoTT as operationalized by the fully coupled CCHM cubical setting used here, the Coherence Window is $d = 2$.
\end{theorem}

The proof splits into an upper bound and a lower bound.

\subsubsection{\texorpdfstring{Upper Bound: $d \leq 2$}{Upper Bound: d <= 2}}

\begin{theorem}[Coherence Upper Bound]
\label{thm:upper}
For any cell presentation $P$ introduced at step $n+1$, every irreducible coherence obligation has historical arity at most~2.
Equivalently, after sealing encapsulation no independent obligation requires simultaneous access to more than layers $L_n$ and $L_{n-1}$.
\end{theorem}

\begin{proof}
\textbf{Stage~1: Unary lookups do not create multi-layer debt.}
Any obligation whose normalized derivation mentions only one prior export is arity~1 in the sense of \cref{def:arity}.
This includes direct variable lookups to arbitrarily old layers: after sealing encapsulation they are substitution morphisms into one opaque interface, hence 1-cells by \cref{lem:arity-dimension}.
They may mention a deep library pointer, but they do not constitute irreducible multi-layer coherence.

\medskip\noindent\textbf{Stage~2: Genuine sealing pressure is binary.}
New independent obligations arise when the candidate must cohere with a compatibility witness relating two prior exports.
These are naturality squares or homotopies: one historical input supplies the current operator interface, the other supplies the previously exported compatibility data.
Hence irreducible cross-layer sealing obligations have arity~2 and live on the 2-skeleton.

\medskip\noindent\textbf{Stage~3: Arity $\geq 3$ carries no independent data.}
Suppose an irreducible obligation had arity $k \geq 3$.
By \cref{lem:arity-dimension}, stating it would require an independent $k$-cell.
But the weak $\omega$-groupoid semantics of identity types (Lurie~\cite{lurie}, Prop.~1.2.5.1; Lumsdaine~\cite{lumsdaine}; van den Berg--Garner~\cite{vdberg-garner}) and the cubical Kan primitives imply that higher fillers are determined by their 2-dimensional boundaries.
Operationally, CCHM computes those fillers via \texttt{hcomp}/\texttt{comp}; they add no new independent sealing data.
Therefore every apparent arity-$k$ obligation with $k \geq 3$ either reduces to binary boundary data or is already packaged inside a previously sealed interface by \cref{rem:encapsulation}.
\end{proof}

\begin{remark}[Sealing Encapsulation]
\label{rem:encapsulation}
Stage~3 above uses a general principle: when layer $L_k$ is sealed, its resolved obligations become \emph{opaque exports}.
Future types interact with $L_k$'s interface---the theorems it proved---not with the underlying layers from which those theorems were derived.
In the notation of Stage~3, the coherence data for $(q, r)$ is an $L_{n-1}$ theorem, not raw $L_{n-2}$ data.
This encapsulation is essential for the Fibonacci recurrence: without it, inherited obligations would reference earlier layers directly, breaking the two-term structure.
The principle is verified by machine-checked Cubical Agda: at step~8, Group~B obligations are discharged from an opaque $L_7$~record with no PropTrunc import (\texttt{Saturation/AbstractionBarrier.agda}); at step~9, inherited obligations for the Hopf fibration are discharged from an opaque $L_8$~record (\texttt{Saturation/AbstractionBarrier9.agda}).
\end{remark}

\begin{remark}[Arity-to-dimension correspondence]
\label{rem:dim-temporal}
\begin{center}
\small
\begin{tabular}{lll}
\toprule
\textbf{Historical arity} & \textbf{Content} & \textbf{Effect on depth} \\
\midrule
1 & Lookup, substitution, transport & One prior export; can target any old layer \\
2 & Naturality square / homotopy & Binary interaction of two prior exports \\
$\geq 3$ & Higher coherence fillers & No independent data beyond 2-boundaries \\
\bottomrule
\end{tabular}
\end{center}
The point is categorical rather than chronological: a deep pointer is still arity~1, whereas an obligation linking $k$ distinct historical interfaces requires a $k$-cell to state.
In CCHM, independent data truncates at arity~2 because higher fillers are computed from the 2-skeleton.
\end{remark}

\subsubsection{\texorpdfstring{Lower Bound: $d \geq 2$}{Lower Bound: d >= 2}}

\begin{theorem}[Coherence Lower Bound]
\label{thm:lower}
There exist cell presentations whose coherence obligations irreducibly span two layers.
\end{theorem}

\begin{proof}
\textbf{Example (Hopf fibration).}
Let $L_{n-1}$ contain $S^1$ and $L_n$ contain $S^2$.
The Hopf fibration $h : S^3 \to S^2$ classifies a principal $S^1$-bundle over $S^2$.
Its clutching function $\gamma : S^1 \to \mathrm{Aut}(S^1)$ encodes how the fiber twists over the equator.
The coherence condition for the elimination principle involves compatibility of the clutching function with \emph{both} the $S^1$-action (from $L_{n-1}$) and the $S^2$-surface (from $L_n$).
This is a dimension-2 obligation that irreducibly references both layers.
\end{proof}

\begin{corollary}
By \cref{thm:upper} ($d \leq 2$) and \cref{thm:lower} ($d \geq 2$), the Coherence Window of the fully coupled CCHM/HoTT lane is exactly $d = 2$.
\end{corollary}

The preceding upper and lower bounds establish $d = 2$ for the fully coupled CCHM/HoTT lane from dimensional analysis of coherence cells.
We now derive the same result from three independent categorical and topological perspectives, elevating it from an empirical observation to a theorem of mathematical logic.

\subsection{The Adjunction Depth Barrier}
\label{sec:event-horizon}

In categorical logic, logical connectives and their inference rules are universally characterized by adjoint functors (e.g., Introduction $\dashv$ Elimination).

\begin{theorem}[The Adjunction Barrier]
\label{thm:adjunction-barrier}
A formal system with Coherence Window $d = 1$ cannot verify the triangle identities required for an adjunction. Therefore, $d = 2$ is the \emph{adjunction depth barrier}---the minimum historical depth required to support self-consistent mathematical operators.
\end{theorem}

\begin{proof}
Defining an adjunction $L \dashv R$ requires data spanning three categorical dimensions:
\begin{enumerate}[label=(\arabic*), nosep]
    \item \textbf{0-cells (Depth 0):} The functors $L$ and $R$.
    \item \textbf{1-cells (Depth 1):} The Unit $\eta : 1 \to R \circ L$ and Counit $\varepsilon : L \circ R \to 1$.
    \item \textbf{2-cells (Depth 2):} The Triangle Identities, $\varepsilon L \circ L\eta \simeq \mathrm{id}_L$ and $R\varepsilon \circ \eta R \simeq \mathrm{id}_R$.
\end{enumerate}

Suppose a candidate structure at step $n+1$ introduces an operator $L$ adjoint to a library operator $R \in L_n$. A system with $d = 1$ can establish the 1-cells ($\eta$ and $\varepsilon$) because it can interact with the immediately preceding layer. However, the triangle identities are 2-dimensional homotopies equating the composition of 1-cells to the identity.

Verifying these identities requires simultaneous visibility of the candidate's operators, the intermediate natural transformations in $L_n$, and the identity structure of the base categories in $L_{n-1}$. A $d = 1$ system is topologically blind to these 2-dimensional constraints spanning three layers (candidate $+ 2$ historical). It cannot verify the adjunction unless it enforces Uniqueness of Identity Proofs (UIP) to trivially collapse all 2-cells (\cref{thm:ext-window}). To natively compute adjoint coherence without truncating the topology, the fully coupled intensional lane analyzed here strictly requires $d = 2$.
\end{proof}

\subsection{Cohomological Upper Bound: Spectral Degeneration}
\label{sec:spectral-upper}

To translate this categorical requirement into a rigorous upper bound ($d \le 2$), we evaluate the homology of the candidate's coherence obligations.

\begin{definition}[Obligation Complex and Historical Filtration]
Let $\mathcal{O}_\bullet(X)$ be the cubical Kan complex of normalized coherence judgments $\Gamma \vdash t : T$ required to preserve global canonicity when sealing candidate $X_{n+1}$ against the library.
In the CCHM reading, $\mathcal{O}_\bullet(X)$ is the cellular Kan object generated by the candidate's boundary data, so its historical filtration may be compared to the usual skeletal/Postnikov filtration of a Kan complex.
We define a historical filtration indexed by the arity of irreducible historical support:
\begin{equation}
    F_0 \mathcal{O}_\bullet \subseteq F_1 \mathcal{O}_\bullet \subseteq F_2 \mathcal{O}_\bullet \subseteq \cdots \subseteq \mathcal{O}_\bullet
\end{equation}
where $F_p \mathcal{O}_\bullet$ is the subcomplex generated by obligations whose normalized contexts use only interface variables exported by the most recent $p$ layers $\Gamma_{n-p+1}, \ldots, \Gamma_n$ after each sealed layer is collapsed to its opaque interface.
The filtration is bounded: $F_2 = F_3 = \cdots$, as we now prove.
\end{definition}

\begin{theorem}[Spectral Degeneration at $E_2$]
\label{thm:spectral-degeneracy}
The homological spectral sequence associated with the historical filtration of the Obligation Complex degenerates at the $E_2$ page. Consequently, $F_2 \mathcal{O}_\bullet \simeq \mathcal{O}_\bullet$, establishing the precise topological meaning of $d \le 2$.
\end{theorem}

\begin{proof}
The filtration yields a spectral sequence $E^r_{p,q}$ converging to the total obligation homology. The $E_1$ page is given by the relative homology:
\begin{equation}
    E^1_{p,q} = H_{p+q}(F_p \mathcal{O}_\bullet / F_{p-1} \mathcal{O}_\bullet)
\end{equation}
An element of $F_p / F_{p-1}$ represents a coherence condition whose normalized typing judgment irreducibly depends on exactly $p$ historical layers.
By the arity-to-dimension correspondence (\cref{lem:arity-dimension,rem:dim-temporal}), such an obligation manifests as a $p$-dimensional coherence cell.

Now specialize to the cubical/CCHM operational semantics.
The homogeneous composition primitive \texttt{hcomp} computes fillers for open boxes of dimension $p \ge 3$ deterministically from the 2-dimensional boundary data already present in $F_2 \mathcal{O}_\bullet$.
Operationally, this means that the relative complex $F_p \mathcal{O}_\bullet / F_{p-1} \mathcal{O}_\bullet$ has no independent generators for $p \ge 3$: every apparent $p$-cell is the boundary of a canonical filler produced by \texttt{hcomp}.
Equivalently, these relative complexes are acyclic, so
\[
  E^1_{p,q} = H_{p+q}(F_p \mathcal{O}_\bullet / F_{p-1} \mathcal{O}_\bullet)=0
  \qquad (p \ge 3).
\]

Thus the $E_1$ page already vanishes in all columns $p \ge 3$.
The sequence advances to the $E_2$ page with differentials $d^r_{p,q} : E^r_{p,q} \to E^r_{p-r, q+r-1}$.
For $r \ge 2$, any differential touching a column $p \ge 3$ is forced to vanish because those columns are zero.
The only possible survivor is $d^2_{2,q} : E^2_{2,q} \to E^2_{0, q+1}$.
But $F_0 \mathcal{O}_\bullet$ consists only of isolated 0-cells (candidate constructors before coherence data), so its higher homology vanishes: $H_k(F_0)=0$ for $k \ge 1$.
Therefore $E^2_{0,q+1}=0$ for $q \ge 0$, and $d^2_{2,q}=0$ as well.

Thus all higher differentials vanish and the spectral sequence collapses at $E_2 \cong E_\infty$.
The historical filtration therefore stabilizes at the 2-skeleton: $F_2 \mathcal{O}_\bullet \simeq \mathcal{O}_\bullet$, so no obligation irreducibly spans more than two historical layers.
\end{proof}

\begin{remark}[Cubical Agda Mechanization]
This spectral degeneration precisely models the behavior of Cubical Agda. For the geometric test cases ($S^1$, $S^2$, $T^2$, Hopf), all 3-dimensional coherence obligations ($p \ge 3$) are automatically discharged by the \texttt{hcomp} (homogeneous composition) primitive. Because \texttt{hcomp} computes higher fillers from the 2-dimensional boundary data already present in $F_2$, the typechecker accepts the definition without requiring explicit imports from $L_{n-2}$. This is the computational shadow of the vanishing $E^1_{p,q}=0$ columns for $p \ge 3$.
\end{remark}

\subsection{Topological Lower Bound: The Clutching Family}
\label{sec:clutching-family}

To prove that $d = 2$ is not merely an upper bound but a strict structural requirement universally saturated by geometry, we exhibit an infinite family of topological structures where every member requires exactly depth-2, and none require depth-3.

\begin{theorem}[Fibrational Lower Bound]
\label{thm:clutching-bound}
The family of principal $G$-bundles over spheres $S^m$ with non-trivial clutching functions irreducibly requires exactly $d = 2$ coherence.
\end{theorem}

\begin{proof}
Let the base sphere $S^m$ reside in layer $L_n$, and the topological group $G$ (the fiber) reside in layer $L_{n-1}$. A principal $G$-bundle $P \to S^m$ is classified by its clutching function:
\begin{equation}
    \gamma : S^{m-1} \to \mathrm{Aut}(G)
\end{equation}
Sealing the total space $P$ at step $n+1$ requires the eliminator to verify that the transport of the fiber $G$ along the $m$-cell of $S^m$ correctly induces the specified automorphism $\gamma$.

This forms a dimension-2 coherence obligation. It irreducibly cross-references the path geometry of the base ($L_n$) with the internal algebraic symmetry of the fiber ($L_{n-1}$). A 1-layer window can access the equator or the fiber independently, but cannot cross-reference them to verify a 2-dimensional homotopy mapping one over the other. Thus, $d \ge 2$.

The Hopf fibration is the case $G = S^1$, $m = 2$, subsuming the concrete lower bound of \cref{thm:lower}.

Crucially, this family never requires $d = 3$. In \v{C}ech cohomology, the transition data for a bundle over a cover $\{U_i\}$ consists of 1-cocycles $g_{ij} : U_i \cap U_j \to G$ and 2-cocycles $g_{ij}g_{jk}g_{ki} = 1$ on triple intersections $U_i \cap U_j \cap U_k$.

Any $m$-sphere admits a minimal open cover of exactly two contractible charts (the northern and southern hemispheres, $U_1, U_2$). Their intersection is the equator ($U_1 \cap U_2 \simeq S^{m-1}$), carrying the 1-cocycle $\gamma$. Because there are only two charts, there are no triple intersections ($U_1 \cap U_2 \cap U_3 = \emptyset$). The \v{C}ech nerve of the minimal cover contains no 2-simplices.

Therefore, the 2-cocycle condition---which would demand a 3-dimensional coherence spanning three layers ($d=3$)---is structurally absent. The gluing occurs purely on a 2-layer interface. Every member of this infinite family requires exactly $d = 2$, perfectly locking the lower bound.
\end{proof}

% ============================================
% SECTION 6: COMPLEXITY SCALING
% ============================================
\section{The Complexity Scaling Theorem}
\label{sec:scaling}

\begin{theorem}[Complexity Scaling]
\label{thm:scaling}
For a foundation with Coherence Window $d$, evolving under PEN:
$\Delta_{n+1} = \sum_{j=0}^{d-1} \Delta_{n-j}$.
\end{theorem}

\begin{proof}
The interface is $I^{(d)}_n = \biguplus_{j=0}^{d-1} S(L_{n-j})$.
By disjointness and the Integration Trace Principle (\cref{lem:trace}), $|S(L_k)| = \Delta_k$, so
$\Delta_{n+1} = |I^{(d)}_n| = \sum_{j=0}^{d-1} \Delta_{n-j}$.
\end{proof}

\begin{corollary}
\label{cor:fibonacci}
For $d = 2$ with $\Delta_1 = \Delta_2 = 1$: $\Delta_n = F_n$ and $\tau_n = F_{n+2} - 1$.
\end{corollary}

\begin{remark}[Internal vs.\ external complexity]
\label{rem:int-ext-complexity}
The obligation count $\Delta$ measures \emph{external} complexity: the number of laws the interface imposes on future candidates.
The specification size $\kappa$ measures \emph{internal} complexity: the number of atomic specification clauses (type formations, introductions, equations) required to define the candidate over the library (\cref{def:effort}).
These can diverge sharply.
For example, the Riemann curvature tensor on an $N$-manifold has $N^4$ components, but the validation space for Connections (step~12) has $\Delta = 144$ basis elements---89 inherited from Connections' interface and 55 from Cohesion---because the interface is organized by the 144 independent algebraic symmetries, not by the raw component count.
The Fibonacci recurrence governs $\Delta$ (external), not $\kappa$ (internal).
\end{remark}

\subsection{Stagnation of Class~1 Systems}
\label{sec:stagnation}

For $d = 1$: $\Delta_{n+1} = \Delta_n$, so $\Delta_n = C$ (constant).
Inflation $\Phi_n = 1$; time $\tau_n = nC$ (linear).
The Cumulative Baseline $\Omega_{n-1}$ converges to a finite limit.
New candidates must clear a fixed threshold with diminishing novelty returns.

\subsection{Acceleration of Class~2 Systems}

For $d = 2$: $\Delta_n = F_n \sim \varphi^n$.
Inflation $\Phi_n \to \varphi$ from below: the dip at $\Phi_4 = 1.50$ provides breathing room for the infrastructure step ($\rho_4 = 1.67$).
Time $\tau_n \sim \varphi^{n+2}/\sqrt{5}$ (exponential).
The bar $\mathrm{Bar}_n = \Phi_n \cdot \Omega_{n-1}$ grows steadily, but the Combinatorial Novelty Theorem (\cref{sec:exponentiality}) ensures efficiency permanently outpaces it.

% ============================================
% SECTION 7: COMBINATORIAL NOVELTY
% ============================================
\section{The Combinatorial Novelty Theorem}
\label{sec:exponentiality}

The Scaling Theorem established $\Delta_n \sim \varphi^n$.
If novelty scaled only linearly with cost, efficiency would converge while the bar rises.
We prove novelty grows superlinearly.

\subsection{OIT Exponentiality}

\begin{theorem}[OIT Exponentiality]
\label{thm:oit-exponentiality}
An ordinary inductive type $X$ with $\Delta_0$ point constructors enables $2^{\Delta_0}$ distinct predicates $X \to \mathbf{2}$ at $O(\Delta_0)$ effort.
\end{theorem}

\begin{proof}
Each constructor independently maps to $\{\mathrm{true}, \mathrm{false}\}$; disjointness ensures semantic distinctness.
\end{proof}

\subsection{HIT Constraints}

\begin{remark}
\label{rem:hit-constraints}
For a HIT $X$, path constructors constrain eliminators into sets: $f : X \to \mathbf{2}$ must map path-connected points to the same value, giving $2^{|\pi_0(X)|}$ maps instead of $2^{\Delta_0}$.
Example: $S^1$ has $\Delta = 2$ but only $2^1 = 2$ maps to $\mathbf{2}$.
\end{remark}

\subsection{Restoring Superlinear Growth}

Three mechanisms compensate for HIT constraints.

\paragraph{Mechanism 1: Dependent Elimination.}
A type family $P : X \to \U$ chooses a fiber $P(c_i) \in \mathcal{B}$ for each point constructor and a transport equivalence for each path constructor.
This yields $|\mathcal{B}|^{\Delta_0} \cdot \prod_j |\mathrm{Aut}(P(s_j))|$ type families---far richer than maps to $\mathbf{2}$.

\paragraph{Mechanism 2: Library Cross-Interaction.}
Each prior type $T \in \mathcal{B}_n$ enables at least three new constructions ($X \to T$, $X \times T$, $\Sigma_{x:X} P(x)$), giving $\nu_n = \Omega(n)$.

\paragraph{Mechanism 3: Composite Constructions.}
Products and function types are superadditive: $\nu(X \times Y) \geq \nu_X + \nu_Y$; for OITs, multiplicative.

\subsection{Combinatorial Schema Synthesis}
\label{sec:ltp}

The most dramatic instance of superlinear novelty is the DCT's synthesis mechanism.

\begin{definition}[Syntactic LTL-Cohesive Extension]
\label{def:dct}
A \emph{temporal-cohesive Step-15 extension} is a type theory equipped with:
\begin{enumerate}[label=(\arabic*), nosep]
\item \textbf{Spatial Logic (Cohesion):} The adjoint string $(\flat \dashv \sharp,\; \Pi \dashv \Disc)$ from $R_{10}$.
\item \textbf{Temporal Core:} $\bigcirc$ (``next'') and $\Diamond$ (``eventually'') from LTL, together with the guarded core rules detected in the Step-15 AST.
\item \textbf{Spatial-Temporal Exchange:} Explicit commuting axioms between the temporal operators and the cohesive interface, exemplified by $\bigcirc(\flat X) \simeq \flat(\bigcirc X)$ and $\sharp(\Diamond X) \simeq \Diamond(\sharp X)$.
\item \textbf{Guarded Coherence:} The temporal self-interaction and elimination clauses audited in \cref{tab:step15}, which make the package executable as a typed AST rather than only semantically described.
\end{enumerate}
Construction effort: $\kappa = 8$ clauses, exactly as audited in \cref{tab:step15}.
No primitive infinitesimal object $\D$ appears in the selected AST; adding one would incur extra $\kappa$ with no syntactic justification from the engine's chosen specification.
In the remainder, ``DCT'' refers to the semantic completion of this temporal-cohesive shell once \cref{lem:infinitesimal-temporal} identifies its derived tangent behavior.
\end{definition}

\begin{theorem}[Combinatorial Schema Synthesis]
\label{thm:tensor}
When a synthesis candidate introduces $k$ new unary type formers into a library containing $|\mathcal{B}|$ structures, the number of non-trivial type-inhabitation schemas at depth~$d$ grows as $\Omega(|\mathcal{B}|^d \cdot k)$, yielding superlinear novelty.
\end{theorem}

\noindent \textbf{Application.}
The selected temporal-cohesive shell introduces temporal modalities ($\bigcirc$, $\Diamond$) into a library of 14~structures already equipped with spatial modalities ($\flat$, $\sharp$, $\Disc$, $\Pi_{\mathrm{coh}}$).
In the current bounded live atomic lane, this shell is evaluated at conceptual cost $\kappa=8$ and yields the canonical executable value $\nu = 103$, hence $\rho = 12.88$.
What matters for the present theorem is that the temporal shell clears the active bar by a wide margin and acts as a library-wide endofunctor over the already selected cohesive and differential layers.

\subsection{Divergence of Efficiency}

\begin{lemma}[Bounded Effort Growth]
\label{lem:log-effort}
The construction effort $\kappa_n$ grows at most linearly in the number of new type formers or modalities introduced at step~$n$.
Because candidates at late steps compose existing library abstractions rather than defining new primitives, the clause count is bounded by the number of compatibility axioms plus import declarations---a quantity independent of the library's cumulative size.
\end{lemma}

\begin{proof}
A candidate at step~$n$ is specified over a library of $n - 1$ entries.
Each clause either introduces a new type former (bounded by the candidate's intrinsic structure) or states a compatibility axiom with an existing library entry.
The number of compatibility axioms is bounded by the candidate's arity (how many existing modalities it must commute with), not by $n$ itself.
In the MBTT bit-length representation (\cref{sec:mbtt-audit}), each library pointer $\textsc{Lib}(i)$ costs $O(\log i)$ bits, giving logarithmic growth of the \emph{bit-length} per clause, but the clause count $\kappa$ depends only on the mathematical structure, not on library size.
\end{proof}

\begin{theorem}[Divergence of Efficiency]
\label{thm:divergence}
In a Class~2 foundation evolving under PEN, assuming an infinite supply of structurally distinct, continuous geometric primitives with nontrivial homotopy content, $\lim_{n \to \infty} \rho_n = \infty$.
\end{theorem}

\begin{proof}
By Combinatorial Schema Synthesis (\cref{thm:tensor}), the generative capacity $\nu_n$ grows at least polynomially as $\Omega(n^c)$ for some $c > 0$, due to library cross-interaction.
By \cref{lem:log-effort}, the construction effort $\kappa_n$ is bounded: each synthesis candidate composes existing library entries via a fixed number of compatibility axioms.
Therefore the efficiency $\rho_n = \nu_n / \kappa_n$ grows without bound as $\nu_n \to \infty$.

The bar $\mathrm{Bar}_n = \varphi \cdot \Omega_{n-1}$ is a cumulative average, growing at most linearly.
The ratio $\rho_n / \mathrm{Bar}_n$ is eventually increasing, so efficiency permanently clears the bar.
The temporal-cohesive endpoint is structurally guaranteed to be high-yield in this regime: combinatorial novelty dominates bounded specification cost.
\end{proof}

\begin{proposition}[Asymptotic rejection of discrete candidates]
\label{prop:discrete-asymptote-main}
Let $X$ be a candidate family with $\nu_H(X)=0$ and $\nu_G(X)=O(\kappa(X))$.
Then in the fully coupled $d=2$ PEN regime, $X$ eventually fails the selection bar permanently.
\end{proposition}

\begin{proof}[Proof sketch]
If $\nu_H(X)=0$, the candidate contributes no homotopy-driven superlinear term.
Under the stated hypothesis, novelty grows at most linearly in specification size, so $\rho(X)=\nu(X)/\kappa(X)$ remains bounded.
By \cref{thm:scaling,thm:divergence}, the active geometric frontier can sustain efficiencies that continue to outpace the Fibonacci-governed bar, while purely discrete families cannot.
Hence once the bar rises past the bounded discrete efficiency ceiling, all later members of the family remain subcritical.
\end{proof}

\subsection{Why Discrete Families Lose}
\label{sec:discrete-frontier}

\Cref{prop:discrete-asymptote-main} gives the asymptotic statement. The same mechanism is already visible at first admissibility, so the rejection of arithmetic and classical set-level packages is a main prediction of the objective rather than an appendix-only afterthought.

At Step~5 the active bar is already $\mathrm{Bar}_5 = 2.14$. Natural numbers enter with $\kappa=3$ and $\nu \le 4$, hence $\rho \le 1.33$; classical logic and power-set style axioms are even lower. By the time the library reaches the geometric API regime, candidates whose novelty remains essentially discrete are farther below the frontier. \Cref{tab:rejected} summarizes representative losing branches under the same fixed evaluator and shared vocabulary.

\begin{table}[tbp]
\centering
\caption{Rejected candidates at first admissibility under the same fixed evaluator and vocabulary. Rejection here means ``not an efficient frontier point under PEN,'' not ``mathematically invalid.''}
\label{tab:rejected}
\small
\begin{tabular}{@{}l rrrrl@{}}
\toprule
Candidate & $\nu$ & $\kappa$ & $\rho$ & Bar & Margin \\
\midrule
Natural numbers $\N$ & $\leq 4$ & 3 & $\leq 1.33$ & 2.14 (step 5) & $\leq -0.81$ \\
Classical logic (LEM) & $\leq 1$ & 1 & $\leq 1.00$ & 2.14 (step 5) & $\leq -1.14$ \\
Power set axiom & $\leq 2$ & 2 & $\leq 1.00$ & 2.14 (step 5) & $\leq -1.14$ \\
Ordinal arithmetic & $\leq 6$ & 5 & $\leq 1.20$ & 2.14 (step 5) & $\leq -0.94$ \\
Measure theory & $\approx 12$ & 6 & $\approx 2.0$ & 4.46 (step 10) & $\approx -2.5$ \\
Elementary topos axioms & $\approx 15$ & 8 & $\approx 1.9$ & 4.46 (step 10) & $\approx -2.6$ \\
Lie groups Lie$(S^3)$ & 9 & 6 & 1.50 & 4.46 (step 10) & $-2.96$ \\
Scheme theory & $\approx 18$ & 7 & $\approx 2.6$ & 5.42 (step 12) & $\approx -2.8$ \\
Galois theory & $\approx 10$ & 5 & $\approx 2.0$ & 5.42 (step 12) & $\approx -3.4$ \\
Symplectic geometry & $\approx 25$ & 5 & $\approx 5.0$ & 5.99 (step 13) & $\approx -1.0$ \\
\bottomrule
\end{tabular}
\end{table}

\paragraph{Structural reading.}
$\N$, classical logic, power-set style axioms, and ordinal arithmetic are all trapped by the same mechanism: they remain 0-groupoid or near-0-groupoid structures, so $\nu_H=0$ and novelty grows at most linearly in specification size. For $\N$ in particular, addition and multiplication do not create new $\kappa$-worthy novelty because they are definitionally derivable from $\mathsf{rec}_\N$.

\paragraph{Generous bundling still loses.}
Even an intentionally over-generous arithmetic package with metric-scale clause count $\kappa=7$ still lacks the homotopy term. Its best-case efficiency remains subcritical once the bar reaches the geometric regime, so geometry wins by topology-driven interaction growth rather than by curated bundling.

\paragraph{Higher but losing packages.}
Measure theory, elementary toposes, scheme theory, and Galois theory generate genuine structure but remain too discrete in the present scoring regime. Symplectic geometry comes closest to survival, but by Step~13 it still falls below the bar cleared by metric structure.

\begin{remark}[Divergence vs.\ termination]
\label{rem:divergence-vs-termination}
\Cref{thm:divergence} describes the \emph{theoretical trajectory} of efficiency growth: given an inexhaustible supply of structurally distinct geometric primitives, the combinatorial novelty of successive candidates grows without bound.
The algorithmic halting theorem (\cref{thm:end-of-history}) describes what \emph{actually happens} in PEN's current bounded live atomic lane: after Step~15, continuous internal continuations contribute zero marginal novelty up to equivalence, while disconnected extensions are starved by interface debt.
The two results are complementary, not contradictory: divergence guarantees that the selection mechanism \emph{never stalls prematurely} (efficiency always outpaces the bar while geometric primitives remain available), while halting identifies the first step at which PEN's internal/external squeeze leaves no admissible bar-clearer.
\end{remark}

% ============================================
% SUBSECTION: THE TANGENT TOPOS FIXED POINT
% ============================================
\subsection{The Tangent Topos Hypothesis and the Combinatorial Squeeze}
\label{sec:tangent-topos}

The abrupt termination of the PEN Synthesis Trajectory after Step~15 is not a heuristic failure of the search space, but an algorithmic phase transition inside PEN's current bounded live atomic lane. The selected temporal-cohesive shell internalizes continuous deformation inside the connected component of $\U_0$, while any disconnected extension must pay the full interface debt of the accumulated library. We formalize this by proving that the semantic DCT completion acts as a tangent-style internalization of the historical library and then showing that Step~16 faces a combinatorial squeeze.

\begin{theorem}[Syntactic-Semantic Transfer: Emergence of the Infinitesimal]
\label{lem:infinitesimal-temporal}
In the semantic completion of the Step-15 temporal-cohesive extension, the temporal ``next'' modality $\bigcirc$ is naturally equivalent to a synthetic tangent endofunctor. Equivalently, there exists a derived infinitesimal object $\D$ and natural equivalences
\[
  \theta_X : \bigcirc X \overset{\sim}{\longrightarrow} X^\D .
\]
The infinitesimal structure is therefore a semantic consequence of the selected temporal-cohesive AST, not an additional primitive clause in the Step-15 syntax.
\end{theorem}
\begin{proof}
The audited AST of \cref{def:dct,tab:step15} contributes three ingredients and nothing more: a guarded temporal endofunctor $\bigcirc$, the eventual modality $\Diamond$, and explicit exchange laws between the temporal operators and the cohesive interface.  In a cohesive semantics, these exchange laws force $\bigcirc$ to preserve the same discrete/shape boundary data preserved by the first-order tangent endofunctor.
\begin{itemize}[nosep]
    \item The exchange law with $\flat$ identifies $\bigcirc$ with a first-order neighborhood operator on discrete points: $\flat(\bigcirc X) \simeq \bigcirc(\flat X)$, so $\bigcirc$ adds no new macroscopic points and behaves like an infinitesimal thickening.
    \item The audited spatial-temporal interaction clause forces $\bigcirc$ to respect the pre-existing cohesive shape/discreteness decomposition exported by Cohesion, so $\bigcirc$ preserves first-order path data rather than introducing a disconnected second dynamics.
    \item The guarded self-composition and elimination clauses in \cref{tab:step15} provide exactly the first-order coherence needed to iterate $\bigcirc$ without introducing a second independent infinitesimal direction.
\end{itemize}
\noindent In the semantic completion used throughout the paper, such a first-order cohesive endofunctor is representable by an infinitesimal object $\D$, yielding a natural equivalence $\theta_X : \bigcirc X \simeq X^\D$.  This is the transfer from the executable temporal AST to the differential-geometric reading: the engine selected the temporal-cohesive syntax, and the infinitesimal object appears only as its exact semantic shadow.
\end{proof}

\begin{lemma}[Internalization of the Meta-Algorithm]
\label{lem:internalization}
The DCT internalizes its own structural evolution.  The metamathematical process of extending the foundational library becomes derivable as an object-level geometric flow.
\end{lemma}
\begin{proof}
By incorporating Nakano's temporal logic, the DCT inherits guarded recursion with L\"ob's rule
\[
  \mathsf{fix}_A : (\bigcirc A \to A) \to A .
\]
From this one derives the guarded coiterator for streams,
\[
  \mathsf{coiter}_A : (A \to \bigcirc A) \to A \to \mathrm{Str}(A),
\]
where $\mathrm{Str}(A)$ is the guarded stream type characterized by $\mathrm{Str}(A) \cong A \times \bigcirc \mathrm{Str}(A)$.

Now let $\theta_X : \bigcirc X \overset{\sim}{\to} X^\D$ be the temporal--infinitesimal equivalence from \cref{lem:infinitesimal-temporal}.
Given a vector field on the small univalent universe
\[
  v : \U_0 \to \U_0^\D,
\]
define the corresponding guarded step map by transport across~$\theta_{\U_0}$:
\[
  \mathsf{step}_v := \theta_{\U_0}^{-1} \circ v : \U_0 \to \bigcirc \U_0 .
\]
Typing is now explicit:
\[
  v : \U_0 \to \U_0^\D
  \xRightarrow{\;\theta_{\U_0}^{-1}\;}
  \mathsf{step}_v : \U_0 \to \bigcirc \U_0
  \xRightarrow{\;\mathsf{coiter}_{\U_0}\;}
  \mathsf{Flow}_v : \U_0 \to \mathrm{Str}(\U_0).
\]
Thus each vector field on~$\U_0$ unfolds to a guarded stream of successive deformations in the same universe.
Because the codomain is $\U_0$ throughout, the construction is confined to the connected component of~$\U_0$ and does not jump to a higher universe $\U_1$.
The metatheoretic act of ``proposing the next type'' is therefore represented internally as guarded geometric evolution on~$\U_0$.
\end{proof}

\begin{theorem}[Algorithmic Halting via the Combinatorial Squeeze]
\label{thm:end-of-history}
Let $S_{\mathrm{int}}$ be the admissible Step-16 candidates representable by guarded internal flow on the connected component of $\U_0$ determined by $\mathcal{B}_{15}$, and let $S_{\mathrm{ext}}$ be the remaining admissible Step-16 candidates.  Then every $Y \in S_{\mathrm{int}}$ has zero marginal external novelty, while every $Y \in S_{\mathrm{ext}}$ has $\rho(Y) < \mathrm{Bar}_{16}$.  Consequently no Step-16 candidate clears the bar, and the PEN Synthesis Trajectory halts at Step~15.
\end{theorem}
\begin{proof}
At Step~16 the selection bar is strictly positive:
\[
  \mathrm{Bar}_{16} = \Phi_{16} \cdot \Omega_{15} \approx 1.618 \times 5.44 \approx 8.80 .
\]
We partition admissible candidates into the two classes from the statement.

\emph{Internal/continuous side.}
Let $Y \in S_{\mathrm{int}}$.  By \cref{lem:internalization}, any such candidate is representable by a guarded internal flow
\[
  v : \U_0 \to \U_0^\D
\]
which transports to a guarded step map $\mathsf{step}_v : \U_0 \to \bigcirc \U_0$ and unfolds to a stream $\mathsf{Flow}_v : \U_0 \to \mathrm{Str}(\U_0)$.
By Univalence, paths in $\U_0$ are equivalences, so continuous deformation inside this connected component never leaves the existing equivalence class of the library's internal geometry.
Re-postulating such a $Y$ as an opaque external axiom therefore contributes no new marginal derivation rules:
\[
  \nu(Y \mid \mathcal{B}_{15}) = 0,
\qquad
  \rho(Y) = 0 < \mathrm{Bar}_{16}.
\]

\emph{External/discontinuous side.}
Now let $Y \in S_{\mathrm{ext}}$.
To avoid the Univalent Horizon, $Y$ must be a disconnected axiomatic jump.
Then maximal interface density is no longer optional: by the full-coupling regime of \cref{thm:canonicity-preservation,thm:univalent-functoriality} together with Fibonacci scaling (\cref{thm:scaling}), Step~16 demands the full interface burden
\[
  \Delta_{16} = F_{16} = 987 .
\]
Thus any admissible external jump must specify interaction clauses against the exported interface at this scale, forcing a massive specification burden.
At the same time, disconnected candidates do not inherit the temporal-cohesive path algebra that created the Step-15 amplification: they contribute no internal flow and no derived tangent lift, so their novelty is at best linear in the size of the external package.
Linear novelty cannot clear a bar driven by the Step-16 Fibonacci debt.
Equivalently, the opaque-axiom bounds from \cref{rem:axiom-packing} and the full-coupling requirement imply $\rho(Y) < \mathrm{Bar}_{16}$ for all $Y \in S_{\mathrm{ext}}$.

Hence Step~16 is squeezed from both sides: internal continuations are definitionally exhausted up to equivalence, and external jumps are computationally starved by interface debt.
The halting claim is therefore algorithmic and PEN-local, not a claim that mathematics globally terminates.
\end{proof}

% ============================================
% SECTION 8: DCT FIXED-POINT CONSOLIDATION
% ============================================
\section{The Temporal-Cohesive Shell as a Terminal Structural Fixed Point}
\label{sec:decomposition}
\label{sec:dct-discovery}

This section consolidates the manuscript's Step-15 claim in the order the paper intends it to be read: first the selected temporal-cohesive AST, then the executable quantities attached to that shell, and only then its semantic DCT completion.
Under fixed fully coupled settings, the PEN objective realizes the temporal-cohesive shell as the terminal high-yield extension, and the current bounded live atomic engine recovers that shell under bounded search.
The key claim is therefore structural rather than label-driven: a typed search-and-selection engine with no target theorem injection in its scoring path reaches a PEN-local fixed point at the Univalent Horizon of the current $d=2$ lane.

\subsection{Why the Temporal-Cohesive Shell Is a Terminal Signal}

\begin{observation}[Terminal temporal-cohesive fixed-point signal]
\label{thm:spectral}
In the current bounded live atomic mode, the engine selects a Step-15 temporal-cohesive shell with $\kappa=8$, $\nu=103$, and $\rho=12.88$, clearing the step-15 bar (7.40) by an absolute overshoot of $5.48$, the largest in the sequence.
The selection is made from admissible typed candidates under fixed rules and bounded budgets, without step-indexed target conditioning or reference-candidate replay in the scoring-selection lane.
\end{observation}

The interpretive claim is deliberately conservative.
We do \emph{not} claim global uniqueness across all possible logics or objectives.
What the theorems and the executable artifact jointly establish is narrower: in this PEN lane, the temporal-cohesive shell is the terminal selected structure, its accepted Step-15 total is the executable value $\nu=103$, and ``DCT'' names the semantic completion used to interpret that shell.
This keeps the proof boundary clean: selection and arithmetic are structural, while the DCT label is interpretive.

This is methodologically distinct from the dominant theorem-prover workflow.
Classical ITP/ATP systems validate user-supplied conjectures.
PEN instead ranks candidate structural extensions under a fixed typed objective, then halts at an internally explained boundary.
The formal contribution is therefore a metatheorem about typed library growth and fixed-point termination, with the semantic interpretation of the terminal AST left deliberately secondary.

\subsection{Consolidated DCT Property Bundle}

DCT properties proved or formalized across the manuscript and appendices can be read as one package:
\begin{itemize}[nosep]
    \item \textbf{Derived temporal-tangent equivalence} (\cref{lem:infinitesimal-temporal}): $\bigcirc X \simeq X^\D$ in the semantic completion of the Step-15 AST.
    \item \textbf{Internalization of meta-evolution} (\cref{lem:internalization}): extension dynamics become representable as internal geometric flow.
    \item \textbf{Combinatorial-squeeze halting} (\cref{thm:end-of-history}): beyond Step~15, internal continuations have zero marginal novelty and disconnected jumps fail the bar.
    \item \textbf{Internal tangent/temporal machinery} (\cref{thm:tangent,thm:temporal}): tangent lift and temporal evolution are native operators.
    \item \textbf{Symplectic-flow corollary} (\cref{thm:symplectic-flow}): flow/bracket structure is expressible internally once DCT is present.
\end{itemize}

Together, these results explain why DCT is qualitatively different from earlier steps: the selected library begins to express its own evolution mechanisms internally.

\subsection{Mechanized Step-15 Novelty Accounting}

The canonical Step-15 count is the machine value attached to the selected temporal-cohesive shell:
\[
  \nu_{\mathrm{exec}}(R_{15}) = 103.
\]
When the manuscript later writes $\nu(\mathrm{DCT}) = 103$, it is using ``DCT'' as semantic shorthand for this same Step-15 shell under its semantic completion.
Operationally, the current evaluator exposes the structural account
\[
  \nu(R_{15}) = \nu_G + \nu_H + \nu_C = 2 + 15 + 86 = 103
\]
while keeping the temporal shell at conceptual cost $\kappa = 8$.

This is the Step-15 value used throughout the manuscript.
The canonical manuscript value is the executable result that actually drives selection in the artifact.
Explanatory decompositions are retained because they now agree with that value and are exposed by the current code path.

\subsection{Calibration Facts Used by the DCT Analysis}

\begin{theorem}[Topological Projection of the Path Algebra (CCHM)]
\label{thm:topological-projection}
Let $X$ be a higher inductive type with $m$ path constructors, and let $d$ be the maximum path dimension.
In CCHM cubical type theory with connections, the computation-rule contribution satisfies
\[
  \nu_H = m + d^2.
\]
\end{theorem}

\noindent In the current strict evaluator this theorem is used as an exact computational shortcut on the candidate's explicit path data; the code does not dynamically enumerate every surviving reduction one by one when the same value is already fixed by the CCHM interaction matrix.

\begin{proof}[Proof sketch]
The $m$ term counts the explicit boundary reductions attached to the $m$ declared path constructors of~$X$.
Each constructor contributes one independent computation rule when its boundary is reduced against the corresponding eliminator.

For the cubical contribution, fix a maximal dimension-$d$ constructor.
In the CCHM path algebra, its independent computational interactions are generated by:
\begin{enumerate}[nosep]
    \item the $d$ diagonal self-reductions coming from composition along each coordinate, and
    \item the $d(d-1)$ ordered pairwise coordinate interactions coming from the connection operations $\lor,\land$ together with homogeneous composition \texttt{hcomp}.
\end{enumerate}
These assemble into a $d \times d$ Kan interaction matrix.
Because \cref{thm:spectral-degeneracy} collapses all genuinely higher-dimensional data to fillers of 2-dimensional boundaries, no further independent computation rules survive beyond this matrix.
Hence the Kan part contributes exactly
\[
  d + d(d-1) = d^2
\]
independent reductions, and therefore $\nu_H = m + d^2$.
\end{proof}

\begin{remark}[Sensitivity of $\nu_H$]
\label{rem:nuH-sensitivity}
The current bounded live atomic lane uses $m+d^2$ as the canonical calibration for the explicit path-style computation payload that survives into the current hybrid evaluator.
This is the operative formula for the HIT and higher-path parts of the strict canon.
The current ablation replay in \texttt{runs/review\_hardening/} is sensitive in exactly the way the theorem predicts: weakening the shortcut to $m+d$ halts the bounded live atomic lane after Step~6, while removing the dimensional term entirely halts after Step~4 before any geometric HIT survives.
Thus the quadratic term is not a cosmetic bonus on late geometry; it is the operative CCHM shortcut that keeps the geometric branch open at all.
\end{remark}

\begin{proposition}[Schema Canonicality]
\label{prop:grammar-canonical}
Depth-1 schema counting is window-independent after schematization of library atoms.
\end{proposition}

\begin{proof}[Proof sketch]
After collapsing library atoms to a schematic symbol $L$, depth-1 schema sets depend only on candidate constructor signatures.
This is verified computationally for steps~1--7 (\cref{sec:verification}).
\end{proof}

\begin{remark}[Witness and $\Pi/\Sigma$ attribution]
\label{rem:witness-pisigma-attribution}
At bootstrap steps, novelty naturally splits across constructor and elimination contributions:
Witness contributes one constructor plus one elimination action, while $\Pi/\Sigma$ contributes constructor forms plus apply/projection eliminations.
\end{remark}

\begin{proposition}[$S^2 \equiv S^3$ Syntactic Identity]
\label{prop:s2s3}
The depth-1 syntactic schema sets for $S^2$ and $S^3$ are identical, so their syntactic counts agree.
\end{proposition}

\begin{proof}
Both generate the same depth-1 schema shapes over $(L,X)$.
Their novelty gap is therefore not syntactic, but geometric/interactional.
\end{proof}
% ============================================
% SECTION 9: COMPUTATIONAL VERIFICATION
% ============================================
\section{Computational Verification}
\label{sec:verification}

\subsection{The Rust Atomic Engine}

The executable artifact used for the manuscript's canonical claim is the \texttt{pen-atomic} Rust workspace implementing PEN as a deterministic strict synthesis loop over anonymous MBTT telescopes.
Older Haskell code is donor provenance only and does not fix the manuscript's executable canon.
The canonical rerun for this manuscript is:
\begin{enumerate}[nosep]
    \item \texttt{cargo test --workspace}
    \item \texttt{cargo run -p pen-cli -- run --config configs/debug.toml --until-step 15}
    \item \texttt{cargo run -p pen-cli -- resume runs/<run-id> --until-step 15}
    \item \texttt{cargo run -p pen-cli -- export-agda --run-dir runs/<run-id>}
\end{enumerate}
The first command checks the full workspace; the second reproduces the live strict 15-step trajectory reported in \cref{tab:genesis}; the latter two exercise the frontier-backed resume and observational Agda export surfaces used by the current artifact.

\paragraph{Candidate generation.}
The strict engine computes admissibility from the active two-layer interface debt, opens a small set of exact $\kappa$-bands, and searches those bands in priority order.
Within each band it enumerates anonymous MBTT telescopes inside the currently admissible structural lane, then checks them for well-typedness, connectivity, and deterministic canonical distinctness.
The current supported live range is bounded through step~15, but within that range selection is performed by live atomic search rather than by replay or target-conditioned semantic ranking.

\paragraph{Charged cost and persisted artifacts.}
Each accepted candidate is an anonymous telescope with a charged exact structural cost $\kappa$.
Run artifacts record the selected telescope, its canonical hashes, candidate-level evidence, frontier retention summaries, step checkpoints, frontier manifests, and telemetry.
These artifacts make the current verification surface artifact-backed rather than replay-only.

\paragraph{Strict novelty computation.}
For live atomic candidates the strict engine uses a hybrid evaluator:
\[
  \nu_{\mathrm{strict}} = (\nu_G+\nu_C)_{\mathrm{native}} + \nu_{H,\mathrm{explicit}}.
\]
The native term is computed on the evaluated telescope.
The explicit topological term counts only path-style witness payload whose framework-derived computation clauses are physically present in the surviving witness basis.
In the current strict code path this payload is scored by the theorem-backed shortcut $\nu_H=m+d^2$ on the compiled witness basis, rather than by dynamically expanding all CCHM reductions individually.
This is the canonical source of the machine values in \cref{tab:genesis}.

\paragraph{Semantic minimality.}
Every bar-clearing candidate is subjected to semantic minimality filtering.
The engine removes terminal SCC sub-bundles, re-checks the remainder for well-typedness, and rejects any candidate whose proper sub-bundle already clears the active bar.
This operationalizes constructive primality in the current bounded live atomic lane.

\paragraph{Results.}
The bounded live atomic run recovers all 15 structures in the order shown in \cref{tab:genesis}, with machine totals
\[
  (1,1,2,5,7,8,10,18,17,19,26,34,46,62,103)
\]
at conceptual costs
\[
  (2,1,1,3,3,3,3,5,4,4,5,6,7,9,8).
\]
These strict machine values are the manuscript's canonical totals.
Deterministic reruns, frontier-backed resume, compatibility-forced step resume, compatibility-forced reevaluation, and observational Agda export all preserve the same accepted strict canon in stored run artifacts.

\paragraph{Rejected candidates.}
To verify that the trajectory is selective rather than merely permissive, the engine also evaluates natural but losing alternatives spanning arithmetic, classical logic, measure theory, scheme theory, elementary topos axioms, and symplectic geometry.
All of them fall below the active bar at first admissibility; the summary table appears in \cref{sec:discrete-frontier,tab:rejected}, with supplementary notes deferred to \cref{app:rejections}.

\subsection{Schema Canonicality Verification}

Legacy schema diagnostics remain useful as early-step consistency checks.
They are not the manuscript's canonical novelty measure or its current executable evidence surface.
For steps 1--7, the depth-1 schema set using all library atoms is identical to the proof-rank schema set using only the 2-step window.
This supports \cref{prop:grammar-canonical}: after schematization, depth-1 schema counting is effectively window-independent for the small foundational cases.

This diagnostic is not the manuscript's canonical novelty measure.
At depth~2 the schema count explodes combinatorially, so the exact-oracle mode is now treated only as an early-step consistency check rather than as a late-stage proof of $\nu$.

\subsection{Uniform Diagnostic}
\label{sec:uniform-nu}

A type-inhabitation comparison procedure remains useful as an auxiliary diagnostic on the foundational steps.
Its late-stage counts over-amplify API and temporal shells, so the manuscript does not use it to set any value in \cref{tab:genesis}.

The categorical point behind the diagnostic remains valid:
Adjoint Completion (\cref{lem:adjoint-completion}) explains why introduction data forces elimination data even when pure inhabitation counts cannot see it directly.
In this manuscript that principle is realized operationally by the current evaluator together with the strict hybrid accounting.

\subsection{Meta-Theoretic Rule Audit}
\label{sec:inference-nu}

The rule audit is now an explanatory reconstruction of the strict canon rather than an independent oracle.
For the foundational and geometric steps the reconstruction is exact.
For the late API shells it is constrained to agree with the machine total in \cref{tab:genesis}.
For DCT we record only the canonical total unless a finer split is justified directly by the current strict code path.

\paragraph{Results.}
\begin{center}
\small
\begin{tabular}{@{}cl rl@{}}
\toprule
$n$ & Selected AST skeleton & $\nu_{\mathrm{strict}}$ & Audit status \\
\midrule
1  & Universe         & 1  & exact foundational rule count \\
2  & Unit             & 1  & exact foundational rule count \\
3  & Witness          & 2  & exact foundational rule count \\
4  & $\Pi$/$\Sigma$   & 5  & exact foundational rule count \\
5  & $S^1$            & 7  & exact HIT reconstruction \\
6  & PropTrunc        & 8  & exact HIT reconstruction \\
7  & $S^2$            & 10 & exact HIT reconstruction \\
8  & $S^3$            & 18 & exact HIT/H-space reconstruction \\
9  & Hopf             & 17 & exact map reconstruction \\
10 & Modal shell      & 19 & exact API-shell reconstruction \\
11 & Connection shell & 26 & expository API-shell reconstruction \\
12 & Curvature shell  & 34 & expository API-shell reconstruction \\
13 & Endomorphic operator bundle & 46 & expository API-shell reconstruction \\
14 & Hilbert-functional shell & 62 & expository API-shell reconstruction \\
15 & Temporal-cohesive shell & 103 & executable total plus exposed structural split \\
\bottomrule
\end{tabular}
\end{center}

\begin{remark}[Status of the DCT decomposition]
For DCT the machine-verified total $\nu = 103$ is canonical, and the current evaluator also exposes the standard three-way split $\nu_G=2$, $\nu_H=15$, $\nu_C=86$.
This keeps the manuscript aligned with the executable artifact.
\end{remark}

\begin{remark}[Automation scope of the verification]
\label{rem:automation-scope}
The verification of the PEN Synthesis Trajectory decomposes into two logically independent claims: (a)~the \emph{selection ordering and bar clearance} ($\rho_n \ge \mathrm{Bar}_n$ for all $n$), and (b)~the \emph{rule-component accounting} ($\nu = \nu_G + \nu_C + \nu_H$).

Claim~(a) is verified computationally for \textbf{all 15~steps} by the current bounded live atomic run itself.
Given the admissible typed candidates produced by the generation layer, the PEN quantities $\nu$, $\kappa$, $\rho$, and bar clearance are computed automatically with no paper lookup and no target-conditioned evaluation term.

\noindent Claim~(b)---the internal rule-component accounting---now has mixed status.
For steps~1--10 the explanatory audit is exact.
For steps~11--14 it is an expository reconstruction constrained to match the strict total.
For step~15 the executable structural total $\nu = 103$ is canonical.

The implemented finding is that the current bounded live atomic lane recovers the 15-step PEN Synthesis Trajectory and clears every bar mechanically.
The stronger conjecture is that sufficiently powerful, materially less-pruned search over the same admissible space would recover the same trajectory without the present bounded search controls.
\end{remark}

\subsection{Determinism and Comparison Lanes}

The present manuscript supplements theorem-backed accounting with deterministic comparison lanes against the current atomic engine.
The stored artifact comparison surface now includes cold live runs, deterministic reruns, frontier-backed resume, compatibility-forced step resume, compatibility-forced reevaluation, spill-pressure runs, and reference-replay baselines, summarized by \texttt{scripts/compare\_runs.py}.

These comparisons sharpen two reviewer-facing points.
First, the canonical claim is about a live atomic run and stored artifacts, not replay-only reconstruction.
Second, frontier acceleration is subordinate to CPU-authoritative acceptance: resume and pressure handling may change provenance, but they do not change the accepted step-15 result.

\subsection{Cubical Agda Mechanization}

The Complexity Scaling Theorem is proved in Cubical Agda~\cite{cubical-agda}:
for $d = 2$, $\Delta_{n+1} = \Delta_n + \Delta_{n-1}$.
The proof covers initial conditions, the recurrence, the cumulative sum identity $\tau_n = F_{n+2} - 1$, and convergence $\Phi_n \to \varphi$.

Coherence Obligation Experiments trace the obligations for $S^1$, $S^2$, $T^2$, and the Hopf fibration in Cubical Agda.
In all cases: all obligations reference at most 2 layers, at least one genuinely references 2, and none references 3.

\paragraph{Obligation decomposition and the abstraction barrier.}
Explicit obligation enumeration at steps~7--9 verifies the Integration Trace Principle (\cref{lem:trace}):
$S^2$ resolves $13 = 8 + 5$ obligations (8~from PropTrunc, 5~from~$S^1$), $S^3$ resolves $21 = 13 + 8$ (13~from~$S^2$, 8~from~PropTrunc), and the Hopf fibration resolves $34 = 21 + 13$ (21~from~$S^3$, 13~from~$S^2$).
The Hopf fibration (step~9) crosses the phase boundary from types to maps: the 21~domain-side obligations arise from postcomposition constraints on~$S^3$, while the 13~codomain-side obligations arise from pullback constraints on~$S^2$, demonstrating a natural domain/codomain duality within the two-layer window.
A Cubical Agda module (\texttt{Saturation/AbstractionBarrier.agda}, \texttt{-{}-cubical -{}-safe}) defines $L_7$'s interface as an opaque record and discharges all Group~B obligations (8.6--8.10) from record fields alone, with no PropTrunc import.
At step~9, the same methodology extends: inherited obligations for the Hopf fibration are discharged from an opaque $L_8$~record (\texttt{Saturation/AbstractionBarrier9.agda}).
This machine-checks Sealing Encapsulation (\cref{rem:encapsulation}): resolved obligations are depth-1 references to the preceding layer, not depth-2 leaks to earlier layers.

\subsection{Coherence Window Stress-Testing}

The engine accepts a \texttt{--window d} flag parameterizing the Coherence Window.
$d = 1$ (constant costs) stagnates after 4--5 structures.
$d = 2$ (Fibonacci) produces the 15-structure PEN Synthesis Trajectory in the current fully coupled strict engine lane.
$d = 3$ (tribonacci) stalls earlier as costs outpace the horizon.
This provides computational evidence that, among the tested values $d \in \{1, 2, 3\}$, only $d = 2$ supports sustained evolution.

% ============================================
% SECTION 10: DISCUSSION
% ============================================
\section{Discussion}
\label{sec:discussion}

\subsection{What Is Proved, What Is Assumed, What Is Open}

\paragraph{Proved.}
\begin{itemize}[nosep]
    \item The Coherence Window Theorems (\cref{thm:ext-window,thm:int-window}): $d = 1$ for extensional foundations, $d = 2$ for the fully coupled CCHM/HoTT lane studied here.  The value $d = 2$ is derived from categorical first principles via three independent arguments: the Adjunction Barrier (\cref{thm:adjunction-barrier}, lower bound from triangle identities), Spectral Degeneration (\cref{thm:spectral-degeneracy}, upper bound from obligation spectral sequence collapse at $E_2$), and the Clutching Family (\cref{thm:clutching-bound}, tight lower bound from an infinite family of principal bundles).
    \item The Integration Trace Principle (\cref{lem:trace}): each layer exports $|S(L_k)| = \Delta_k$ schemas, verified by explicit obligation enumeration at steps~3--9, including the phase boundary from types to maps at step~9.  Machine-checked abstraction barrier at steps~8--9 in Cubical Agda.
    The Complexity Scaling Theorem (\cref{thm:scaling}) follows: $\Delta_n = F_n$ for $d = 2$.
    \item Sealing Encapsulation (\cref{rem:encapsulation}): resolved obligations become opaque $L_k$ exports; machine-checked at steps~8--9 in the Cubical Agda abstraction-barrier modules.
    \item The Divergence of Efficiency (\cref{thm:divergence}): $\rho_n \to \infty$, strengthened by the Logarithmic Effort Growth lemma (\cref{lem:log-effort}): combinatorial novelty $\Omega(n^c)$ dominates logarithmic specification cost $O(\log n)$.
    \item Schema Canonicality (\cref{prop:grammar-canonical}): depth-1 schemas are window-independent, verified computationally for steps 1--7.
    \item The $S^2 \equiv S^3$ Syntactic Identity (\cref{prop:s2s3}).
    \item \textbf{The Topological Projection} (\cref{thm:topological-projection}): The formula $\nu_H = m + d^2$ is derived from the cubical path algebra within the CCHM framework.  The derivation relies on framework-specific properties (connection maps, Kan operations, spectral degeneration at $E_2$) and should be read as a structural derivation within CCHM cubical type theory rather than a theorem of arbitrary cubical frameworks.
    \item The terminal DCT fixed-point signal (\cref{thm:spectral}): in the current bounded live atomic lane, the temporal-cohesive shell is selected at step~15 with $\nu=103$, $\kappa=8$, and the largest absolute overshoot in the sequence.
\end{itemize}

\paragraph{Assumed.}
\begin{itemize}[nosep]
    \item \textbf{Fully coupled evaluation lane} (\cref{thm:canonicity-preservation,thm:univalent-functoriality,rem:maximal-coupling}): PEN evaluates every candidate under maximal interface density. Within that lane, full coupling is forced by canonicity preservation for foundational constructors and by the economic-topological squeeze of univalent functoriality for API packages. What remains assumed is that the search stays in this monolithic, efficiency-optimized lane throughout.
    \item The Integration Trace Principle is verified for steps~3--9 and the abstraction barrier is machine-checked at steps~8--9, but a general proof for all $k$ remains open.  Steps~10--15 (modal operators and library specifications) are extrapolated: the obligation structure for axioms differs from HITs and maps, and verification of the trace principle for these steps remains open.
    \item The $\nu$ values in \cref{tab:genesis}: the canonical values are the machine totals produced by the current bounded live atomic run.  The expository rule audits for steps~11--15 remain subordinate to those totals, although individual components are now justified where named structural theorems apply (notably \cref{thm:telescopic-elimination} for explicit-reference API bundles).
    \item The DCT novelty claim is local to the current strict evaluator: the canonical statement is $\nu=103$ at conceptual cost $\kappa=8$, with the current code path exposing the standard split $\nu_G=2$, $\nu_H=15$, $\nu_C=86$.
    \item The executable lane is bounded: the current recovery artifact searches only the currently admissible atomic MBTT region through step~15 under fixed exact-band, connectivity, and retention controls.
\end{itemize}

\paragraph{Open.}
\begin{itemize}[nosep]
    \item \textbf{Specification clause boundaries for steps~1--4.}
    For steps~5--15, the clause count $\kappa$ is unambiguous: each clause corresponds to a single inference rule, and the specification selection (\cref{rem:kolmogorov-ambiguity}) is governed by the minimal-overshoot criterion.
    For the bootstrap steps~1--4 (Universe, Unit, Witness, $\Pi/\Sigma$), the clause boundaries depend on how one factors the foundational judgments.
    The values in \cref{tab:genesis} follow the convention that each atomic judgment form (formation, introduction, or computation rule) is one clause; the MBTT encoding audit (\cref{sec:mbtt-audit}) provides the explicit AST for verification.
    \item \textbf{How strong is the Step-15 claim?}
    The present claim is a structural terminal fixed-point theorem within PEN, not a universal novelty theorem over all existing mathematics.
    A stronger semantic novelty claim would require a broader comparative formalization program and an explicit external corpus criterion for structural equivalence classes.
    \item \textbf{Late-stage value decomposition.}
    For explicit-reference API bundles, the manuscript now proves the $\nu_C$ term by telescopic elimination (\cref{thm:telescopic-elimination}).
    What remains open is a fully uniform theorem-level decomposition of all late steps---especially mixed-shell cases and Step~15---into independent $\nu_G,\nu_C,\nu_H$ parts without reintroducing target-steered structural bonuses.
\end{itemize}

\paragraph{Conjecture versus implemented finding.}
The manuscript's strongest mathematical claim is a \emph{conjecture} about search-power, not a theorem about the current runtime profile: that sufficiently powerful MBTT-first search over the same admissible space would recover the same PEN Synthesis Trajectory without materially changing the present evaluator.
The \emph{implemented finding} is narrower and fully auditable: the current codebase recovers that trajectory in a bounded live atomic strict lane using exact-band admissibility, atomic MBTT enumeration, hybrid novelty evaluation, and semantic minimality filtering.
The paper should therefore be read as proving PEN's objective-level theorems and reporting a bounded live atomic recovery result for the current executable artifact.

\subsection{Limitations}

\paragraph{Scope of the Step-15 claim.}
The Step-15/DCT claim is intentionally local to the PEN setup: fixed objective, fixed admissibility, fixed bounded typed search, and fixed strict selection lane.
This is sufficient for a reproducible terminal-selection result within one principled synthesis regime, but not sufficient to conclude that no alternative objective or foundation could yield a different terminal structure.
The result should therefore be read as a strong mechanized existence claim in one principled synthesis regime.

\paragraph{Objective dependence.}
The decisive quantities $\nu$, $\kappa$, admissibility, and minimal overshoot are PEN's operational invariants of library growth, not canonical invariants of all possible foundational dynamics.
The manuscript's theorems are therefore conditional: once these quantities are fixed by constructive irreducibility, canonicity preservation, and universal coding, the induced trajectory is deterministic.

The topological projection $\nu_H = m + d^2$ is now a derived quantity (\cref{thm:topological-projection}), but the derivation relies on Spectral Degeneration at $E_2$ and the cubical interaction matrix, which are structural properties of the path algebra rather than consequences of the PEN axioms alone.
In the present manuscript that formula is used only for the explicit path-style payload preserved by the strict evaluator, where it functions as a theorem-backed shortcut rather than as a hand-tuned late-stage heuristic.

\paragraph{Component decomposition for axiomatic steps.}
For step~9 (Hopf fibration), $\nu_G = \nu_H = 0$ and $\nu = \nu_C$: as a map introducing no new type formers, all novelty is Elimination rules.
For steps 10--14 (library specifications), the decomposition is dominated by $\nu_C$ (Elimination schemas and structural interactions), with small but nonzero $\nu_G$ (2--5 Introduction schemas per step).
Step~15 (DCT) is different: the current evaluator canonizes the executable structural total $\nu = 103$ and exposes the standard split $\nu_G=2$, $\nu_H=15$, $\nu_C=86$.
The temporal shell is therefore the first late step whose full tripartite structure is now asserted directly from the executable artifact.

\paragraph{The current Step-15 computation.}
The selected Step-15 shell has canonical executable Generative Capacity
\[
  \nu(R_{15}) = 103.
\]
That value is produced by the current evaluator from the conceptual temporal shell together with its structural temporal-cohesive interaction payload.
Whenever the manuscript writes $\nu(\mathrm{DCT})=103$, it is using the DCT reading as semantic shorthand for this same executable Step-15 total.

\subsection{Mathematical Implications}

If the PEN Synthesis Trajectory faithfully identifies a locally optimal foundational trajectory under PEN, then cohesive and differential-geometric libraries arise as a consequence of efficiency optimization within the fully coupled CCHM/HoTT $d = 2$ coherence regime.
The sequence exhibits a coherent arc: dependent types $\to$ spheres $\to$ fibrations $\to$ cohesion $\to$ differential geometry $\to$ temporal-cohesive synthesis with derived infinitesimal semantics.

The absorption of arithmetic is instructive: natural numbers ($\rho \approx 1.5$) and Lie groups ($\rho = 1.50$) fail the rising bar and are subsumed by more efficient geometric and modal frameworks.
Within the model, efficient foundations prioritize geometric generality over discrete utility.

After DCT, PEN's external synthesis mechanism is squeezed: internal continuations are definitionally exhausted up to equivalence, while disconnected jumps are too inefficient to clear the Fibonacci-scaled bar.
This suggests a computational phase transition inside the model: the search shifts from external axiomatization to internal exploration within the DCT semantic completion.

\subsection{Falsifiability}

The PEN framework makes three testable predictions:

\begin{enumerate}[nosep]
    \item \textbf{Dimensional limit.}
    If a type-theoretic formalization in this target family required Class~3 coherence (non-trivial 3-paths not reducible to 2-path coherence), the claim that $d = 2$ for the fully coupled CCHM/HoTT lane would be falsified.
    \item \textbf{No early efficiency monsters.}
    If a structure existed with $\kappa < 4$ and $\nu > 20$ before homotopy theory, the sequence would be falsified.
    \item \textbf{The DCT novelty count.}
    $\nu(R_{15}) = 103$ is the canonical machine value in the current bounded live atomic lane.
    If independent analysis reduces this below the bar-clearing threshold $\kappa \cdot \mathrm{Bar}_{15} \approx 60.10$, the terminal fixed-point claim disappears and the framework is falsified.
\end{enumerate}

% ============================================
% SECTION 12: CONCLUSION
% ============================================
\section{Conclusion}
\label{sec:conclusion}

The Principle of Efficient Novelty determines a 15-step PEN Synthesis Trajectory---15 structures constituting an emergent geometric type-theoretic framework, derived from an initially empty accepted library under a fixed signature---governed by three theorems:
\begin{enumerate}[nosep]
    \item \textbf{Coherence Windows.} In the fully coupled CCHM/HoTT intensional lane analyzed here, $d = 2$; extensional type theory has $d = 1$.
    \item \textbf{Complexity Scaling.} The Integration Trace Principle---each layer's exported interface equals its integration trace---forces Fibonacci costs ($\Delta_n = F_n$) in that $d = 2$ regime; $d = 1$ stagnates.
    \item \textbf{Combinatorial Novelty.} Superlinear novelty growth ensures efficiency permanently outpaces the selection bar.
\end{enumerate}
and two structural results:
\begin{enumerate}[nosep, start=4]
    \item \textbf{Terminal temporal-cohesive fixed point.} In PEN's current bounded live atomic lane, the temporal-cohesive Step-15 AST is selected with $\nu=103$, $\kappa=8$, and maximal absolute overshoot; its semantic completion is read as DCT, and Step~16 then halts by the internal/external combinatorial squeeze.
    This identifies the terminal fixed point of PEN within the connected component explored by the engine.
    \item \textbf{Mechanical Verification.} The bounded live atomic engine recovers the full 15-step sequence with the machine values reported in \cref{tab:genesis}. Auxiliary audits validate portions of that computation but do not override those totals.
\end{enumerate}

The paper's final claim splits in two.
\emph{Conjecture:} with sufficiently powerful search over the same admissible MBTT space, the same structural PEN trajectory, including the same temporal-cohesive endpoint and the same semantic identifications, would emerge without the current tractability guidance.
\emph{Finding:} the present executable artifact recovers that trajectory in a bounded live atomic strict lane using exact-band admissibility, atomic MBTT enumeration, hybrid novelty evaluation, and semantic minimality filtering.
Taken together, these results support a conditional metatheory of typed library growth rather than an anthropomorphic discovery claim: under PEN's fixed MBTT vocabulary, admissibility policy, minimal-overshoot rule, and maximal-interface-density requirement, the fully coupled CCHM/HoTT $d=2$ coherence regime forces Fibonacci integration debt, superlinear novelty, asymptotic rejection of purely discrete families, and convergence to a Step-15 temporal-cohesive fixed point whose semantic completion is DCT and whose external continuation is halted inside the connected component of $\U_0$ explored by the engine.
Reproducibility details, extended arithmetic checks, boundary-case clarifications, and full MBTT audits are deferred to the appendices.

% ============================================
% APPENDIX: DCT DETAILS
% ============================================
\appendix

\section{Boundary Cases and Algorithmic Limits}
\label{sec:boundaries}

This appendix records three scope clarifications deferred from the main text, plus a short supplementary note linking back to the main-text discrete rejection argument.
They address reviewer-facing boundary questions, but are logically downstream of the core derivation.

\subsection{Algorithmic Criticality}
\label{sec:criticality}

The 0.11 margin at Step~6 does not indicate continuous hyperparameter tuning.
In the strict MBTT lane, the relevant quantities are discrete: $\kappa$ is an audited AST-clause count (\cref{def:effort,tab:mbtt-grammar-main}) and $\nu_H$ is locked by the cubical path algebra via $\nu_H = m + d^2$ (\cref{thm:topological-projection}).

\begin{theorem}[Rigidity of the $d=2$ critical window]
\label{thm:rigidity}
Within the strict MBTT lane, the survival region of the PEN Synthesis Trajectory is discrete rather than continuously tunable.
Hence the narrow Step-6 margin is a rigidity phenomenon, not a fitted threshold.
\end{theorem}

\begin{proof}[Proof sketch]
The sweep of \cref{rem:nuH-sensitivity} ranges over integer triplets, not real-valued hyperparameters.
Likewise, admissible changes to $\kappa$ are discrete clause insertions or deletions in a fixed grammar.
At Step~6, increasing the PropTrunc cost by one clause already moves $\rho_6$ from $8/3$ to $2.00 < 2.56$, killing the geometric branch.
The 0.11 margin is therefore the exact phase transition at which geometry first outpaces Fibonacci integration debt.
\end{proof}

\subsection{The Algorithmic Inefficiency of the Discrete}
\label{sec:discrete}

The main proof spine already establishes the discrete frontier in \cref{prop:discrete-asymptote-main,sec:discrete-frontier} and tabulates representative losing candidates in \cref{tab:rejected}.
This appendix records only the boundary-level reading: arithmetic and set-theoretic packages are not ignored by the evaluator, but rejected because in the fully coupled $d=2$ regime they lack the homotopy bonus and remain confined to linear novelty growth.

\subsection{The Univalent Horizon}
\label{sec:godelian}

The Step-15 halt theorem is local to the connected component of $\U_0$ explored by the guarded internal flow.
It is not a claim that mathematics globally halts.

\begin{definition}[The Univalent Horizon]
\label{def:godelian-horizon}
The \emph{Univalent Horizon} of a library $\mathcal{B}$ is the boundary of the connected component of $\mathrm{Mod}(\mathcal{B})$ reachable from $\mathcal{B}$ by internal continuous deformation.
Crossing that boundary requires an external axiomatic jump.
Within PEN, the halting theorem further shows that these jumps are not merely external but also inefficient once the Step-16 interface debt is charged.
\end{definition}

\begin{corollary}[Connected-component confinement]
\label{cor:confinement}
The DCT's guarded corecursion principle (\cref{lem:internalization}) generates structures only within the connected component of $\U_0$.
Higher universes, large-cardinal style axioms, and non-geometric logics lie beyond that component and are unreachable by internal flow alone.
\end{corollary}

This preserves G\"odelian open-endedness.
The claim is exhaustion of the largest connected fragment reachable from the initially empty accepted library under PEN's fixed signature, not exhaustion of mathematics as such.

\subsection{The DCT Novelty Count and Adjoint Completion}
\label{sec:adjoint-completion}

The DCT count $\nu=103$ is not a raw string count.
It is the canonical value produced by the current evaluator on the temporal shell selected at step~15.
Where explicit path-style payload is present, that evaluator uses the theorem-backed CCHM shortcut $\nu_H=m+d^2$ on the compiled witness basis rather than an ad hoc semantic bonus.
That total is the one used throughout the manuscript and in the executable selection loop.

Likewise, Adjoint Completion (\cref{lem:adjoint-completion}) is not a retrofitted scoring bonus.
It is a measurement principle forced by adjoint structure in categorical logic: introduction data determines corresponding elimination structure.
For bootstrap types that eliminative payload appears in $\nu_C$; for bare HITs it is absorbed into $\nu_H$ by dimensional orthogonality, which is why the HIT rows can have $\nu_C=0$ without contradicting adjoint completion.
In the current manuscript that principle is realized through the current evaluator and strict structural accounting.

\section{Artifact Map and Reproduction Checklist}
\label{app:artifact-map}

The executable artifact used for this manuscript is the \texttt{pen-atomic} workspace accompanying the paper.

\paragraph{Repository layout.}
\begin{itemize}[nosep]
    \item \texttt{crates/pen-core}, \texttt{pen-type}, \texttt{pen-eval}, \texttt{pen-search}: hot-path syntax, admissibility, evaluation, and live search.
    \item \texttt{crates/pen-store}, \texttt{pen-cli}, \texttt{pen-agda}: persisted artifacts, CLI orchestration, and observational Agda export.
    \item \texttt{scripts/compare\_runs.py}: deterministic cross-run comparison over stored artifacts.
    \item \texttt{tests/integration/}: end-to-end run, resume, comparison, and export checks.
    \item \texttt{tex/pen\_paper.tex}: manuscript source.
\end{itemize}

\paragraph{Minimal rerun checklist.}
\begin{enumerate}[nosep]
    \item \texttt{cargo test --workspace}
    \item \texttt{cargo run -p pen-cli -- run}\\
    \texttt{--config configs/debug.toml --root runs --run-id paper-run --until-step 15}
    \item \texttt{cargo run -p pen-cli -- resume runs/paper-run --until-step 15}
    \item \texttt{cargo run -p pen-cli -- inspect runs/paper-run}
    \item \texttt{cargo run -p pen-cli -- export-agda --run-dir runs/paper-run}
    \item \texttt{cargo run -p xtask -- export-reference-agda 15} \textit{(optional reference export)}
\end{enumerate}
These commands are sufficient to inspect the implementation, rerun live strict synthesis, exercise resume and export, and regenerate the reference Agda payload when needed.

\section{Worked Arithmetic and Rejected Candidates}
\label{app:rejections}

\subsection{Worked arithmetic for Steps 5--6}
\label{app:worked-arithmetic}

The transition from the Circle to Propositional Truncation is the tightest part of the trajectory and is therefore the most useful arithmetic checkpoint.
At Step~5, the Circle has $\Delta_5=F_5=5$, $\nu_5=7$, $\kappa_5=3$, and $\rho_5=7/3=2.33$.
The structural inflation is $\Phi_5=\Delta_5/\Delta_4=5/3=1.67$, the cumulative baseline is $\Omega_4=(1+1+2+5)/(2+1+1+3)=9/7=1.29$, and the bar is $\mathrm{Bar}_5=\Phi_5\Omega_4=2.14$.
Thus $S^1$ survives with margin $0.19$.

At Step~6, Propositional Truncation has $\Delta_6=F_6=8$, $\nu_6=8$, $\kappa_6=3$, and $\rho_6=8/3=2.67$.
Now $\Phi_6=8/5=1.60$, $\Omega_5=(1+1+2+5+7)/(2+1+1+3+3)=16/10=1.60$, and $\mathrm{Bar}_6=2.56$.
Hence Step~6 survives by only $0.11$.
If PropTrunc required one extra clause, its efficiency would fall to $2.00$ and the geometric branch would terminate immediately.

\subsection{Rejected candidates}

The summary table now appears in the main text (\cref{sec:discrete-frontier,tab:rejected}).
This appendix records only the supplementary interpretation notes used by the implementation discussion.

\paragraph{Discrete candidates.}
$\N$, classical logic, power-set style axioms, and ordinal arithmetic are all trapped by the same mechanism: they remain 0-groupoid or near-0-groupoid structures, so $\nu_H=0$ and novelty grows at most linearly in specification size, exactly as in \cref{prop:discrete-asymptote-main}.
For $\N$ in particular, addition and multiplication do not increase $\kappa$-worthy novelty because they are definitionally derivable from $\mathsf{rec}_\N$.

\paragraph{Generous arithmetic bundling.}
Even an intentionally over-generous arithmetic package with metric-scale clause count $\kappa=7$ still lacks the homotopy term.
Its best-case efficiency remains subcritical once the bar reaches the geometric regime, so geometry wins by topology-driven interaction growth rather than by curated bundling.

\paragraph{Higher but losing geometric or algebraic packages.}
Measure theory, elementary toposes, scheme theory, and Galois theory generate genuine structure but remain too discrete in the present scoring regime.
Symplectic geometry comes closest to survival, but by Step~13 it still falls below the bar cleared by metric structure.

\section{DCT: Key Theorems and Derived Construction Families}
\label{app:dct}

\subsection{Internal Tangent Bundle}

\begin{theorem}
\label{thm:tangent}
For any type $X : \U$ in DCT, the tangent bundle is:
$TX := \sum_{x : X} (X^{\D})_x$,
where $(X^{\D})_x$ is the type of infinitesimal curves through $x$.
The projection $\pi : TX \to X$ is a fiber bundle, and any flow on $X$ lifts to $TX$.
\end{theorem}

\subsection{Temporal Evolution}

\begin{theorem}
\label{thm:temporal}
Any type $X : \U$ in DCT admits a temporal evolution operator $E_X : \R \to (X \to X)$ satisfying identity, group property, smoothness, and discrete-structure preservation.
The next-modality $\bigcirc X$ is the type of fixed points of infinitesimal evolution.
\end{theorem}

\subsection{Symplectic Flow Corollary}

\begin{theorem}
\label{thm:symplectic-flow}
For a smooth type $M$ with symplectic form $\omega$:
any $H : M \to \R$ generates a unique symplectic vector field $X_H$;
the resulting flow preserves $\omega$;
the Poisson bracket makes $C^\infty(M)$ a Lie algebra.
\end{theorem}

\begin{corollary}[PDE-style evolution as internal flow]
Any linear PDE $\partial_t u = \mathcal{L}u$ on a smooth type is an instance of temporal evolution:
\[
E(t) = e^{t\mathcal{L}}.
\]
Heat, wave, and Schr\"odinger-type equations are distinguished by algebraic properties of~$\mathcal{L}$.
\end{corollary}

\subsection{Representative Formalization Families}

Temporal logic formulas, guarded recursion, geometric flow schemata, and symplectic transport schemata are all instances of DCT's temporal evolution operator.
Representative families include:
\begin{itemize}[nosep]
    \item \textbf{Classical mechanics:} on $T^*Q$, Hamiltonians generate flows obeying Hamilton's equations via \cref{thm:symplectic-flow}.
    \item \textbf{Gauge dynamics:} time-sliced Yang--Mills evolution is represented by flow on connection data coupled to cohesive transport structure.
    \item \textbf{Geometric flows:} Ricci/mean-curvature-style equations are instances of vector-field-generated evolution on moduli spaces.
    \item \textbf{Temporal verification:} LTL constructors ($\bigcirc,\Diamond$) and guarded recursion internalize safety/liveness-style reasoning.
\end{itemize}

\subsection{Relation to Existing Frameworks}

The ingredients of DCT each have established antecedents:
\begin{itemize}[nosep]
    \item Cohesive toposes provide $\flat,\sharp,\Pi,\Disc$ but no native temporal modality~\cite{lawvere,schreiber}.
    \item Guarded temporal type theory provides the ``next'' modality for recursion but not the cohesive semantic transfer studied here~\cite{nakano}.
    \item Synthetic differential geometry provides infinitesimals and tangent reasoning but not the temporal-cohesive exchange laws selected by the Step-15 AST~\cite{kock}.
\end{itemize}
What is specific to DCT in this manuscript is the unified typed synthesis of a temporal-cohesive AST together with the proof that its semantic completion carries internal tangent, temporal, and symplectic-flow consequences.

\subsection{Semantic Audit of DCT Novelty}
\label{tab:nu-audit}

The current manuscript canon is the executable value
\[
  \nu(\text{DCT}) = 103.
\]
This value is produced by the current evaluator and is the DCT novelty total used in the main text.
The same code path also exposes the standard split $\nu_G=2$, $\nu_H=15$, $\nu_C=86$.

\paragraph{What the current code supports.}
The selected temporal shell still justifies the same qualitative construction families:
\begin{itemize}[nosep]
    \item temporal endomorphism and guarded recursion patterns built from $\bigcirc$ and $\Diamond$,
    \item spatial--temporal exchange laws linking temporal operators to cohesive modalities,
    \item transport and flow-like interaction with the already selected cohesive and differential layers.
\end{itemize}
What changes is the measurement claim.
The paper now records the executable structural total together with the evaluator-exposed split rather than relying on a disconnected semantic bucketization.

\paragraph{Representative schema examples.}
The following shapes remain representative of the mathematical content carried by the temporal shell:
\begin{enumerate}[nosep]
\item \textbf{Temporal endomorphism flow:} $\Pi_{A : \U}\; \bigcirc(\flat A \to A)$, schematized as $\bigcirc(L \to L)$.
\item \textbf{Connection transport shape:} $\Pi_{A : \U}\; \sharp(\nabla(A)) \to \bigcirc(A)$, schematized as $\sharp(L) \to \bigcirc(L)$.
\item \textbf{Guarded recursion / step-indexing:} $\Pi_{A : \U}\; \bigcirc(A) \to A$, schematized as $\bigcirc(L) \to L$.
\end{enumerate}
These examples motivate the DCT interpretation, and the canonical novelty claim is the executable structural total $103$.

% ============================================
% APPENDIX: FORMAL TYPE SIGNATURES (STEPS 10-14)
% ============================================
\section{Post-Hoc Semantic Interpretations of AST Skeletons (Steps 10--14)}
\label{app:formal-specs}

When the synthesis engine selects the next bar-clearing candidate under the minimal-overshoot rule, the raw representation is an anonymous MBTT AST (\cref{sec:mbtt-audit}).
The algorithm itself does not assign semantic labels such as ``metric'' or ``Hilbert''; those names are post-hoc identifications.
This appendix documents that translation layer for Steps~10--14 by giving semantic interpretations of the selected AST skeletons and clarifying their axiomatic status.

The PEN Synthesis Trajectory's first nine steps---Universe through Hopf fibration---are standard constructions in HoTT, formalizable directly in Cubical Agda.
Steps 10--14 extend the type theory with differential-geometric structure.
This appendix therefore does \emph{not} claim that the MBTT clauses already encode the full semantic content of Riemannian or Hilbertian mathematics.
Rather, it records the minimal typed API shapes for which those readings are natural, and it explains how the corresponding inference-rule counts ($\nu_G$, $\nu_C$) arise in the manuscript's scoring model.

\paragraph{Axiomatic status.}
Steps 10--14 are \emph{library specifications} (\cref{rem:library-api}): packages of types, operations, and axioms that expand the library's derivation repertoire.
Unlike defined terms (which have $\nu = 0$ because they name existing derivable judgments), these specifications introduce structure that is \emph{not derivable} from the prior library.
Each step adds Introduction schemas (unit maps, constructors, type-formation side-effects) and Elimination schemas (counits, structural operations, cross-interactions).
The $\nu$ counts reflect the total number of derivation schemas natively unlocked by the specification, scored by the same methodology used for foundational steps~1--9.
The interpretive names below should therefore be read as semantic glosses on structural skeletons, not as evidence that the raw AST syntax independently reconstructs the entire downstream theory.

\subsection{Step 10: Cohesion (Adjoint Quadruple of Modalities)}
\label{app:cohesion-spec}

Following Lawvere~\cite{lawvere} and Schreiber~\cite{schreiber}, the cohesive structure extends HoTT with an adjoint string of modalities:
\[
    \Pi_{\!\infty} \;\dashv\; \flat \;\dashv\; \sharp
\]
where $\Pi_{\!\infty}$ is shape (fundamental $\infty$-groupoid), $\flat$ is flat (discrete), and $\sharp$ is sharp (codiscrete).

\paragraph{Type formation axioms ($\kappa = 4$ clauses).}
\begin{alignat*}{3}
    &\flat\text{-form:}   &\quad& A : \U \;\vdash\; \flat A : \U \\
    &\sharp\text{-form:}  && A : \U \;\vdash\; \sharp A : \U \\
    &\mathrm{Disc}\text{-form:} && A : \U \;\vdash\; \mathrm{Disc}(A) : \U \\
    &\Pi_{\mathrm{coh}}\text{-form:} && A : \U \;\vdash\; \Pi_{\!\infty}(A) : \U
\end{alignat*}

\paragraph{Derived inference rules ($\nu = 19$: $\nu_G = 2$, $\nu_C = 17$).}
The adjoint structure generates three families of rules:

\emph{(i) Unit/counit maps (4 rules: 2 Introduction, 2 Elimination):}
\begin{alignat*}{3}
    &\flat\text{-counit:}  &\quad& \flat A \to A  && \quad (\nu_C) \\
    &\sharp\text{-unit:}   && A \to \sharp A  && \quad (\nu_G) \\
    &\Pi_{\!\infty}\text{-unit:}   && A \to \Pi_{\!\infty}(A)  && \quad (\nu_G) \\
    &\mathrm{Disc}\text{-counit:}   && \mathrm{Disc}(\Gamma(A)) \to A  && \quad (\nu_C)
\end{alignat*}
The unit maps ($\eta_\sharp$, $\eta_{\Pi_\infty}$) are Introduction rules: they construct new term inhabitants.
The counit maps ($\varepsilon_\flat$, $\varepsilon_{\mathrm{Disc}}$) are Elimination rules: they extract underlying data from a modal type.

\emph{(ii) Modal structural rules (9 rules, all $\nu_C$):}
Idempotency ($\flat\flat A \simeq \flat A$, $\sharp\sharp A \simeq \sharp A$, $\Pi_{\!\infty}\Pi_{\!\infty} A \simeq \Pi_{\!\infty} A$); distribution over $\Pi$ and $\Sigma$ types ($\flat(\Pi_{(a:A)} B) \to \Pi_{(a:\flat A)} \flat B$, and analogously for $\sharp$, $\Pi_{\!\infty}$); these yield 3~idempotency + 6~distribution = 9 structural rules.

\emph{(iii) Cross-interactions with HITs (6 rules, all $\nu_C$):}
Each modality interacts with the library's higher inductive types ($S^1$, $S^2$, $S^3$, Hopf): $\flat(S^1) \simeq \mathbf{1}$ (the discrete circle is contractible), $\Pi_{\!\infty}(S^1) \simeq B\Z$ (the shape of the circle classifies $\Z$-torsors), and analogous interactions for the other HITs.
These produce 6~rules from 3~modalities $\times$ 2~independent HIT interactions.

\emph{Total:} $2\;(\nu_G) + 2 + 9 + 6\;(\nu_C) = 19$.

\subsection{\texorpdfstring{Step 11: Connections ($\nabla$ on Cohesive Types)}{Step 11: Connections (nabla on Cohesive Types)}}
\label{app:connections-spec}

A connection provides the differential-geometric machinery for parallel transport.
The formal content is a splitting of the de~Rham coefficient sequence in the cohesive setting.

\paragraph{Axioms ($\kappa = 5$ clauses).}
\begin{alignat*}{3}
    &\nabla\text{-form:}        &\quad& \nabla : \flat(TX) \to TX
        && \text{(connection operator)} \\
    &\mathrm{transport}_\nabla\text{:} && \mathrm{Id}_X(x, y) \to (F_x \simeq F_y)
        && \text{(parallel transport)} \\
    &\mathrm{cov}_\nabla\text{:}       && \flat(A) \to A
        && \text{(covariant derivative)} \\
    &\mathrm{lift}_\nabla\text{:}      && \flat(TX) \to TX
        && \text{(horizontal lift)} \\
    &\mathrm{Leibniz}\text{:}          && \nabla(f \cdot s) = df \otimes s + f \cdot \nabla s
        && \text{(product rule)}
\end{alignat*}

\paragraph{Derived inference rules ($\nu = 26$: $\nu_G = 2$, $\nu_C = 24$).}
Of the five charged clauses, two contribute genuinely new carrier families, so $\nu_G = 2$.
The cross-interface term is determined by the dominant imported carrier, namely cohesion from step~10. By the telescopic elimination law,
\[
  \nu_C = \nu_{10} + \kappa_{11} = 19 + 5 = 24.
\]
Total: $2\;(\nu_G) + 24\;(\nu_C) = 26$.

\paragraph{Why not definable.}
A connection on a general cohesive type $X$ is not derivable from the modalities $\flat$, $\sharp$ alone.
The de~Rham coefficient sequence $\flat X \to X \to {\Pi_{\!\infty}}_{\mathrm{dR}} X$ is exact, but a splitting of this sequence (i.e., a connection) is additional data requiring an axiom.

\subsection{\texorpdfstring{Step 12: Curvature ($R = d\nabla + \nabla \wedge \nabla$)}{Step 12: Curvature (R = d nabla + nabla wedge nabla)}}
\label{app:curvature-spec}

The curvature measures the failure of a connection to be flat.

\paragraph{Axioms ($\kappa = 6$ clauses).}
\begin{alignat*}{3}
    &R\text{-form:}             &\quad& R : \nabla \to \Omega^2(X)
        && \text{(curvature 2-form)} \\
    &\mathrm{Bianchi}\text{:}   && dR + [\nabla, R] = 0
        && \text{(Bianchi identity)} \\
    &\mathrm{hol}\text{:}       && \Omega(X, x) \to \mathrm{Aut}(F_x)
        && \text{(holonomy)} \\
    &\mathrm{CW}\text{:}        && R \to H^*(X)
        && \text{(Chern--Weil map)} \\
    &\mathrm{char}\text{:}      && \Omega^*(BG) \to H^*(X)
        && \text{(characteristic classes)} \\
    &R\text{-comp:}             && \nabla_1, \nabla_2 \mapsto R(\nabla_1 \circ \nabla_2)
        && \text{(composition)}
\end{alignat*}

\paragraph{Derived inference rules ($\nu = 34$: $\nu_G = 1$, $\nu_C = 33$).}
Only the curvature carrier itself contributes a genuinely new introduction family, so $\nu_G = 1$.
The shell depends on two imported layers, connections and cohesion, with the connection shell the dominant reference. By the telescopic elimination law,
\[
  \nu_C = \nu_{11} + \kappa_{12} + (2-1) = 26 + 6 + 1 = 33.
\]
Total: $1\;(\nu_G) + 33\;(\nu_C) = 34$.

\paragraph{Why not definable.}
While the curvature formula $R = d\nabla + \nabla \wedge \nabla$ is algebraically determined by a connection, the \emph{elimination rules} it generates (Bianchi identity, holonomy, Chern--Weil theory) require the curvature as a first-class object with its own structural rules.
These rules expand the derivation logic in ways not accessible from the connection alone.

\subsection{\texorpdfstring{Step 13: Endomorphic Operator Bundle (metric-equivalent reading)}{Step 13: Endomorphic Operator Bundle (metric-equivalent reading)}}
\label{app:metric-spec}

A direct inspection of the machine-verified MBTT AST for the Step-13 winner (\cref{tab:step13}) shows that the selected syntax for the leading clause is
\[
  \mathtt{AST}_g = \Sigma(\Pi(\mathrm{Var}_1, \mathrm{Var}_1), \Pi(\mathrm{Var}_1, \mathrm{Var}_1)),
\]
where $\mathrm{Var}_1$ resolves to the tangent space~$TX$.
The selected shell is therefore endomorphic and operator-theoretic before it is ever read as ``metric'' in the semantic layer.
This shell-first reading is the reason Step~13 survives the scalar objection.

\paragraph{Resolution of the Scalar Penalty: The Operator-Theoretic Bypass.}
A natural and mathematically precise objection to the Step~13 metric reading asks how the specification can formulate a symmetric bilinear form $g : TX \otimes_s TX \to \R$ without paying a large scalar penalty $\kappa_{\mathrm{scalar}}$ to discretely axiomatize the continuum and its commutative ring structure.
Attempting to extract $\R$ geometrically from the cohesive shape of the circle fails in standard cohesive HoTT: the circle's universal-cover/fundamental-groupoid route classifies $\mathbb{Z}$-torsors rather than supplying a free real field structure \cite{licata-shulman}.
Obtaining a genuine continuum object along that route requires extra real-cohesive axioms not selected by PEN \cite{shulman-real-cohesive}.

Furthermore, even if an additive continuum carrier were present, function composition is associative but not definitionally commutative.
Proving $x \cdot y = y \cdot x$ requires propositional equality witnesses, so a commutative ring presentation cannot be algorithmically free.
Under the $d=2$ Fibonacci constraints, any positive scalar penalty inflates the denominator and drives the Step-13 efficiency below the survival bar.

This objection is type-theoretically decisive, but it attacks an anthropomorphic reading that the engine itself did not adopt.
The winning shell therefore bypasses the commutative-continuum penalty entirely: the engine selected an operator-theoretic, endomorphic encoding rather than a literal scalar-valued bilinear form.

Within the semantic-interpretation layer of this appendix, such an endomorphic pair can be read as the forward and inverse musical maps that mediate metric behavior operationally.
The classical gloss ``$g : TX \otimes_s TX \to \R$'' is therefore a readability translation of the selected operator bundle, not a literal claim that PEN imported a commutative scalar field into the AST.
Consequently, $\kappa_{\mathrm{scalar}} = 0$ is not a claim of magical arithmetic deduction.
It is the narrower statement that the machine-selected MBTT shell pays no separate scalar-field import because the winning Step-13 syntax formalizes metric behavior endomorphically rather than through a discretely axiomatized continuum.

\paragraph{How $\R$ is typed (explicit dependency accounting).}
For completeness, if one nevertheless insists on a literal scalar-valued metric presentation, the metric codomain $\R$ is not available by default.
We evaluate explicit scalar routes under strict first-use charging: Cauchy-style arithmetic completion, topological arithmetic compression, and synthetic continuum insertion.
For a route with scalar burden $\kappa_{\mathrm{scalar}}$, Step~13 is scored by
\[
  \rho_{13}^{\mathrm{strict}} = \frac{46 + \Delta \nu_{\mathrm{mech}}}{7 + \kappa_{\mathrm{scalar}}},
  \qquad
  \Delta \nu_{\min} = \left\lceil 5.99\,(7+\kappa_{\mathrm{scalar}})-46\right\rceil.
\]

\paragraph{Axioms ($\kappa = 7$ clauses).}
The following display is the classical semantic reading of the seven selected clauses, not a claim that the raw MBTT syntax literally contains an explicit scalar-field object.
\begin{alignat*}{3}
    &g\text{-form:}           &\quad& g : TX \otimes_s TX \to \R
        && \text{(symmetric bilinear form)} \\
    &\mathrm{LC}(g)\text{:}  && g \mapsto \nabla_g
        && \text{(Levi-Civita connection)} \\
    &\gamma\text{-form:}     && [0,1] \to X
        && \text{(geodesic equation)} \\
    &\mathrm{vol}_g\text{:}  && g \mapsto \omega \in \Omega^n(X)
        && \text{(volume form)} \\
    &\star_g\text{:}          && \Omega^k(X) \to \Omega^{n-k}(X)
        && \text{(Hodge star)} \\
    &\Delta_g\text{:}         && \Omega^*(X) \to \Omega^*(X)
        && \text{(Laplace--Beltrami operator)} \\
    &\mathrm{Ric}_g\text{:}  && g \mapsto (\mathrm{Ric}, R_{\mathrm{scalar}}) \in \Omega^2(X) \times \R
        && \text{(Ricci tensor / scalar curvature)}
\end{alignat*}

\paragraph{Derivation schemas ($\nu = 46$: $\nu_G = 4$, $\nu_C = 42$).}
Of the 7~self-interaction schemas, 4 are Introduction-type (the metric $g$, Levi-Civita connection $\nabla_g$, geodesic equation, and volume form each construct new objects), while 3 are Elimination-type (Hodge star, Laplacian, Ricci extraction).
The remaining schemas are Elimination-type.
By the telescopic elimination law (\cref{thm:telescopic-elimination,ex:telescopic-elimination}), the structural interaction count is
\[
  \nu_C
  =
  \max_{\mathrm{ref}} \nu + \kappa + (\text{distinct refs} - 1)
  =
  34 + 7 + 1
  =
  42,
\]
because Metric directly imports Curvature and Connections, with Curvature the unique maximal imported interface.
Total: $4\;(\nu_G) + 42\;(\nu_C) = 46$.
These schemas are counted as \emph{derivation patterns natively unlocked} by the metric specification, not as inference rules of the core type theory (\cref{rem:library-api}).
The Hodge star, for example, becomes available only after the metric is specified; it is not derivable from prior structures.

\paragraph{Why not definable.}
In cohesive HoTT, types carry smooth structure via the adjoint modalities, but no intrinsic notion of distance or angle.
A metric is additional structure: the Levi-Civita theorem (uniqueness of the torsion-free metric connection) guarantees that each metric determines a unique connection, but the metric itself is not derivable from cohesion.

\paragraph{Derivability of components vs.\ $\N$'s addition.}
A natural objection asks why the Levi-Civita connection, Hodge star, Laplacian, and Ricci scalar are counted as specification clauses rather than derivable consequences with $\nu = 0$---the standard applied to $\N$'s addition from $\mathrm{rec}_\N$.
The distinction is \emph{constructive availability within the type theory}.
Addition on~$\N$ is a definitional term: $\mathrm{add} := \mathrm{rec}_\N\,(\lambda n.\,\mathrm{succ}\,n)$ is a closed expression in MLTT requiring no additional axioms.
By contrast, the Levi-Civita connection is the solution to a system of partial differential equations (torsion-freeness + metric compatibility) whose existence and uniqueness require analytic machinery---the Fundamental Theorem of Riemannian Geometry---that is not constructively available from the bare metric tensor $g$ within cohesive HoTT.
The Hodge star requires extending the inner product induced by~$g$ from~$TX$ to the full exterior algebra $\Lambda^k T^*X$, a construction that depends on the orientation and dimension of~$X$; the Laplacian and Ricci scalar each require the connection as input.
These operations form a \emph{dependency chain}---$g \to \nabla_g \to R \to \mathrm{Ric}$---in which each link requires existence theorems that are not reducible to a single recursive definition.
The specification bundles this chain as a coherent API package (\cref{rem:admissibility-constraints}, structural unity): the components are interdependent and cannot be meaningfully decomposed.
A stripped-down specification with $\kappa = 1$ (bare $g$-tensor alone) would have $\nu_G = 1$, $\nu_C \leq 2$, giving $\rho \leq 3.0$---below the Step~13 bar of $5.99$.
The metric clears the bar precisely because its API specification natively unlocks the full chain of derived geometric operations as new derivation schemas.

\subsection{Step 14: Hilbert Functional (Operator Algebra)}
\label{app:hilbert-spec}

The Hilbert functional introduces inner-product structure, completeness, and spectral theory---a high-interaction analytic package over the geometric library.

\paragraph{Axioms ($\kappa = 9$ clauses).}
\begin{alignat*}{3}
    &\langle\cdot,\cdot\rangle\text{-form:} &\quad& V \times V \to \R
        && \text{(inner product)} \\
    &\mathrm{complete}\text{:}               && \mathrm{Cauchy}(V) \to V
        && \text{(Cauchy completeness)} \\
    &\perp\text{-decomp:}                    && V \to V_1 \oplus V_2
        && \text{(orthogonal decomposition)} \\
    &\mathrm{spec}\text{:}                   && \mathrm{Op}(V) \to \Sigma(\lambda : \R).\, V_\lambda
        && \text{(spectral decomposition)} \\
    &C^*\text{-alg:}                         && (V \to V) \times (V \to V) \to (V \to V)
        && \text{($C^*$-algebra structure)} \\
    &g\text{-compat:}                        && \mathrm{Metric} \to \langle\cdot,\cdot\rangle
        && \text{(metric compatibility)} \\
    &R\text{-op:}                            && \mathrm{Curv} \to \mathrm{Op}(V)
        && \text{(curvature operator)} \\
    &\nabla\text{-op:}                       && \mathrm{Conn} \to \mathrm{Op}(V)
        && \text{(connection operator)} \\
    &\delta/\delta\phi\text{:}               && (V \to \R) \to V
        && \text{(functional derivative)}
\end{alignat*}

\paragraph{Derived inference rules ($\nu = 62$: $\nu_G = 5$, $\nu_C = 57$).}
Of the spectral-theoretic rules, 5 are Introduction-type (inner product, completeness/limits, orthogonal decomposition, spectral decomposition, and the functional derivative each construct new objects).
The remaining 57 rules are Elimination schemas.
By \cref{thm:telescopic-elimination,ex:telescopic-elimination}, the structural interaction count is
\[
  \nu_C
  =
  \max_{\mathrm{ref}} \nu + \kappa + (\text{distinct refs} - 1)
  =
  46 + 9 + 2
  =
  57,
\]
because Hilbert directly imports Metric, Curvature, and Connections, with Metric the unique maximal imported interface.
Total: $5\;(\nu_G) + 57\;(\nu_C) = 62$.

\paragraph{Why not definable.}
Hilbert-space structure---inner products, completeness, spectral theory---requires analytic axioms beyond the algebraic and geometric content of steps 10--13.
The functional derivative $\delta/\delta\phi$ provides variational-calculus machinery that couples operator structure to the metric/connection/curvature layer.

\paragraph{Selection margin at Step~14.}
The Hilbert functional clears the selection bar by $0.21$ (\cref{tab:genesis}): $\rho = 62/9 = 6.89$ against Bar $= 6.68$.
While not the tightest margin (that distinction belongs to Propositional Truncation at step~6 with margin $0.11$), it remains structurally significant: the analytic content of Hilbert-space theory ($\nu = 62$) survives only because its 9~specification clauses achieve sufficient efficiency.
If the analytic axioms required even two additional specification clauses ($\kappa = 11$, giving $\rho \approx 5.64 < 6.68$), the sequence would permanently stagnate at Riemannian geometry without reaching this operator-analytic layer.

% ============================================
% APPENDIX: MBTT ENCODING AUDIT
% ============================================
\section{The MBTT Encoding Audit}
\label{sec:mbtt-audit}

This appendix opens the $\kappa$ ``black box'' by reporting the encoding-specific audit data for all 15 genesis steps.
The shared MBTT grammar is fixed in \cref{tab:mbtt-grammar-main}; here we restrict attention to prefix-cost conventions, Elias-gamma pointer costs, and the full per-step bit-length audit.
The prefixes in \cref{tab:mbtt-grammar-main} satisfy the Kraft--McMillan inequality by construction.

\noindent\textbf{Elias Gamma coding.}
For a positive integer $n$, $|\gamma(n)| = 2\lfloor \log_2 n \rfloor + 1$ bits.
Examples: $\gamma(1) = 1$~bit; $\gamma(2)\text{--}\gamma(3) = 3$~bits; $\gamma(4)\text{--}\gamma(7) = 5$~bits; $\gamma(8)\text{--}\gamma(15) = 7$~bits.
Library pointers use 1-based indexing: $\textsc{Lib}(1)$ references the first library entry (Universe) and costs $3 + 1 = 4$~bits; $\textsc{Lib}(12)$ references Curvature and costs $3 + 7 = 10$~bits.

\subsection{Full Audit: Clause Counts vs.\ MBTT Bit-Lengths}

\begin{table}[tbp]
\centering
\caption{Complete encoding audit.  $\kappa$ is the clause count (used in the main text); MBTT is the total bit-length of all clauses in the prefix-free encoding.  The specification shown for each step is the one \emph{selected} by the minimal-overshoot criterion (\cref{ax:selection}).}
\label{tab:mbtt-audit}
\scriptsize
\begin{tabular}{@{}cr lrr >{\raggedright\arraybackslash}p{3.4cm}@{}}
\toprule
$n$ & $\kappa$ & Selected AST skeleton & Clauses & MBTT (bits) & Selected specification \\
\midrule
1  & 2 & Universe           & 1 & 4   & $\U$-formation \\
2  & 1 & Unit               & 1 & 10  & $\mathbf{1}$-formation \\
3  & 1 & Witness            & 2 & 27  & $\star$ + $\mathrm{ind}_1$ \\
4  & 3 & $\Pi/\Sigma$       & 5 & 93  & $\lambda$, pair, app, fst, snd \\
5  & 3 & $S^1$              & 3 & 21  & base + loop \\
6  & 3 & PropTrunc          & 3 & 35  & $\|\cdot\|_0$ + squash \\
7  & 3 & $S^2$              & 3 & 23  & base + surf$^\dagger$ \\
8  & 5 & $S^3$              & 5 & 66  & H-space: base, 3-cell, $\mu$, unit$^\dagger$ \\
9  & 4 & Hopf               & 4 & 78  & $h : S^3 \to S^2$ + fiber + section \\
10 & 4 & Modal shell        & 4 & 45  & $\flat, \sharp, \mathrm{Disc}, \Pi_\mathrm{coh}$ \\
11 & 5 & Connection shell   & 5 & 79  & $\nabla$, transport, covariant lift, Leibniz \\
12 & 6 & Curvature shell    & 6 & 121 & $R$, Bianchi, holonomy, Chern--Weil, char.\ class \\
13 & 7 & Endomorphic bundle & 7 & 138 & See \cref{tab:step13} \\
14 & 9 & Hilbert-functional shell & 9 & 169 & See \cref{tab:step14} \\
15 & 8 & Temporal-cohesive shell & 8 & 229 & See \cref{tab:step15} \\
\bottomrule
\end{tabular}
\end{table}

\noindent ${}^\dagger$For $S^2$ and $S^3$, shorter suspension encodings exist ($\Sigma S^1$ and $\Sigma S^2$, each 13~bits / 1~clause), but the minimal-overshoot criterion (\cref{ax:selection}) selects the native or enriched specifications because they yield tighter efficiency matches to the bar (\cref{rem:kolmogorov-ambiguity}).

\medskip
\noindent The MBTT bit-lengths range from 4 to 229 across the PEN Synthesis Trajectory, while clause counts range from 1 to 9.
The two measures are positively correlated (Spearman $r \approx 0.87$) but not identical: steps that reference deep library entries (like DCT at 229~bits) have disproportionately large bit-lengths due to Elias gamma costs.
The main text uses clause counts for $\kappa$ because they are encoding-independent; the bit-lengths serve as an independent verification layer.

\subsection{Per-Clause Breakdown: Steps 13--15}

We now provide the explicit AST and bit-cost for the three steps with the tightest selection margins.

\begin{table}[tbp]
\centering
\caption{Step 13 selected AST skeleton: endomorphic operator bundle ($\kappa = 7$ clauses, 138~bits total).  Library references: Connections at index 11 ($\textsc{Lib}(11) = 10$~bits), Curvature at index 12 ($\textsc{Lib}(12) = 10$~bits).  The classical metric reading is post-hoc and discussed in \cref{app:metric-spec}.}
\label{tab:step13}
\small
\begin{tabular}{@{}clrl@{}}
\toprule
\# & Clause & AST & Bits \\
\midrule
1 & Endomorphic pair shell           & $\Sigma(\Pi(\mathrm{Var}_1, \mathrm{Var}_1), \Pi(\mathrm{Var}_1, \mathrm{Var}_1))$ & 26 \\
2 & Levi-Civita connection           & $\Pi(\Sigma(\mathrm{Var}_1, \mathrm{Var}_2), \textsc{Lib}(11))$ & 27 \\
3 & Geodesic equation                & $\Pi(\mathrm{Var}_1, \Pi(\mathrm{Var}_1, \mathrm{Var}_1))$ & 18 \\
4 & Volume form                      & $\mathrm{Lam}(\mathrm{App}(\mathrm{Var}_1, \mathrm{Var}_2))$ & 14 \\
5 & Hodge star                       & $\Pi(\textsc{Lib}(12), \textsc{Lib}(12))$ & 23 \\
6 & Laplacian                        & $\mathrm{Lam}(\Pi(\mathrm{Var}_1, \mathrm{Var}_1))$ & 13 \\
7 & Ricci scalar                     & $\Pi(\textsc{Lib}(12), \mathrm{Var}_1)$ & 17 \\
\midrule
  & \textbf{Total} & & \textbf{138} \\
\bottomrule
\end{tabular}
\end{table}

\begin{table}[tbp]
\centering
\caption{Step 14: Hilbert-functional shell ($\kappa = 9$ clauses, 169~bits total).  Library references: Connection shell at index 11, Curvature shell at index 12, endomorphic operator bundle at index 13 ($\textsc{Lib}(13) = 10$~bits).}
\label{tab:step14}
\small
\begin{tabular}{@{}clrl@{}}
\toprule
\# & Clause & AST & Bits \\
\midrule
1 & Inner product               & $\Sigma(\Pi(\mathrm{Var}_1, \Pi(\mathrm{Var}_1, \U)), \mathrm{Var}_1)$ & 26 \\
2 & Completeness                & $\Pi(\mathrm{Var}_1, \mathrm{Var}_1)$ & 11 \\
3 & Orthogonal decomposition    & $\Pi(\mathrm{Var}_1, \Sigma(\mathrm{Var}_1, \mathrm{Var}_1))$ & 19 \\
4 & Spectral decomposition      & $\Pi(\mathrm{Lam}(\mathrm{Var}_1), \Sigma(\mathrm{Var}_1, \mathrm{Var}_2))$ & 23 \\
5 & $C^*$-algebra structure     & $\Sigma(\Pi(\mathrm{Var}_1, \mathrm{Var}_1), \Pi(\mathrm{Var}_1, \mathrm{Var}_1))$ & 26 \\
6 & Metric compatibility        & $\Pi(\textsc{Lib}(13), \mathrm{Var}_1)$ & 17 \\
7 & Curvature operator          & $\Pi(\textsc{Lib}(12), \mathrm{Var}_1)$ & 17 \\
8 & Connection lift             & $\Pi(\textsc{Lib}(11), \mathrm{Var}_1)$ & 17 \\
9 & Functional derivative       & $\mathrm{Lam}(\Pi(\mathrm{Var}_1, \U))$ & 13 \\
\midrule
  & \textbf{Total} & & \textbf{169} \\
\bottomrule
\end{tabular}
\end{table}

\begin{table}[tbp]
\centering
\caption{Step 15: temporal-cohesive shell (semantic DCT completion) ($\kappa = 8$ clauses, 229~bits total).  Library reference: modal shell at index 10 ($\textsc{Lib}(10) = 10$~bits).  The audited AST contains temporal operators and exchange clauses only; the infinitesimal reading is derived semantically rather than encoded as a primitive node.}
\label{tab:step15}
\small
\begin{tabular}{@{}clrl@{}}
\toprule
\# & Clause & AST & Bits \\
\midrule
1 & Next modality $\bigcirc X$           & $\bigcirc(\mathrm{Var}_1)$ & 13 \\
2 & Eventually modality $\Diamond X$     & $\Diamond(\mathrm{Var}_1)$ & 13 \\
3 & $\bigcirc \to \Diamond$ axiom        & $\Pi(\bigcirc(\mathrm{Var}_1), \Diamond(\mathrm{Var}_1))$ & 29 \\
4 & Spatial--temporal interaction         & $\mathrm{Lam}(\mathrm{App}(\textsc{Lib}(10), \bigcirc(\mathrm{Var}_1)))$ & 27 \\
5 & $\flat\bigcirc \leftrightarrow \bigcirc\flat$ exchange & $\Pi(\flat(\bigcirc(\mathrm{Var}_1)), \bigcirc(\flat(\mathrm{Var}_1)))$ & 43 \\
6 & $\sharp\Diamond \leftrightarrow \Diamond\sharp$ exchange & $\Pi(\sharp(\Diamond(\mathrm{Var}_1)), \Diamond(\sharp(\mathrm{Var}_1)))$ & 43 \\
7 & $\Diamond$-elimination               & $\mathrm{Lam}(\mathrm{App}(\Diamond(\mathrm{Var}_1), \mathrm{Var}_2))$ & 23 \\
8 & $\bigcirc\bigcirc \to \bigcirc$ witness & $\Pi(\bigcirc(\bigcirc(\mathrm{Var}_1)), \bigcirc(\mathrm{Var}_1))$ & 38 \\
\midrule
  & \textbf{Total} & & \textbf{229} \\
\bottomrule
\end{tabular}
\end{table}

\subsection{Canonicality of the Encoding}

The MBTT grammar is not gerrymandered to produce the PEN Synthesis Trajectory.
Three properties establish this:

\begin{enumerate}[nosep]
\item \textbf{Standard prefix-free design.}
The prefix lengths follow the standard entropy-coding principle: constructors that appear most frequently in dependent type theory (\textsc{App}, \textsc{Lam}, \textsc{Pi}) receive the shortest prefixes (2--3~bits), while domain-specific operators (modal, temporal) receive longer ones (7--9~bits).
This is the same design principle used in Huffman and arithmetic codes.
Crucially, the modal/temporal nodes are not ``rewarded'' by this encoding; they are assigned \emph{higher} local description cost than core constructors.

\item \textbf{Separation of concerns.}
The main text uses $\kappa$ = clause count, which is entirely encoding-independent.
The MBTT bit-lengths serve only as an independent audit: they verify that no clause is degenerate (encoding a trivial operation in a complex AST) and that the total specification complexity is monotonically increasing across the PEN Synthesis Trajectory.

\item \textbf{Reproducibility.}
The complete MBTT encoding is implemented in \texttt{Kolmogorov.hs} (under \texttt{engine/src/}).
The manuscript's MBTT audit table is retained as a deterministic representation audit in the current repository.
Because all candidate families are scored in this single grammar, the result is an apples-to-apples optimum test over combinations, not a hand-picked proof that one pre-named structure must win.
\end{enumerate}

\subsection{Search-Space Tractability and Strict Recovery}
\label{sec:search-tractability}

A recurring objection is that the MBTT audit lengths (e.g., 138~bits for Metric, 229~bits for DCT) imply an infeasible raw space of $2^{138}$ or $2^{229}$ binary strings.
That objection would be valid for blind bitstring search; it does not describe the implemented engine.
The runtime domain is not ``all admissible MBTT ASTs'' at once, but a bounded exact-band search over the currently admissible anonymous MBTT region.

The runtime search is constrained by construction:
\begin{enumerate}[nosep]
    \item \textbf{Typed manifold, not raw strings.} Candidates are generated as well-formed MBTT telescopes, not as unconstrained binaries.
    \item \textbf{Exact-band control.} Enumeration runs under strict admissibility caps and prioritized exact $\kappa$-bands rather than over the full raw grammar at once.
    \item \textbf{Atomic generation.} Enumeration searches anonymous MBTT telescopes directly inside the current admissible structural lane.
    \item \textbf{Structural pruning.} Ill-typed and inadmissible telescopes are removed by type checks, derivability checks, and semantic minimality before final selection.
    \item \textbf{Deterministic quotienting.} Canonical-key deduplication and deterministic tie-breaking keep the live lane reproducible.
    \item \textbf{Bounded frontier evidence.} Retained valid candidates, frontier checkpoints, and resume artifacts are persisted without changing acceptance truth.
    \item \textbf{Diagnostics separated from canon.} Auxiliary diagnostics may still be evaluated for validation, but the manuscript's canonical values are set by the strict hybrid evaluator.
\end{enumerate}

Hence the implementation is a bounded live atomic strict lane over an admissible, bounded, and reproducible candidate manifold---not exhaustive brute-force over the entire binary code space.
The executable claim is therefore an objective optimum test over this bounded atomic basis, not a claim that the engine searched the full binary ambient space or invented its own late-stage semantic vocabulary.
The MBTT bit-lengths in \cref{tab:mbtt-audit} are post-hoc representation audits of selected specifications, not the dimensionality of the runtime enumeration domain.

\section{A Detailed Proof of the 15-Step Genesis Sequence}
\label{sec:genesis-sequence-proof}

This section gives a step-by-step proof of the current 15-step PEN genesis
sequence in a form intended for a journal proof spine rather than for a
developer report.
The global mathematical input comes from the established PEN theorems in
\cref{thm:scaling,ax:selection,def:derivability,def:minimal-api,thm:end-of-history},
while the step-local selection facts come from the verified strict trace
generated by
\texttt{cargo run -p pen-cli -- run --config configs/debug.toml --until-step 15}.
By \cref{rem:automation-scope}, that strict run is the canonical source for the
selection ordering, the bar clearances, and the quantitative values
$(\nu,\kappa,\rho)$ used below.

% Keep proof headings input-safe across LaTeX setups without changing the host
% document's sectioning semantics after this file ends.
\let\genesissequenceorigparagraph\paragraph
\renewcommand{\paragraph}[1]{\par\medskip\noindent\textbf{#1.}\ }

\subsection{Proof Contract}
\label{subsec:genesis-proof-contract}

Three conventions govern the proof.

\begin{enumerate}[nosep]
    \item \textbf{Charged cost.} Throughout this section, $\kappa$ means the
    charged conceptual/desugared cost used by the strict selector, not the raw
    size of every compiled auxiliary clause.
    This is the cost notion fixed in \cref{def:effort,ax:selection} and exposed
    operationally in the strict engine description of
    \cref{sec:verification}.
    \item \textbf{Compiled payload.} Especially for the HIT steps, the compiled
    telescope may contain framework-derived eliminators or computation payload
    beyond the charged conceptual shell.
    Those extra compiled clauses can increase the evaluated novelty, but they do
    not increase the charged $\kappa$.
    This is the conceptual/compiled separation described in
    \cref{sec:verification}.
    \item \textbf{Uniqueness.} Every uniqueness claim below means
    ``unique winner in the current bounded live atomic lane under admissibility,
    semantic minimality, and minimal overshoot''.
    It does \emph{not} assert a foundation-independent uniqueness theorem over
    all logics or all search policies.
    \item \textbf{Shell-first naming.} For steps~10--15, the quantitative claim
    attaches to the selected AST skeleton first.
    Human names such as ``metric'' and ``DCT'' are retained only as semantic
    readings when they help orient the discussion.
\end{enumerate}

Two further scope boundaries matter for the late stages.
First, the codebase contains additional family-specific recovery and quality
machinery in non-canonical lanes, but the present section does not use those
heuristics as evidence for the strict claim.
Whenever the verified strict trace already yields a singleton viable set, we use
that stronger fact directly.
Second, the proof-strength boundary from \cref{rem:automation-scope} is
maintained verbatim:
steps~1--10 admit exact or exact-reconstruction language,
steps~11--14 are expository reconstructions constrained to match the strict
totals, and step~15 canonizes the executable structural total $\nu=103$
attached to the selected temporal-cohesive shell.

\subsection{A Global Bar Lemma}
\label{subsec:genesis-bar-lemma}

\begin{lemma}[Bar computation along the strict genesis trace]
\label{lem:genesis-bar-lemma}
For every step $n \ge 2$ in the strict genesis trace,
\[
  \mathrm{Bar}_n
  =
  \Phi_n \Omega_{n-1}
  =
  \frac{F_n}{F_{n-1}}
  \cdot
  \frac{\sum_{i=1}^{n-1}\nu_i}{\sum_{i=1}^{n-1}\kappa_i}.
\]
Consequently the active bars for steps~$2$ through~$15$ are exactly the values
listed in \cref{tab:genesis-sequence-trace}.
\end{lemma}

\begin{proof}
By \cref{thm:scaling}, the current $d=2$ lane satisfies $\Delta_n=F_n$, hence
\[
  \Phi_n = \Delta_n/\Delta_{n-1} = F_n/F_{n-1}.
\]
By \cref{def:omega},
\[
  \Omega_{n-1} = \frac{\sum_{i=1}^{n-1}\nu_i}{\sum_{i=1}^{n-1}\kappa_i}.
\]
By \cref{ax:selection}, the selection bar is
$\mathrm{Bar}_n = \Phi_n \Omega_{n-1}$.
Substituting the already realized strict-trace values
$(\nu_i,\kappa_i)$ for $i<n$ therefore determines each bar uniquely.
The worked arithmetic in \cref{app:worked-arithmetic} checks this explicitly for
steps~5 and~6, and the same substitution yields the full bar column in
\cref{tab:genesis-sequence-trace}.
\end{proof}

\begin{table}[tbp]
\centering
\caption{Strict genesis trace used in the proof.
The `cap/bands' column reports the strict admissibility cap together with the
searched exact $\kappa$-band(s).
The `raw/dist.' column reports raw and deduplicated candidate counts.
`Runner-up' and `near miss' mean the top displayed strict-lane competitor in the
bar-clearing and subcritical pools respectively; a dash means that no distinct
competitor was displayed by the verified strict trace.}
\label{tab:genesis-sequence-trace}
\tiny
\setlength{\tabcolsep}{2pt}
\begin{tabular}{@{}c l r r r r r r l l >{\raggedright\arraybackslash}p{1.6cm} >{\raggedright\arraybackslash}p{1.6cm}@{}}
\toprule
Step & Selected AST skeleton & $\Delta_n$ & $\nu$ & $\kappa$ & $\rho$ & Bar & Margin & cap/bands & raw/dist. & Runner-up & Near miss \\
\midrule
1  & Universe    & 1   & 1  & 2 & 0.50 & 0.50 & 0.00 & 2 / 2   & 2144 / 552 & --- & --- \\
2  & Unit        & 1   & 1  & 1 & 1.00 & 0.50 & 0.50 & 1 / 1   & 96 / 78    & --- & --- \\
3  & Witness     & 2   & 2  & 1 & 2.00 & 1.33 & 0.67 & 1 / 1   & 96 / 81    & --- & found.\ cand.\ (1.00) \\
4  & $\Pi/\Sigma$ & 3  & 5  & 3 & 1.67 & 1.50 & 0.17 & 3 / 3   & 1 / 1      & --- & --- \\
5  & $S^1$       & 5   & 7  & 3 & 2.33 & 2.14 & 0.19 & 3 / 3   & 1 / 1      & --- & --- \\
6  & PropTrunc   & 8   & 8  & 3 & 2.67 & 2.56 & 0.11 & 4 / 3   & 3 / 3      & $S^2$ (3.33) & $S^1$ (2.33) \\
7  & $S^2$       & 13  & 10 & 3 & 3.33 & 3.00 & 0.33 & 4 / 3   & 2 / 2      & --- & $S^1$ (2.33) \\
8  & $S^3$       & 21  & 18 & 5 & 3.60 & 3.43 & 0.17 & 5 / 5   & 1 / 1      & --- & --- \\
9  & Hopf        & 34  & 17 & 4 & 4.25 & 4.01 & 0.24 & 5 / 4   & 1 / 1      & --- & --- \\
10 & Modal shell & 55  & 19 & 4 & 4.75 & 4.46 & 0.29 & 5 / 4   & 2 / 2      & --- & Hopf (4.25) \\
11 & Connection shell & 89  & 26  & 5 & 5.20  & 4.91 & 0.29 & 6 / 5   & 2 / 2      & --- & $S^3$ (3.60) \\
12 & Curvature shell & 144 & 34  & 6 & 5.67  & 5.42 & 0.24 & 6 / 5|6 & 3 / 3      & --- & $S^3$ (3.60) \\
13 & Endomorphic bundle & 233 & 46  & 7 & 6.57  & 5.99 & 0.58 & 8 / 7   & 3 / 3      & --- & HIT (2.29) \\
14 & Hilbert-functional shell & 377 & 62  & 9 & 6.89  & 6.68 & 0.21 & 9 / 8|9 & 4 / 4      & --- & HIT (3.00) \\
15 & Temporal-cohesive shell & 610 & 103 & 8 & 12.88 & 7.40 & 5.48 & 9 / 8   & 3 / 3      & --- & HIT (3.00) \\
\bottomrule
\end{tabular}
\end{table}

\subsection{Bootstrap: Steps 1--4}
\label{subsec:genesis-bootstrap-proof}

The bootstrap steps establish the minimum typed infrastructure from which later
homotopical and geometric structure can be sealed.
The foundational rule counts used here are the exact strict values recorded in
\cref{rem:automation-scope,tab:genesis-sequence-trace,tab:mbtt-audit}.

\begin{proposition}[Step 1: Universe]
\label{prop:genesis-step1}
At Step~1 the strict genesis trace selects the universe package.
\end{proposition}

\begin{proof}
\paragraph{What it is.}
The selected structure is the first universe $\U_0$, namely the ambient type of
small types.
In the strict trace it is the unique selected bootstrap object at step~1; see
\cref{tab:genesis-sequence-trace}.

\paragraph{Why $\kappa$ has this value.}
The charged cost is $\kappa=2$, exactly the two foundational clauses audited in
\cref{tab:mbtt-audit}: universe formation and the universe-level witness.
This is the minimal non-empty package that realizes a type universe rather than
an opaque constant, so the cost is fixed by \cref{def:effort}.

\paragraph{How it builds on the established library.}
It builds on nothing, because the prior library is empty.
Precisely for that reason it is the only admissible way to create a typed growth
regime at all: every later candidate presupposes some ambient type formation.

\paragraph{Why the novelty count has this value.}
Its novelty is the exact foundational count $\nu=1$ from
\cref{rem:automation-scope}: one new formation rule for small types.
There is no earlier eliminator to interact with, so no additional
cross-interface novelty is available at this step.

\paragraph{Why it is the unique selected winner at this step.}
The strict trace opens cap/band $2/2$, evaluates $2144$ raw candidates, and
selects Universe with $\rho=0.50$ exactly at the initial bar $0.50$; see
\cref{tab:genesis-sequence-trace}.
By the bootstrap admissibility conditions encoded in the strict search and by
\cref{ax:selection}, this is the unique least-supercritical bootstrap package.
\end{proof}

\begin{proposition}[Step 2: Unit]
\label{prop:genesis-step2}
At Step~2 the strict genesis trace selects the unit type $\mathbf{1}$.
\end{proposition}

\begin{proof}
\paragraph{What it is.}
The selected structure is the terminal one-point type $\mathbf{1}$.
It is the first concrete small type over $\U_0$, and therefore the first place
where inhabitation can later be witnessed.

\paragraph{Why $\kappa$ has this value.}
The charged cost is $\kappa=1$ because only the unit-type formation clause is
irreducibly new at this step; see \cref{tab:mbtt-audit}.
No extra eliminator clause is charged yet, since the selected package here is
the type former itself, not its witness.

\paragraph{How it builds on the established library.}
It leverages Step~1 exactly once: $\mathbf{1}$ is introduced as a small type in
$\U_0$.
This makes Step~2 the first genuinely inhabited target available for subsequent
introduction rules.

\paragraph{Why the novelty count has this value.}
The strict total is again exact and equals $\nu=1$; see
\cref{rem:automation-scope,tab:genesis-sequence-trace}.
At this stage the unit type contributes one new foundational rule, but still no
non-trivial interface interaction.

\paragraph{Why it is the unique selected winner at this step.}
The strict trace opens cap/band $1/1$ and selects Unit at
$\rho=1.00 > 0.50$, with no displayed runner-up or near miss; see
\cref{tab:genesis-sequence-trace}.
Since the winner clears the bar with the smallest available cost and bootstrap
admissibility still restricts the search to one-clause packages over the new
universe, \cref{ax:selection} forces this choice.
\end{proof}

\begin{proposition}[Step 3: Witness]
\label{prop:genesis-step3}
At Step~3 the strict genesis trace selects the unit witness
$\star:\mathbf{1}$.
\end{proposition}

\begin{proof}
\paragraph{What it is.}
The selected structure is the canonical inhabitant of the unit type.
Semantically it is the first realized term, not merely the first realized type.

\paragraph{Why $\kappa$ has this value.}
The charged cost is $\kappa=1$, because the selected specification adds exactly
one irreducible introduction clause, namely the witness itself; see
\cref{tab:mbtt-audit}.
The associated elimination behavior is not charged separately at this step
because it is accounted for via adjoint completion rather than as an additional
primitive clause.

\paragraph{How it builds on the established library.}
It uses the already selected unit type as its codomain and thus turns the
bootstrap library from a pure type-formation fragment into an inhabited theory.
This is exactly the bridge from Step~2's abstract type to a concrete term.

\paragraph{Why the novelty count has this value.}
By \cref{lem:adjoint-completion,rem:witness-pisigma-attribution}, the witness
contributes one introduction rule and one elimination action, giving the exact
strict total $\nu=2$ recorded in \cref{tab:genesis-sequence-trace}.
This is the canonical early example where eliminative payload exceeds the raw
charged clause count.

\paragraph{Why it is the unique selected winner at this step.}
The strict trace opens cap/band $1/1$ and selects Witness with
$\rho=2.00$ against bar $1.33$; the displayed top near miss is only a
foundational candidate at $\rho=1.00$; see
\cref{tab:genesis-sequence-trace}.
Hence \cref{ax:selection} picks the witness package as the unique bar-clearer
with minimal overshoot in the foundational one-clause regime.
\end{proof}

\begin{proposition}[Step 4: Dependent function and sum types]
\label{prop:genesis-step4}
At Step~4 the strict genesis trace selects the dependent function/sum package
($\Pi/\Sigma$ types).
\end{proposition}

\begin{proof}
\paragraph{What it is.}
The selected structure is the minimal former package supplying dependent
functions, dependent sums, and the core application/pairing interface.
Although the engine names the winner ``Pi'' in the trace, the mathematical
structure is the joint $\Pi/\Sigma$ former package described in
\cref{tab:mbtt-audit}.

\paragraph{Why $\kappa$ has this value.}
The charged cost is $\kappa=3$, corresponding to the selected specification
listed in \cref{tab:mbtt-audit}: $\lambda$-introduction, pair formation, and
application.
This is the minimal cost at which the library acquires bona fide dependent
former structure.

\paragraph{How it builds on the established library.}
It builds on the previous three bootstrap entries by turning the bare inhabited
universe into a theory with dependent abstraction and pairing.
That infrastructure is the prerequisite for all later HIT, map, and API-shell
steps.

\paragraph{Why the novelty count has this value.}
By \cref{rem:witness-pisigma-attribution}, the exact foundational total is
$\nu=5$: two introduction rules (abstraction and pairing) together with three
elimination rules (application, first projection, second projection).
This is the strict total reported in \cref{tab:genesis-sequence-trace}.

\paragraph{Why it is the unique selected winner at this step.}
The strict trace opens cap/band $3/3$ and finds a singleton viable set, with
$\rho=1.67$ against bar $1.50$; see \cref{tab:genesis-sequence-trace}.
Since Step~4 is the first stage where the library needs a genuine former layer,
and since the selected package is already constructively prime in the sense of
\cref{ax:selection,def:minimal-api}, it is uniquely selected.
\end{proof}

\subsection{Geometric Ascent: Steps 5--9}
\label{subsec:genesis-geometric-proof}

The next five steps move from bootstrap syntax to genuine higher structure.
Here the key theorem-level inputs are the two-layer coherence results
\cref{lem:trace,thm:scaling}, the HIT/map exactness status from
\cref{rem:automation-scope}, and the Agda abstraction-barrier evidence reported
in \cref{sec:verification}.

\begin{proposition}[Step 5: The circle]
\label{prop:genesis-step5}
At Step~5 the strict genesis trace selects the circle $S^1$.
\end{proposition}

\begin{proof}
\paragraph{What it is.}
The selected structure is the first higher inductive type, the circle $S^1$,
with a point and a non-trivial loop.
This is the library's first genuinely homotopical object.

\paragraph{Why $\kappa$ has this value.}
The charged cost is $\kappa=3$, because the selected conceptual shell is the
minimal three-clause HIT signature from \cref{tab:mbtt-audit}: formation,
base point, and loop.
The compiled telescope is larger in the strict trace, but by the proof contract
that extra compiled payload is not extra charged cost.

\paragraph{How it builds on the established library.}
It is the first higher object that can be expressed only after the
$\Pi/\Sigma$ former layer has been installed.
Thus it turns the bootstrap library into a setting where non-trivial path data
can be generated and reused.

\paragraph{Why the novelty count has this value.}
By \cref{ex:foundation-steps}, the exact total is $\nu=7$:
five introduction-style schemas together with two computation/path rules coming
from the loop constructor.
This is the strict HIT reconstruction recorded in
\cref{rem:automation-scope,tab:genesis-sequence-trace}.

\paragraph{Why it is the unique selected winner at this step.}
The strict trace opens cap/band $3/3$ and the viable set is singleton, with
$\rho=2.33$ against bar $2.14$; see \cref{tab:genesis-sequence-trace}.
By \cref{lem:genesis-bar-lemma,ax:selection}, the circle is therefore the
unique least-supercritical entry into the geometric regime.
\end{proof}

\begin{proposition}[Step 6: Propositional truncation]
\label{prop:genesis-step6}
At Step~6 the strict genesis trace selects propositional truncation.
\end{proposition}

\begin{proof}
\paragraph{What it is.}
The selected structure is the propositional truncation operator
$\|{-}\|_0$, which collapses higher structure to a mere proposition while still
remembering existence.
This is the first operation that systematically controls homotopical complexity
rather than only creating it.

\paragraph{Why $\kappa$ has this value.}
The charged cost is again $\kappa=3$, with the conceptual shell audited in
\cref{tab:mbtt-audit} as truncation formation, truncation introduction, and the
squash witness.
As with $S^1$, the strict compiled telescope is larger than the charged
conceptual shell, but the selector charges only the conceptual package.

\paragraph{How it builds on the established library.}
It builds directly on the new geometric branch opened by $S^1$.
In the library-growth dynamics this is the first point where PEN can both
generate higher data and quotient it to mere existence, which is why the
subsequent lift to $S^2$ becomes available.

\paragraph{Why the novelty count has this value.}
The strict total is the exact HIT reconstruction $\nu=8$ recorded in
\cref{rem:automation-scope}.
Operationally this counts the new truncation former, its introduction and
elimination payload, and the squash-induced computation structure physically
present in the strict HIT compilation.

\paragraph{Why it is the unique selected winner at this step.}
This is the tightest selection point in the whole trajectory.
The strict trace opens cap/band $4/3$ and finds exactly two viable candidates:
Trunc at $\rho=8/3$ and $S^2$ at $\rho=10/3$; see
\cref{tab:genesis-sequence-trace}.
Since both have the same charged $\kappa=3$, \cref{ax:selection} reduces to
minimal overshoot, and Trunc wins because its margin is only $0.11$ whereas
$S^2$ overshoots by about $0.77$.
The worked arithmetic in \cref{app:worked-arithmetic} shows that even one extra
charged clause would kill this step, so the winner is both unique and
structurally critical.
\end{proof}

\begin{proposition}[Step 7: The 2-sphere]
\label{prop:genesis-step7}
At Step~7 the strict genesis trace selects the sphere $S^2$.
\end{proposition}

\begin{proof}
\paragraph{What it is.}
The selected structure is the first genuinely two-dimensional sphere, realized
as a minimal higher inductive package.
It is the next geometric lift after the circle/truncation bridge has been
installed.

\paragraph{Why $\kappa$ has this value.}
The charged cost remains $\kappa=3$, because the selected conceptual shell is
the three-clause $S^2$ signature from \cref{tab:mbtt-audit}: formation, base
point, and 2-path surface witness.
Again the strict selector charges the conceptual shell, not every derived
compiled clause.

\paragraph{How it builds on the established library.}
By \cref{lem:trace,thm:scaling}, Step~7 inherits the two-layer interface debt
created by Steps~5 and~6.
The geometric meaning is exactly that $S^2$ is the first post-trunc higher lift:
the library now has both non-trivial loop data and a way to compress it.

\paragraph{Why the novelty count has this value.}
The strict total is the exact HIT reconstruction $\nu=10$ listed in
\cref{rem:automation-scope,tab:genesis-sequence-trace}.
It is larger than the circle's total because a 2-dimensional HIT carries a
richer path/computation payload in the strict evaluator.

\paragraph{Why it is the unique selected winner at this step.}
The strict trace opens cap/band $4/3$ and the viable set is singleton:
$S^2$ clears the bar at $\rho=3.33$, while the displayed near miss is only the
replayed $S^1$ package at $\rho=2.33$; see
\cref{tab:genesis-sequence-trace}.
Thus \cref{ax:selection} again yields a unique winner.
\end{proof}

\begin{proposition}[Step 8: The H-space on $S^3$]
\label{prop:genesis-step8}
At Step~8 the strict genesis trace selects the H-space package on $S^3$.
\end{proposition}

\begin{proof}
\paragraph{What it is.}
The selected structure is not merely a bare 3-sphere, but the minimal
multiplicative H-space package on $S^3$.
This is the first point where PEN prefers a slightly richer geometric shell to a
bare sphere because the additional interaction is worth the extra cost.

\paragraph{Why $\kappa$ has this value.}
The charged cost is $\kappa=5$, exactly the H-space specification selected in
\cref{tab:mbtt-audit}.
The paper's comparison in \cref{rem:kolmogorov-ambiguity} explains why this is
the relevant charged shell: bare 3-sphere encodings exist at lower $\kappa$, but
they overshoot the bar too much and are therefore not selected.

\paragraph{How it builds on the established library.}
By \cref{lem:trace}, Step~8 resolves the Fibonacci debt
$21=13+8$ inherited from Steps~7 and~6.
Geometrically, it uses the just-installed $S^2$ layer together with the
truncation bridge to support the first multiplicative higher sphere.

\paragraph{Why the novelty count has this value.}
The strict total is the exact HIT/H-space reconstruction $\nu=18$ from
\cref{rem:automation-scope}.
More specifically, \cref{rem:kolmogorov-ambiguity} shows that the bare 3-sphere
already contributes $15$ rules, and the H-space laws add $3$ further
elimination-style rules, yielding $18$ in total.

\paragraph{Why it is the unique selected winner at this step.}
The verified strict trace already has a singleton viable set at cap/band $5/5$;
see \cref{tab:genesis-sequence-trace}.
The paper's stronger specification comparison in
\cref{rem:kolmogorov-ambiguity} then explains the uniqueness semantically:
among the natural $S^3$ presentations that clear the bar, minimal overshoot
selects the H-space package rather than the bare suspension or the bare native
HIT.
\end{proof}

\begin{proposition}[Step 9: The Hopf fibration]
\label{prop:genesis-step9}
At Step~9 the strict genesis trace selects the Hopf fibration.
\end{proposition}

\begin{proof}
\paragraph{What it is.}
The selected structure is the map package
$h:S^3\to S^2$ classifying a principal $S^1$-bundle over $S^2$.
This is the first step where the genesis sequence crosses from higher types to a
non-trivial higher map.

\paragraph{Why $\kappa$ has this value.}
The charged cost is $\kappa=4$, exactly the selected Hopf specification audited
in \cref{tab:mbtt-audit}: map, fiber, total-space witness, and classifying map.
Unlike the earlier HIT steps, the conceptual shell and the charged shell agree
here.

\paragraph{How it builds on the established library.}
It leverages the full preceding homotopy package:
the domain $S^3$, the codomain $S^2$, and the fiber $S^1$ are all earlier
library entries.
By \cref{lem:trace} and the Agda evidence summarized in
\cref{sec:verification}, the inherited obligation count is exactly
$34=21+13$, with a domain/codomain split reflecting the $S^3/S^2$ interface.

\paragraph{Why the novelty count has this value.}
For Step~9 the paper gives a sharp qualitative description:
\cref{sec:discussion} records that $\nu_G=\nu_H=0$ and $\nu=\nu_C$.
Thus the exact strict total $\nu=17$ from
\cref{rem:automation-scope,tab:genesis-sequence-trace} is entirely interactional
map novelty rather than new type-former novelty.

\paragraph{Why it is the unique selected winner at this step.}
The strict trace opens cap/band $5/4$ and the viable set is singleton, so Hopf
is already the unique strict bar-clearer at $\rho=4.25$ against bar $4.01$; see
\cref{tab:genesis-sequence-trace}.
This is the point where the codebase's structural regime shifts from higher
types to higher maps, but in the canonical bounded live atomic lane no additional rescue
heuristic is needed: the viable set has already collapsed to a single winner.
\end{proof}

\subsection{Framework Abstraction: Steps 10--14}
\label{subsec:genesis-framework-proof}

From Step~10 onward the sequence grows by library specifications rather than by
new foundational judgment forms; see \cref{rem:library-api}.
The semantic force of each step comes from constructive irreducibility and
structural unity, not from opaque axiom packing.
For Step~10 the rule audit is exact; for Steps~11--14 it is expository but
constrained to match the strict totals by \cref{rem:automation-scope}.

\begin{proposition}[Step 10: Cohesion]
\label{prop:genesis-step10}
At Step~10 the strict genesis trace selects the cohesive modality package.
\end{proposition}

\begin{proof}
\paragraph{What it is.}
The selected structure is the adjoint quadruple of modalities
$(\flat,\sharp,\Disc,\Pi_{\infty})$ described in
\cref{app:cohesion-spec}.
This is the first abstraction step: PEN stops adding one more homotopy object
and instead adds a modal interface for the whole geometric library.

\paragraph{Why $\kappa$ has this value.}
The charged cost is $\kappa=4$, exactly the four modality-formation clauses
listed in \cref{app:cohesion-spec,tab:mbtt-audit}.
By \cref{def:derivability,def:minimal-api}, these are irreducible API clauses:
the adjoint string cannot be obtained definitionally from the prior HoTT
library.

\paragraph{How it builds on the established library.}
By \cref{rem:library-api}, a late API shell is valuable only insofar as it
interacts with earlier structure.
That interaction is explicit here:
\cref{app:cohesion-spec} records cross-interactions with $S^1$, $S^2$, $S^3$,
and Hopf, so Cohesion is a global abstraction of the entire geometric ascent.

\paragraph{Why the novelty count has this value.}
For Step~10 the rule audit is exact:
\cref{app:cohesion-spec} gives $\nu=19=2+17$, namely two introduction-style unit
maps together with seventeen elimination/interaction rules.
This exact total is the one selected in the strict trace of
\cref{tab:genesis-sequence-trace}.

\paragraph{Why it is the unique selected winner at this step.}
The strict trace opens cap/band $5/4$ and again finds a singleton viable set:
Cohesion clears the bar at $\rho=4.75$, while the displayed near miss is the old
Hopf package at $\rho=4.25$; see \cref{tab:genesis-sequence-trace}.
Hence \cref{ax:selection} selects the cohesive shell uniquely, and the
non-survival of Hopf shows that PEN is now in a genuinely new abstraction
regime.
\end{proof}

\begin{proposition}[Step 11: Connections]
\label{prop:genesis-step11}
At Step~11 the strict genesis trace selects the connection package.
\end{proposition}

\begin{proof}
\paragraph{What it is.}
The selected structure is the differential-geometric package centered on a
connection $\nabla$ together with transport, covariant lift, horizontal lift,
and Leibniz interaction; see \cref{app:connections-spec}.
This is the first differential shell built over the cohesive interface.

\paragraph{Why $\kappa$ has this value.}
The charged cost is $\kappa=5$, exactly the five API clauses listed in
\cref{app:connections-spec,tab:mbtt-audit}.
By \cref{def:derivability}, a connection is not derivable from cohesion alone,
so the package must be charged explicitly.

\paragraph{How it builds on the established library.}
It leverages Cohesion as its immediate carrier and thereby translates the modal
library into differential transport data.
This is the first late step whose entire meaning is cross-interface leverage:
without the cohesive layer from Step~10 there is no differential package to
select.

\paragraph{Why the novelty count has this value.}
The expository rule audit in \cref{app:connections-spec} reconstructs the strict
total as $\nu=26=2+24$, where the new interaction comes from cohesive
cross-rules, function-space rules, and library interactions with the previously
selected HITs and Hopf.
By \cref{rem:automation-scope}, this reconstruction is subordinate to, but
consistent with, the strict total in \cref{tab:genesis-sequence-trace}.

\paragraph{Why it is the unique selected winner at this step.}
The strict trace opens cap/band $6/5$ and the viable set is singleton:
Connections clears the bar at $\rho=5.20$, while the displayed near miss is only
the old $S^3$ package at $\rho=3.60$; see
\cref{tab:genesis-sequence-trace}.
Thus the bounded live atomic lane does not need any external steering here: the first
differential shell is already the unique bar-clearing continuation.
\end{proof}

\begin{proposition}[Step 12: Curvature]
\label{prop:genesis-step12}
At Step~12 the strict genesis trace selects the curvature package.
\end{proposition}

\begin{proof}
\paragraph{What it is.}
The selected structure is the curvature layer
$R=d\nabla+\nabla\wedge\nabla$ together with Bianchi, holonomy, Chern--Weil, and
characteristic-class structure; see \cref{app:curvature-spec}.
It is the first shell whose meaning is ``second-order differential structure''
rather than just differential transport.

\paragraph{Why $\kappa$ has this value.}
The charged cost is $\kappa=6$, exactly the six irreducible clauses audited in
\cref{app:curvature-spec,tab:mbtt-audit}.
These clauses are constructively irreducible because the elimination behavior
they expose is not definitionally recoverable from a bare connection package;
see \cref{def:derivability,app:curvature-spec}.

\paragraph{How it builds on the established library.}
It leverages Connections as its immediate carrier and Cohesion as its ambient
modal base.
This is the canonical next abstraction after Step~11: once the library can
differentiate, the next non-derivable interaction is precisely the failure of
flatness.

\paragraph{Why the novelty count has this value.}
The expository audit in \cref{app:curvature-spec} reconstructs the strict total
as $\nu=34=1+33$, with the extra weight coming from interactions among
curvature, connection, and cohesive structure.
By \cref{rem:automation-scope}, that reconstruction is explanatory rather than a
second oracle, but it agrees with the canonical strict total in
\cref{tab:genesis-sequence-trace}.

\paragraph{Why it is the unique selected winner at this step.}
The strict trace opens cap/bands $6/(5|6)$ and nevertheless yields a singleton
viable set: Curvature clears the bar at $\rho=5.67$, while the displayed near
misses are far below the bar; see \cref{tab:genesis-sequence-trace}.
Thus the switch from differential to curvature structure is not a naming
convention but a mechanically unique strict continuation.
\end{proof}

\begin{proposition}[Step 13: Endomorphic operator bundle]
\label{prop:genesis-step13}
At Step~13 the strict genesis trace selects the endomorphic operator bundle rather than a
bare tensor or scalar-heavy detour.
\end{proposition}

\begin{proof}
\paragraph{What it is.}
The selected structure is the endomorphic operator bundle described in
\cref{app:metric-spec,tab:step13}: an endomorphic pair shell together with
Levi-Civita, geodesic, volume, Hodge, Laplacian, and Ricci operators.
Its classical metric reading is semantic and post hoc, not a literal claim
that the raw AST already contains a scalar-valued tensor field.
This is the first step where PEN selects an explicitly geometric API bundle
rather than a single operator.

\paragraph{Why $\kappa$ has this value.}
The charged cost is $\kappa=7$, exactly the seven clauses audited in
\cref{app:metric-spec,tab:step13}.
By \cref{def:minimal-api}, this package is charged as a single structural bundle:
the bare metric tensor is too weak, whereas the coupled bundle is the smallest
constructively prime package that can survive the Step~13 bar.

\paragraph{How it builds on the established library.}
It leverages Curvature and Connections directly and uses the cohesive/differential
layers indirectly through those references.
The key structural point, emphasized in \cref{app:metric-spec}, is that the
metric package is selected because it unlocks the whole chain
$g\to\nabla_g\to R\to\mathrm{Ric}$ inside the already accumulated differential
library.

\paragraph{Why the novelty count has this value.}
The expository reconstruction in \cref{app:metric-spec} gives
$\nu=46=4+42$, where the dominant term is the bounded coupling with the
curvature and connection layers.
That elimination term is not a local heuristic: \cref{thm:telescopic-elimination}
identifies it exactly from the import DAG, the seven internal clauses, and the
single bridge needed to fuse Curvature and Connections into one telescope.
The same section explains why a bare metric tensor would have
$\rho\le 3.0$, far below the Step~13 bar, so the strict total in
\cref{tab:genesis-sequence-trace} is genuinely bundle-driven rather than
tensor-driven.

\paragraph{Why it is the unique selected winner at this step.}
The strict trace opens cap/band $8/7$ and the viable set is singleton:
Metric clears the bar at $\rho=6.57$, whereas the displayed near miss is only an
unnamed HIT candidate at $\rho=2.29$; see
\cref{tab:genesis-sequence-trace}.
Hence the selected operator bundle is uniquely forced by the strict objective, and
\cref{def:minimal-api} explains why smaller fragments are excluded.
\end{proof}

\begin{proposition}[Step 14: Hilbert functional]
\label{prop:genesis-step14}
At Step~14 the strict genesis trace selects the Hilbert-functional package.
\end{proposition}

\begin{proof}
\paragraph{What it is.}
The selected structure is the operator-analytic package described in
\cref{app:hilbert-spec,tab:step14}: inner product, completeness, orthogonal and
spectral decomposition, $C^*$-algebra structure, cross-links to Metric,
Curvature, and Connections, and a functional derivative.
This is the analytic culmination of the framework-abstraction phase.

\paragraph{Why $\kappa$ has this value.}
The charged cost is $\kappa=9$, exactly the nine clauses audited in
\cref{app:hilbert-spec,tab:step14}.
By \cref{def:derivability}, none of this operator-analytic structure is
definitionally free from the preceding geometric library.

\paragraph{How it builds on the established library.}
It leverages the three immediately relevant layers: Metric, Curvature, and
Connections.
This is a textbook PEN abstraction step: the new package is not selected for
opaque existence but because it compresses and re-exposes the interaction of the
whole geometric stack in analytic form.

\paragraph{Why the novelty count has this value.}
The expository reconstruction in \cref{app:hilbert-spec} gives
$\nu=62=5+57$, dominated by the coupling to Metric, Curvature, and Connections.
Here again \cref{thm:telescopic-elimination} supplies the exact $\nu_C$ law:
the Hilbert shell inherits the Metric derivation base, contributes nine internal
eliminators from its own clauses, and pays two bridge schemas to amalgamate the
three imported interfaces into one analytic telescope.
The same section also shows why the margin is small: even two extra charged
clauses would drop the package below the Step~14 bar.

\paragraph{Why it is the unique selected winner at this step.}
The strict trace opens cap/bands $9/(8|9)$ and still yields a singleton viable
set: Hilbert clears the bar at $\rho=6.89$, while the displayed near miss is
only an unnamed HIT candidate at $\rho=3.00$; see
\cref{tab:genesis-sequence-trace}.
Thus the operator-analytic layer is the unique strict continuation of the metric
bundle, with no competing bar-clearer in the verified run.
\end{proof}

\subsection{Terminal Synthesis: Step 15}
\label{subsec:genesis-terminal-proof}

Step~15 is special.
The proof therefore proceeds shell-first: the strict trace selects a
temporal-cohesive AST, the executable artifact fixes the canonical total
$\nu=103$ attached to that shell, and ``DCT'' names the semantic completion
used to interpret it; see \cref{sec:decomposition,rem:automation-scope}.

\begin{proposition}[Step 15: Temporal-cohesive shell]
\label{prop:genesis-step15}
At Step~15 the strict genesis trace selects the temporal-cohesive shell whose
semantic completion is read as DCT.
\end{proposition}

\begin{proof}
\paragraph{What it is.}
The selected AST skeleton is the temporal-cohesive shell audited in
\cref{tab:step15}: next modality, eventually modality, the map
$\bigcirc\to\Diamond$, and the explicit exchange laws between temporal and
cohesive structure.
Its semantic completion is the DCT reading established in
\cref{sec:decomposition}.

\paragraph{Why $\kappa$ has this value.}
The charged cost is $\kappa=8$, exactly the eight temporal-cohesive clauses
audited in \cref{tab:step15}.
As stressed in \cref{tab:step15,sec:decomposition}, the infinitesimal reading is
\emph{not} charged as an extra primitive clause; it is derived semantically from
the temporal-cohesive shell already selected.

\paragraph{How it builds on the established library.}
It builds directly on Cohesion and sits atop the already selected Hilbert layer.
This is the terminal synthesis step because it fuses the earlier spatial logic
of Step~10 with temporal operators in a way that internalizes the library's own
evolution; see \cref{thm:end-of-history,sec:decomposition}.

\paragraph{Why the novelty count has this value.}
The executable structural total $\nu=103$ is canonical here.
By \cref{sec:verification,sec:adjoint-completion}, the strict engine computes
that total from the conceptual temporal shell together with the explicit
path-style witness payload preserved by compilation.
This is exactly the number reported in \cref{tab:genesis-sequence-trace}.

\paragraph{Why it is the unique selected winner at this step.}
The strict trace opens cap/band $9/8$ and the viable set is singleton:
DCT clears the bar at $\rho=12.88$ against bar $7.40$, while the displayed near
miss is only an unnamed HIT candidate at $\rho=3.00$; see
\cref{tab:genesis-sequence-trace}.
This is also the largest absolute overshoot in the entire sequence.
Thus the temporal-cohesive shell is not merely admissible but uniquely dominant
within the current bounded live atomic lane at the terminal step.
\end{proof}

\subsection{Conclusion of the Step-by-Step Proof}
\label{subsec:genesis-proof-conclusion}

\begin{theorem}[The strict genesis sequence]
\label{thm:strict-genesis-sequence}
In the current bounded live atomic PEN lane, the unique step-by-step sequence of
selected structures through Step~15 is
\[
  \begin{aligned}
    &\U_0,\ \mathbf{1},\ \star,\ \Pi/\Sigma,\ S^1,\ \|{-}\|_0,\ S^2,\ S^3,\\
    &\mathrm{Hopf},\ \mathrm{Modal},\ \mathrm{Connections},\\
    &\mathrm{Curvature},\ \mathrm{EndoOp},\ \mathrm{HilbertShell},\\
    &\mathrm{TemporalShell}.
  \end{aligned}
\]
For each $1\le n\le 15$, the selected structure has the corresponding values of
$\Delta_n$, $\nu_n$, $\kappa_n$, $\rho_n$, and $\mathrm{Bar}_n$ listed in
\cref{tab:genesis-sequence-trace}.
\end{theorem}

\begin{proof}
Combine the global bar formula of \cref{lem:genesis-bar-lemma} with the fifteen
step propositions
\cref{prop:genesis-step1,prop:genesis-step2,prop:genesis-step3,prop:genesis-step4,prop:genesis-step5,prop:genesis-step6,prop:genesis-step7,prop:genesis-step8,prop:genesis-step9,prop:genesis-step10,prop:genesis-step11,prop:genesis-step12,prop:genesis-step13,prop:genesis-step14,prop:genesis-step15}.
Each step proposition identifies the selected AST skeleton, justifies its charged
cost and novelty total at the appropriate proof-strength level, and verifies
that it is the unique selected winner in the current bounded live atomic lane at its step.
\end{proof}

\begin{corollary}[No Step~16 continuation in the current bounded live atomic lane]
\label{cor:genesis-no-step16}
No Step~16 candidate clears the selection bar in the current bounded live atomic
lane.
\end{corollary}

\begin{proof}
By \cref{thm:strict-genesis-sequence}, Step~15 is the selected temporal-cohesive
shell with the canonical executable total $\nu=103$.
By the halting theorem \cref{thm:end-of-history}, every internal Step~16
continuation has zero marginal external novelty and every disconnected
continuation fails the positive Step~16 bar.
Hence no Step~16 candidate clears the bar.
\end{proof}

\begin{thebibliography}{99}

\bibitem{hott}
Univalent Foundations Program.
\textit{Homotopy Type Theory: Univalent Foundations of Mathematics}.
Institute for Advanced Study, 2013.

\bibitem{cubical}
C.~Cohen, T.~Coquand, S.~Huber, A.~M\"ortberg.
Cubical Type Theory: a constructive interpretation of the univalence axiom.
\textit{TYPES 2015}, 2015.

\bibitem{schreiber}
U.~Schreiber.
Differential Cohomology in a Cohesive Infinity-Topos.
arXiv:1310.7930, 2013.

\bibitem{lawvere}
F.~W.~Lawvere.
Axiomatic Cohesion.
\textit{Theory and Applications of Categories}, 19(3), 2007.

\bibitem{nakano}
H.~Nakano.
A Modality for Recursion.
\textit{Proceedings of LICS}, 2000.

\bibitem{cchm-hits}
T.~Coquand, S.~Huber, A.~M\"ortberg.
On Higher Inductive Types in Cubical Type Theory.
\textit{LICS 2018}, 2018.

\bibitem{cubical-agda}
A.~Vezzosi, A.~M\"ortberg, A.~Abel.
Cubical Agda: A Dependently Typed Programming Language with Univalence and Higher Inductive Types.
\textit{ICFP}, 2019.

\bibitem{licata-shulman}
D.~R.~Licata, M.~Shulman.
Calculating the Fundamental Group of the Circle in Homotopy Type Theory.
\textit{LICS}, 2013.

\bibitem{shulman-real-cohesive}
M.~Shulman.
Brouwer's Fixed-Point Theorem in Real-Cohesive Homotopy Type Theory.
\textit{Mathematical Structures in Computer Science}, 28(6), 2018.

\bibitem{alphageometry}
T.~Trinh, Y.~Wu, Q.~Le, H.~He, and collaborators.
Solving Olympiad Geometry without Human Demonstrations.
\textit{Nature}, 2024.

\bibitem{leandojo}
K.~Yang, A.~Selsam, K.~Pang, et al.
LeanDojo: Theorem Proving with Retrieval-Augmented Language Models.
\textit{NeurIPS}, 2023.

\bibitem{openai-lean}
OpenAI.
Advancing AI-Assisted Formal Mathematics.
Research update, 2025.
\url{https://openai.com/index/advancing-ai-assisted-formal-mathematics/}

\bibitem{quickspec}
N.~Smallbone, M.~Johansson, K.~Claessen, M.~Algehed.
Quick Specifications for the Busy Programmer.
\textit{Journal of Functional Programming}, 27:e18, 2017.

\bibitem{isacosy}
M.~Johansson, L.~Dixon, A.~Bundy.
Conjecture Synthesis for Inductive Theories.
\textit{Journal of Automated Reasoning}, 47:251--289, 2011.

\bibitem{hipster}
M.~Johansson, D.~Ros\'en, N.~Smallbone, K.~Claessen.
Hipster: Integrating Theory Exploration in a Proof Assistant.
\textit{CICM}, LNCS 8543, pp.~108--122, 2014.

\bibitem{synquid}
N.~Polikarpova, I.~Kuraj, A.~Solar-Lezama.
Program Synthesis from Polymorphic Refinement Types.
\textit{PLDI}, 2016.

\bibitem{lenat-am}
D.~B.~Lenat.
\textit{AM: Discovery in Mathematics as Heuristic Search}.
McGraw-Hill, 1983.

\bibitem{colton-atf}
S.~Colton.
\textit{Automated Theory Formation in Pure Mathematics}.
Springer, 2002.

\bibitem{lurie}
J.~Lurie.
\textit{Higher Topos Theory}.
Annals of Mathematics Studies, vol.~170, Princeton University Press, 2009.

\bibitem{lumsdaine}
P.~L.~Lumsdaine.
Weak $\omega$-Categories from Intensional Type Theory.
\textit{TLCA}, LNCS 6690, pp.~172--187, 2010.

\bibitem{vdberg-garner}
B.~van den Berg, R.~Garner.
Types are Weak $\omega$-Groupoids.
\textit{Proc.\ London Math.\ Soc.}, 102(2):370--394, 2011.

\bibitem{kuratowski}
K.~Kuratowski.
Sur l'op\'eration de l'$\bar{A}$.
\textit{Fundamenta Mathematicae}, 3, 1922.

\bibitem{pnueli}
A.~Pnueli.
The Temporal Logic of Programs.
\textit{Proceedings of FOCS}, 1977.

\bibitem{kock}
A.~Kock.
\textit{Synthetic Differential Geometry}.
Cambridge University Press, 2nd edition, 2006.

\bibitem{maclane-coherence}
S.~Mac Lane.
Natural Associativity and Commutativity.
\textit{Rice University Studies}, 49(4):28--46, 1963.

\bibitem{stasheff}
J.~D.~Stasheff.
Homotopy Associativity of H-Spaces~I, II.
\textit{Trans.\ Amer.\ Math.\ Soc.}, 108(2):275--312, 1963.

\bibitem{kraus-vonraumer}
N.~Kraus, J.~von Raumer.
Coherence via Well-Foundedness.
\textit{Proceedings of LICS}, 2019.

\bibitem{kolmogorov}
M.~Li, P.~Vit\'anyi.
\textit{An Introduction to Kolmogorov Complexity and Its Applications}.
Springer, 4th edition, 2019.
(The Invariance Theorem: Theorem~2.1.1.)

\bibitem{lawvere-adjoint}
F.~W.~Lawvere.
Adjointness in Foundations.
\textit{Dialectica}, 23(3/4):281--296, 1969.
Reprinted in \textit{Repr.\ Theory Appl.\ Categ.}, 16:1--16, 2006.

\bibitem{godel}
K.~G\"odel.
\"Uber formal unentscheidbare S\"atze der Principia Mathematica und verwandter Systeme~I.
\textit{Monatshefte f\"ur Mathematik und Physik}, 38:173--198, 1931.

\end{thebibliography}

\end{document}
